%% Template for the submission to:
%%   The Annals of Statistics [AOS]
%%
%%%%%%%%%%%%%%%%%%%%%%%%%%%%%%%%%%%%%%%%%%%%%%
%% In this template, the places where you   %%
%% need to fill in your information are     %%
%% indicated by '???'.                      %%
%%                                          %%
%% Please do not use \input{...} to include %%
%% other tex files. Submit your LaTeX       %%
%% manuscript as one .tex document.         %%
%%%%%%%%%%%%%%%%%%%%%%%%%%%%%%%%%%%%%%%%%%%%%%

\documentclass[aos]{imsart}

%% Packages
\RequirePackage{amsthm,amsmath,amsfonts,amssymb,booktabs,array}
\RequirePackage[numbers]{natbib}
%\RequirePackage[authoryear]{natbib}%% uncomment this for author-year citations
\RequirePackage[colorlinks,citecolor=blue,urlcolor=blue]{hyperref}%% uncomment this for coloring bibliography citations and linked URLs
\RequirePackage{graphicx}%% uncomment this for including figures

\startlocaldefs
%%%%%%%%%%%%%%%%%%%%%%%%%%%%%%%%%%%%%%%%%%%%%%
%%                                          %%
%% Uncomment next line to change            %%
%% the type of equation numbering           %%
%%                                          %%
%%%%%%%%%%%%%%%%%%%%%%%%%%%%%%%%%%%%%%%%%%%%%%
%\numberwithin{equation}{section}
%%%%%%%%%%%%%%%%%%%%%%%%%%%%%%%%%%%%%%%%%%%%%%
%%                                          %%
%% For Axiom, Claim, Corollary, Hypothesis, %%
%% Lemma, Theorem, Proposition              %%
%% use \theoremstyle{plain}                 %%
%%                                          %%
%%%%%%%%%%%%%%%%%%%%%%%%%%%%%%%%%%%%%%%%%%%%%%
\theoremstyle{plain}
\newtheorem{thm}{\textsc{Theorem}}[section]
\newtheorem{prop}[thm]{\textsc{Proposition}}
\newtheorem{cor}[thm]{\textsc{Corollary}}
\newtheorem{lem}[thm]{\textsc{Lemma}}
%%%%%%%%%%%%%%%%%%%%%%%%%%%%%%%%%%%%%%%%%%%%%%
%%                                          %%
%% For Assumption, Definition, Example,     %%
%% Notation, Property, Remark, Fact         %%
%% use \theoremstyle{remark}                %%
%%                                          %%
%%%%%%%%%%%%%%%%%%%%%%%%%%%%%%%%%%%%%%%%%%%%%%
%\theoremstyle{remark}
%\newtheorem{???}{???}
%\newtheorem*{???}{???}
%\newtheorem{???}{???}[???]
%\newtheorem{???}[???]{???}
%%%%%%%%%%%%%%%%%%%%%%%%%%%%%%%%%%%%%%%%%%%%%%
%% Please put your definitions here:        %%
%%%%%%%%%%%%%%%%%%%%%%%%%%%%%%%%%%%%%%%%%%%%%%
\newcommand{\david}[1]{{\small{\color{blue} \sf $\heartsuit$ David: #1}}}
\newcommand{\piotr}[1]{{\small{\color{red} \sf $\clubsuit$ Piotr: #1}}}
\newcommand{\paul}[1]{{\small{\color{green} \sf $\spadesuit$ Paul: #1}}}
\newcommand{\E}{\mathbb{E}}
\newcommand{\N}{\mathbb{N}}
\newcommand{\R}{\mathbb{R}}
\newcommand{\I}{\mathbb{I}}
\newcommand{\M}{\mathcal{M}}
\newcommand{\T}{\mathcal{T}}
\newcommand{\bs}{\boldsymbol}
\newcolumntype{M}[1]{>{\centering\arraybackslash}m{#1}}
\endlocaldefs

\begin{document}

\begin{frontmatter}
%%%%%%%%%%%%%%%%%%%%%%%%%%%%%%%%%%%%%%%%%%%%%%
%%                                          %%
%% Enter the title of your article here     %%
%%                                          %%
%%%%%%%%%%%%%%%%%%%%%%%%%%%%%%%%%%%%%%%%%%%%%%
\title{Improving variable selection properties by integrating external data}
%\title{A sample article title with some additional note\thanksref{T1}}
\runtitle{Improving variable selection properties by integrating external data}
\thankstext{T1}{Working paper, to be submitted this January.}

\begin{aug}
%%%%%%%%%%%%%%%%%%%%%%%%%%%%%%%%%%%%%%%%%%%%%%%
%% Only one address is permitted per author. %%
%% Only division, organization and e-mail is %%
%% included in the address.                  %%
%% Additional information can be included in %%
%% the Acknowledgments section if necessary. %%
%% ORCID can be inserted by command:         %%
%% \orcid{0000-0000-0000-0000}               %%
%%%%%%%%%%%%%%%%%%%%%%%%%%%%%%%%%%%%%%%%%%%%%%%
\author[A]{\fnms{Paul}~\snm{Rognon-Vael}\ead[label=e1]{paul.rognon@gmail.com}},
\author[A]{\fnms{David}~\snm{Rossell}\ead[label=e2]{rosselldavid@gmail.com}}
\and
\author[B]{\fnms{Piotr}~\snm{Zwiernik}\ead[label=e3]{piotr.zwiernik@utoronto.ca}}
%%%%%%%%%%%%%%%%%%%%%%%%%%%%%%%%%%%%%%%%%%%%%%
%% Addresses                                %%
%%%%%%%%%%%%%%%%%%%%%%%%%%%%%%%%%%%%%%%%%%%%%%
\address[A]{Department of Economics and Business, Universitat Pompeu Fabra \printead[presep={ ,\ }]{e1,e2}}
\address[B]{Department of Statistical Sciences, University or Toronto
\printead[presep={,\ }]{e3}}
\end{aug}

\begin{abstract}
We study the variable selection properties of $\ell_0$ penalties depending on external information in the Gaussian sequence model and high-dimensional linear regression. We consider in particular the case where external information helps partition variables into blocks
with potentially distinct characteristics and study non-exchangeable block-dependent penalties. We show that such block informed penalization achieves variable selection consistency under milder conditions than standard penalization and that oracle block penalties converge at faster rates. We propose a data-based method to incorporate external information in selection that is grounded in Bayesian theory and empirical Bayes techniques. We demonstrate it realizes theoretical benefits of block informed penalization without requiring oracles.  
\end{abstract}

\begin{keyword}[class=MSC]
\kwd[Primary ]{62F07}
\kwd[; secondary ]{62C12}
\kwd{62R07}
\end{keyword}

\begin{keyword}
\kwd{variable selection}
\kwd{high-dimensional statistics}
\kwd{$\ell_0$ penalty}
\kwd{empirical Bayes}
\end{keyword}
\end{frontmatter}
%%%%%%%%%%%%%%%%%%%%%%%%%%%%%%%%%%%%%%%%%%%%%%
%% Please use \tableofcontents for articles %%
%% with 50 pages and more                   %%
%%%%%%%%%%%%%%%%%%%%%%%%%%%%%%%%%%%%%%%%%%%%%%
%\tableofcontents

%%%%%%%%%%%%%%%%%%%%%%%%%%%%%%%%%%%%%%%%%%%%%%
%%%% Main text entry area:

\paul{\Large NOT THE CURRENT VERSION}

\section{Introduction} 

High-dimensional inference theory relies on suitable assumptions regarding sparsity and signal strength.
Our motivation is that in many applications one has external information regarding each parameter, e.g. its magnitude or its likelihood of being non-zero, that can be leveraged to relax said assumptions and enhance inference.
%In many modern applications, the sample size is insufficient to reliably do inference. This limitation has motivated the development of high-dimensional techniques that frequently rely on strict assumptions on sparsity and signal strength. Yet, often, external information is available and can be leveraged to enhance inference.
%Our motivation stems from the observation that, often, 
In particular, external information can guide our decisions regarding which parameters to include in a regression model. This information may originate from each variable's inherent nature (e.g., clinical history versus genomic markers, sociodemographics versus job history),
information from previous datasets (e.g., studying related cancer types), or meta-covariates
%or external data that distinguishes relevant variables from irrelevant ones 
(e.g., functional annotations on genes). 
We investigate this concept in the Gaussian sequence model (GSM) and in linear regression, where external information partitions variables into blocks with potentially distinct characteristics. We show that said information allows pushing the mathematical conditions under which consistent model recovery is possible, and improving the associated rates.

Using external information is often advocated within the frameworks of data integration and transfer learning,  as joint learning can lead to more accurate inference than analyzing datasets separately. Indeed, numerous applied works employed external information to guide inference. For instance, \cite{stingo:2011, quintana, cassese:2014} proposed Bayesian variable selection methods for gene expression, where prior probabilities for non-zero coefficients depend on biological knowledge and meta-covariates. Additionally, \cite{chen_tinghuei:2021} predicted disease outcomes by allowing LASSO penalties to depend on functional annotation categories. Beyond regression, node and edge covariates have been incorporated in \cite{peterson} and \cite{jewson2023graphical} to drive edge inclusion in Gaussian graphical models, while \cite{schiavon2022generalized} used meta-covariates to determine non-zero loadings in factor models. In causal analysis, the inclusion of control covariates is driven by their degree of association with the covariates of interest (referred to as treatments) 
%guide the inclusion of control covariates based on their degree of association with the treatments 
\cite{belloni2014inference,antonelli2021bayesian}. Collectively, empirical evidence consistently demonstrates improved structural learning when suitable external data is integrated. However, despite this empirical success, a theoretical framework explaining precisely why this occurs remains elusive.

%We show that incorporating external information allows pushing the theoretical limits under which consistent variable selection is possible, understood as recovering the true model with probability going to 1 as the sample size $n$ and the number of variables $p$ grow. 
Previous literature has extensively explored how high-dimensional variable selection is constrained by inherent characteristics of the dataset, such as the number of samples $n$, the dimension $p$, the signal strength, the correlation between variables and the number of variables truly associated to the outcome \cite{buhlmannvandegeerbook,overviewvarsel,handbookbayesvarsel,wainwrighthighdimbook}. In this paper we analyse how conditions for consistency can be relaxed in the presence of external information. We focus on variable selection but using external information is also expected to benefit parameter estimation and prediction problems. 

To make our ideas concrete, consider variable selection in the Gaussian linear regression
\begin{equation}\label{eq:linearmodel}
	\bs y \;=\; \bs X \bs \beta^* + \bs\epsilon,\qquad \bs \epsilon \sim N(0,\sigma^2 I_n),
\end{equation}
where $\bs X\in \R^{n\times p}$, $\bs y\in \R^n$, and $\sigma >0$. A large class of methods operate by penalizing the size of an estimated $\hat{\bs  \beta}$. For instance, penalized likelihood methods optimize the log-likelihood plus a penalty term driven by the $\ell_q$ "norm" of $\hat{\bs  \beta}$ for $q\in [0,1]$ or folded concave penalties \citep{lasso,overviewvarsel,bertsimas}. 
In Bayesian settings, variable selection is often based on posterior model probabilities that are directly connected to $\ell_0$ penalties \citep{BIC,EBIC,rossell2022concentration}. 

Our interest is in the setting where the external information partitions the $p$ variables into $b$ blocks denoted $B_j \subset \{1,\ldots,p\}$, $j=1,\ldots,b$. The idea is that key characteristics such as the level of sparsity or signal strength may vary across the $B_j$'s. We introduce block $\ell_0$ penalties that allow modulating the strength of the penalty in each block. Unlike standard $\ell_0$ penalties, like BIC \cite{BIC}, RIC \cite{RIC} or EBIC \cite{EBIC}, block penalties are non-exchangeable in the sense that the penalty for adding a variable may depend on its block $B_j$.

Our main contributions are as follows. We show that block-based non-exchangeable $\ell_0$ penalties soften the theoretical conditions for consistent model recovery in the sequence model and in linear regression. Consistent recovery refers to recovering the true model with probability going to 1 as $n$ and $p$ grow. We first 
show that an oracle may set block-based penalties so that variable selection consistency is either attained where otherwise it would not be possible, or is attained at a faster rate. 
We next introduce data-based variable selection procedures, based on empirical Bayes, that realize the theoretical benefits without requiring an oracle. Other more minor contributions include tight necessary conditions on signal strength in the GSM and linear regression.

Our work has connections with multiple hypothesis testing ideas. In this line of research, \cite{genovese2006} proposed a false discovery rate (FDR) procedure in which p-values are weighted based on prior information. They show that if the weights are positively associated to the null hypotheses being false, their procedure improves power. Subsequent works discussed oracle choices of weights either based on external information or derived from the data \cite{roederwasserman,basu}. 
Adaptive FDR methods \cite{blanchard} estimate the
proportion of truly null hypotheses to improve statistical power.
Recently, \cite{ramdas2019prior} combined external information and adaptivity in a group-based FDR control. 
Relative to this work, we study variable selection consistency in high-dimensional regression as well as the associated conditions on sparsity and signal strength.

Also related to our work, \cite{scarlett2012compressed} studied compressed sensing when prior information allows splitting parameters into blocks, where one knows the true proportion of non-zero parameters in each block. They show that sparse signal recovery requires smaller $n$ than with exchangeable penalties. Key differences with our work are their assuming independence across covariates, which renders the results inapplicable to regression. Also, they do not consider the sequence model, nor that the proportions of non-zero parameters are unknown in practice.

The paper is organized as follows. Section~\ref{sec:intro_blockpen} introduces block $\ell_0$ penalization. Section~\ref{sec:properties_orth} presents its
variable selection properties and its benefits in the Gaussian sequence model, in an oracle setting where the sparsity and betamin conditions are known for all blocks.
Those results apply to essentially any penalized likelihood or Bayesian method, including $\ell_0$ and $\ell_1$ penalties, because in that setting selection operates by thresholding \cite{davidomiros,rigollethighdimensional}. Section~\ref{sec:properties_legression} studies block $\ell_0$ penalties in linear regression, and shows analogous benefits to those in Section~\ref{sec:properties_orth}. These results can be extended to a wide class of Bayesian variable selection methods. Section~\ref{sec:informed_l0} presents data-based 
procedures, motivated by empirical Bayes, that achieve faster convergence without requiring an oracle. Section~\ref{sec:simulation} shows empirical examples and Section~\ref{sec:discussion} concludes. Proofs are gathered at the end of the paper.
\medskip

\noindent \textbf{Notation:}  
We denote by $\bs \beta\in \R^p$ the parameters of interest 
and by $\bs \beta^*$ their true values. Let  $V=\{1,\ldots,p\}$. For any $A\subseteq V$ and any vector $\bs x\in \R^p$, $\bs x_A$ denotes the subvector of $\bs x$ with entries corresponding to indices in $A$. For any matrix $\bs X \in \R^{n\times p}$, $\bs X_A$ denotes the submatrix of $\bs X$ obtained by selecting the \emph{columns} with indices in $A$. We denote by
$S=\left\{i\in V:\,\,  \beta^*_i \neq 0 \right\}$ the true support of $\bs \beta^*$, its size is $s=|S|$, and hence $p-s$ is the number of truly inactive parameters.
Denote $\beta_{\min}^*=\min_{i \in S}|\beta^*_i|$ the smallest true signal.
We assume that $V$ is partitioned into a fixed number $b$ of disjoint blocks $B_j \subseteq V$ for $j=1,\ldots,b$.
We denote by $S_j=\left\{i\in B_j:\,\,  \beta^*_i \neq 0\right\}$ the set of active parameters in block $B_j$,  $s_j=|S_j|$ its size, $p_j-s_j$ the number of inactive parameters in the block,
and $\beta_{\min,j}^*=\min_{i \in S_j}|\beta^*_i|$.
Given sequences $f(n) > 0$ and $g(n) >0$, $f(n)=O(g(n))$ means that there exists a constant $c<\infty$ such that $f(n) \leq c g(n)$ for all $n \geq n_0$ and some fixed $n_0$, and $f(n)=\Theta(g(n))$ means that $f(n) = O(g(n))$ and $g(n) = O(f(n))$. 
$f(n)=o(g(n))$ means that $\lim_{n\to\infty}f(n) / g(n) = 0$, and $f(n)=\Omega(g(n))$ that $\lim_{n\to\infty}f(n) / g(n) > 0$. 
\medskip

\noindent\textbf{Our asymptotic regime:} Although not explicitly denoted, $s_j$, $p_j$ and $\beta_{\min,j}^*$ are functions of $n$, and so are $s$, $p$ and $\beta_{\min}^*$. We study asymptotic regimes where \\[5pt]
     (A1)\quad $n \,\to \, \infty$, and the number of blocks $b$ is constant. \\[5pt]
     (A2)\quad For all $j$, $p_j-s_j \,\to \, \infty$.\\[5pt]
Assumption A1 can be slightly relaxed to allow $b$ to grow slowly with $p$, but we assume a constant $b$ for simplicity.
Similarly, A2 is a mild assumption that can be relaxed but simplifies the exposition. We provide results for both the case where the $s_j\geq1$ are fixed and $s_j \rightarrow \infty$.
 At this point we do not make any assumption on the asymptotic regime linking $n$ and $p$, but our main interest is in $n= o(p)$ settings.
%Assumption A2 implies that $p-s$ goes to infinity but we do not make any assumption on the asymptotic regime linking $n$ and $p$. We provide results for both cases $s_j\geq1$ fixed or diverging.

\section{Variable selection via informed block penalization}\label{sec:intro_blockpen}

Consider a set of candidate models $\M$ given by subsets $M\subseteq V$ and their corresponding coordinate subspaces
$$
\mathcal{L}_M\;:=\;\{\beta\in \R^p: \beta_i=0\mbox{ if }i\notin M\}.
$$
Consider a standard $\ell_0$ selection procedure with penalty $\eta(M)$ that depends on $M$ linearly through its cardinality, $\eta(M)=\kappa |M|$ for some $\kappa>0$. The selected model is
\begin{equation}\label{eq:stdl0}
\hat{S}= \arg\max_{M\in \M} \left\{\max_{\beta\in \mathcal{L}_M}\ell(\bs y; \bs \beta)\;-\;\kappa |M|\right\},
\end{equation}
where $\ell(\bs y; \bs \beta)$ is the log-likelihood function. 
We study a block-informed procedure that penalizes differently each block $B_1,\ldots,B_b$,
\begin{equation}\label{eq:Mhat}
  \hat S^b\;\in\;\arg\max_{M\in \M} \left\{\max_{\beta\in \mathcal{L}_M}\ell(\bs y; \bs \beta)-\sum_{j=1}^b\kappa_j |M_j| \right\},
\end{equation}
where $\kappa_1>0,\ldots,\kappa_b>0$. We denote a model by $M=M_1\cup\cdots\cup M_b$, and $M_j\subseteq B_j$ represents the  selected parameters in block $j$.
Note that $\hat{S}^1$ is the standard $\ell_0$ selector in \eqref{eq:stdl0}.

\section{Gaussian Sequence Model}\label{sec:properties_orth}

We now introduce the sequence model, its connection to orthogonal linear regression and discuss that $\hat S^b$ simplifies to performing hard-thresholding. %performs block-wise hard-thresholding in that setting. 
We also provide the variable selection properties of $\hat S^b$, discuss its benefits in lessening the conditions for consistent variable selection and in its convergence rates, and illustrate said advantages under different asymptotic regimes.

In this section, in addition to A1--A2, we also assume\\[5pt]
%we study the selection properties of $\hat{S}^b$ under assumption \\[5pt]
(A3)\quad $\text{for all } j, \;\; s_j \,\to \, \infty$.\\[5pt]
That is, we assume a diverging number of active signals $s_j$ in every block. %Alternatively, the $s_j$'s can be assumed to be finite. 
We present here results for $\hat S^b$ only under A3 to simplify the exposition, and because we derive a new necessary condition on signal strength under that assumption. We obtained analogous results when the $s_j$'s are finite (Supplement in Section~\ref{sec:appendixGSMsfinite}).

\subsection{Sequence model and thresholding}\label{sec:GSMgen}

The Gaussian sequence model assumes:
\begin{equation}\label{eq:GSM}
	\bs y \;=\; \sqrt{n} \bs \beta^* + \bs\epsilon,\qquad \bs \epsilon \sim N(0,\sigma^2 I_p),
\end{equation}
where $\bs y \in \mathbb{R}^p$ and without loss of generality we set $\sigma^2=1$ to streamline notation. It is a popular simplified model for high-dimensional inference, and it is also central to non-parametric statistics. See \cite{johnstone}, Chapter 3 for a complete introduction.

Let $\tilde{\bs \beta}=\frac{1}{\sqrt{n}}\bs y \sim N(\bs \beta^*,\tfrac{1}{n}I_p)$ be the MLE under \eqref{eq:GSM}. The next proposition shows that, by orthogonality, the block-wise $\ell_0$ penalty reduces to hard-thresholding $\tilde{\bs \beta}$, with block-dependent thresholds.
\begin{prop}\label{prop:isthresh}
	In the sequence model \eqref{eq:GSM}, let $\hat{S}^b$ and $\kappa_1,\ldots,\kappa_b$ defined in \eqref{eq:Mhat}. Then $\hat S^b=\hat S^b_1\cup \cdots \cup \hat S^b_b$, where, for each $j=1,\ldots,b$,
	$$
	\hat S_j^b\;:=\;\left\{i\in B_j:\; |\tilde{\beta}_i|> \sqrt{\frac{2\kappa_j}{n}}\right\}.
	$$ 
\end{prop}
Other popular penalties
%Most popular variable selection procedures 
also take this block-based thresholding form in \eqref{eq:GSM} with non-exchangeable penalization. It is the case of the LASSO and adaptive LASSO when letting the penalty vary by block: see Lemma \ref{lem:l1isthreshold} %in Auxiliary results 
in the Supplement. It is also the case of Bayesian procedures under most standard priors \citep{davidomiros} where one sets a different prior inclusion probability in each block. Thus, equivalently, we study generic hard-thresholding model selectors of the form
\begin{equation}\label{eq:thresholdS}
\hat S_j^{b}\;:=\;\{i\in B_j: |\tilde{\beta}_i|>\tau_j\}, \qquad\bs \tau=(\tau_1,\ldots,\tau_b)\in \R^b_{>0}.	
\end{equation}

Results in this section also generalize to orthogonal linear regression with normalized columns, i.e. when $\bs X$ in \eqref{eq:linearmodel} satisfies $\bs X^\top \bs X=n I_p$. In that setting the MLE is also distributed $N(\bs \beta^*,\tfrac{1}{n}I_p)$ and by orthogonality, variable selection with $\hat{S}^b$ also amounts to block thresholding.

\subsection{Selection based on block thresholds}
\label{ssec:block_thresholds}


We now study the statistical performance of the block thresholding operator in \eqref{eq:thresholdS}. Block thresholding was previously studied in the context of parameter estimation for wavelet-based models with equally sized blocks; see \cite{johnstone}, Chapters 7 to 9.
We focus here on its variable selection properties, for arbitrarily-sized blocks. We successively analyze properties relative to false positives, false negatives, exact recovery and rates of convergence. 
\medskip

\noindent\textbf{False positives.} 
A false positive
occurs when there exists a block $j=1,\ldots,b$ and $i\in B_j$ such that $\beta^\star_i=0$
but $|\tilde \beta_i|>\tau_j$. In other words, $\hat S^{b}\not\subseteq S$. 
The next proposition shows that assumption A4 is sufficient, and essentially necessary,
%A sufficient assumption 
to guarantee that asymptotically $\hat S^{b}\subseteq S$.\\[5pt]
(A4)\quad For all $j=1,\ldots,b$ and every sufficiently large $n$,\;\; $\tau_j \;\geq\;  \sqrt{\frac{{2\ln(p_j-s_j)}}{{n}}}$.
\\[5pt]
%The next proposition shows that A4 is also essentially necessary.
\begin{prop}\label{prop:suffnecIblockthresh} In the
sequence model \eqref{eq:GSM}, assume A1 and A2.
\begin{longlist}
    \item If A4 holds, then 
$\lim_{n \rightarrow \infty} P(\hat{S}^{b}\subseteq S) = 1$.\\
\item If for some $j\in\{1,\ldots,b\}$\; $
\lim_{n\to \infty} \frac{\sqrt{n}\tau_j}{\sqrt{2\ln(p_j-s_j)}} \;<\;1$,
then $\lim_{n \rightarrow \infty} P(\hat{S}^{b}\subseteq S ) <1$.
\end{longlist}
\end{prop}
By Proposition~\ref{prop:suffnecIblockthresh},
a necessary assumption to have no false positives asymptotically
is
\begin{equation}\label{cond:necIallblockthresh}
\lim_{n\to \infty} \frac{\tau_j}{\sqrt{\tfrac{2\ln(p_j-s_j)}{n}}}\;\geq\;1\qquad\mbox{for all }j=1,\ldots,b.	
\end{equation}
Proposition~\ref{prop:suffnecIblockthresh} extends to $\hat{S}^{b}$ results known for the standard hard-thresholding $\hat{S}$ in orthogonal linear regression \cite{ slope, buhlmannvandegeerbook, wainwrighthighdimbook} and in the sequence model, \cite{ariascastro} and \cite{johnstone}, Chapter 3.
\bigskip

\noindent\textbf{False negatives.} 
False negatives occur when $\hat S^{b}\not\supseteq S$, i.e. for some block $j \in \{1,\ldots,b\}$ and parameter $i\in B_j$,  we have that $\beta^\star_i\neq 0$ and $|\tilde\beta_i|\leq \tau_j$.
Control of false negatives in a block $B_j$ requires bounding the minimum of $|S_j|=s_j$ independent Gaussians with non-zero means $\beta^*_i$, $i\in S_j$. This can be achieved by setting a betamin condition
on $\beta_{\min,j}^*\;:=\;\min\{|\beta^\star_i|: i\in S_j\}$.
%, and a assumption on the thresholds $\tau_j$ for $j=1,\ldots,b$.  
Intuitively, to have no false negatives (also known as screening property), each $\tau_j$ needs to be smaller than $\beta^*_{\min,j}$.
More precisely, a sufficient assumption is: \\[5pt]
(A5)\quad For all $j=1,\ldots,b$ and sufficiently large $n$, \;
	 $\beta_{\min,j}^* - \tau_j \;\geq\; \sqrt{\frac{2\ln(s_j)}{n}}$.

\begin{prop}\label{prop:suffnecIIblockthresh}
 In the sequence model \eqref{eq:GSM}, assume A1--A3 and A5. Then
$\lim_{n \rightarrow \infty} P(\hat{S}^{b}\supseteq  S)= 1$.
\end{prop}
Assumption A5 was proved to be sufficient for selection consistency in orthogonal linear regression in \cite{wainwrighthighdimbook} for standard thresholding. Proposition~\ref{prop:suffnecIIblockthresh} generalizes the result to block-thresholding. The study of {\it necessary} assumptions
is more complicated. Proposition \ref{prop:necIIblockthresh} below provides first a necessary condition valid for any vector $\bs \beta^*$, and second a tight condition for a worst-case where $\beta_i^*=\beta_{\min}^*$ for all $i\in S$.
\begin{prop}\label{prop:necIIblockthresh}In the
sequence model \eqref{eq:GSM}, assume A1 and A2.
\begin{longlist}
    \item Suppose that $\tau_j<\beta^\star_{\min,j}$ satisfies \eqref{cond:necIallblockthresh} for some $j \in \{1,\ldots,b\}$. If \begin{equation}\label{cond:necIIallblockthresh2}
    \lim_{n\to \infty} \sqrt{n}(\beta_{\min,j}^*-\tau_j)\; <\; \infty
\end{equation}
then $\lim_{n \rightarrow \infty} P(\hat S^{b}\supseteq S  ) < 1$.
\item Suppose there exists $j\in \{1,\ldots,b\}$ such that  $\beta^*_i=\beta_{\min,j}^*$ for all $i \in S_j$, $\tau_j<\beta^\star_{\min,j}$ satisfies \eqref{cond:necIallblockthresh}, and  $s_j/p_j<1$. If A3 holds and, in addition, 
\begin{equation}\label{cond:necIIallblockthresh}
\lim_{n\to \infty} 	\frac{\beta_{\min,j}^*-\tau_j}{\sqrt{\frac{\pi}{2}\frac{\ln(s_j)}{n}}} \;\leq\;  1\end{equation}
then $\lim_{n \rightarrow \infty} P(\hat S^{b}\supseteq S ) < 1$.
\end{longlist}\end{prop}
\cite{abraham2023sharp} showed the necessity of \eqref{cond:necIIallblockthresh2} for consistency in the single-block case
in more restrictive asymptotic regimes satisfying $p/s \to \infty$.
Proposition~\ref{prop:necIIblockthresh} (ii) gives a new precise necessary  rate of growth of $\sqrt{n}\left(\beta_{\min}^*-\tau\right)$ for a worst case in which all signals are equal. It shows that in that case sufficient assumption~A5 is essentially necessary up to a constant factor (replacing 2 by $\sqrt{\pi/2} \approx 1.25$), and that A5 cannot be relaxed in general.
\bigskip


\noindent\textbf{Exact model recovery guarantees.} 
Combining assumption A4 (Proposition~\ref{prop:suffnecIblockthresh}) and A5 (Proposition~\ref{prop:suffnecIIblockthresh}),
we asymptotically recover the correct support if for all $j=1,\ldots,b$,
\begin{equation*}
  \sqrt{\frac{2\ln(p_j-s_j)}{n}}\;\;\leq\;\;\tau_j\;\;\leq\;\;\beta_{\min,j}^*-\sqrt{\frac{2\ln(s_j)}{n}}
\end{equation*}
for sufficiently large $n$. This requires that for all $j=1,\ldots,b$,
\begin{equation}\label{eq:betamin}
    {\beta_{\min}^*}_j \geq \sqrt{\frac{2\ln(p_j-s_j)}{n}} + \sqrt{\frac{2\ln(s_j)}{n}}
\end{equation}
By Propositions~\ref{prop:suffnecIblockthresh} and~\ref{prop:suffnecIIblockthresh} signals as small as the right-hand side in \eqref{eq:betamin} can be recovered in block $B_j$. Moreover, Propositions~\ref{prop:suffnecIblockthresh} (ii) and \ref{prop:necIIblockthresh} show that assumptions A4--A5 are essentially necessary. It follows that betamin condition \eqref{eq:betamin} is near-necessary for selection consistency, in the following sense.
\begin{lem}\label{lem:necbetaminblockthresh}
    If there exists a block $j\in\{1,\ldots,b\}$ with equal non-zero parameter values $\beta^*_i=\beta_{\min,j}^*$ for all $i \in S_j$, where $s_j/p_j<1$ and
    \begin{equation}\label{eq:necbetaminblockthresh}
      \lim_{n\to \infty}\frac{\beta_{\min,j}^*}{\sqrt{\tfrac{2\ln(p_j-s_j)}{n}}+\sqrt{\tfrac{\pi}{2}\tfrac{\ln(s_j)}{n}}}\;<\; 1 .
    \end{equation}
Then under assumptions A1--A3,\;  $\lim_{n \rightarrow \infty} P(\hat{S}^{b} = S) < 1$.
\end{lem}



\bigskip

\noindent\textbf{Rate of convergence.} From the
preceding results, we can bound the rate of convergence of $P(\hat S^{b}\neq S)$.
\begin{thm}\label{theo:fullrecovGS}
Assume A4 and A5. Then, for sufficiently large $n$,
\begin{equation} \label{eq:rateconvgsm}
P(\hat{S}^{b} \neq S)\;\leq\;\sum_{j=1}^b\,\,\frac{e^{-\tfrac{n}{2}\big[\tau_j^2-\tfrac{2\ln(p_j-s_j)}{n}\big]}}{\sqrt{\pi\ln(p_j-s_j)}}+\frac{e^{-\tfrac{n}{2}\big[(\beta^\star_{\min,j}-\tau_j)^2-\tfrac{2\ln(s_j)}{n}\big]}}{\sqrt{\pi\ln(s_j)}}.
\end{equation}
 Further, %On the other hand, 
if $\beta^\star_{\min,j}$ satisfy \eqref{eq:betamin} for every block $j$, and the thresholds $\tau_j$ take oracle values
	\begin{equation} \label{eq:optimtauj}
\tau^\star_j \;=\; \frac{1}{2}\beta^\star_{\min,j}+\frac{1}{n\beta^\star_{\min,j}}(\ln({p_j-s_j})-\ln({s_j})),
\end{equation}
then A4 and A5 hold and, for sufficiently large $n$,
\begin{equation}\label{eq:optimrateconvgsm}
    P(\hat{S}^{b} \neq S) \; \leq \; 2 \sum_{j=1}^be^{-[\frac{n}{8}{\beta^\star_{\min,j}}^2-\max\{\ln(p_j-s_j),\ln(s_j)\}]}.
\end{equation}
\end{thm}
In \eqref{eq:rateconvgsm} the first term of each summand decreases exponentially in $n\tau_j^2/2$, while the second term decreases exponentially in $n(\beta_{\min,j}^*-\tau_j)^2/2$. The choice $\tau^*_j$ ensures that both terms are equal and hence approximates the values minimizing \eqref{eq:rateconvgsm}. We refer to these ideal $\tau^*_j$ as {\it oracle thresholds} because they depend on quantities $s_j$ and $\beta_{\min,j}$ that are unknown in practice. 
The resulting bound in \eqref{eq:optimrateconvgsm} is tightest in the worst-case scenario where $\beta_i^*=\beta_{\min}^*$ for all $i \in S$, in that it approximates the fastest rate achievable in that worst case.

\subsection{Benefits of block penalties}\label{sec:prosconsblock}


We now examine the benefits of block penalization. We start with a simple example to build intuition and then discuss
two types of benefit: softening the %those relative to 
conditions for model selection consistency and 
improving the %those relative to 
rates of convergence.
%\bigskip

Intuitively, block penalties allow adapting 
to varying sparsity $s_j$  or signal strengths $\beta^\star_{\min,j}$ across blocks. 
When the number of inactive parameters $p_j-s_j$ is small one may use smaller thresholds $\tau_j$ and still control false positives, and when $\beta^\star_{\min,j}$ is large one may use larger $\tau_j$ and still control false negatives.
As the blocks become more discriminative, that is, more distinct $p_j-s_j$ and $\beta^\star_{\min,j}$ across blocks, the larger the potential benefits.
Consider a somewhat extreme example with $b=2$ blocks, a large sparse block with $s_1=0$, $p_1-s_1=80$ and a smaller non-sparse block with $s_2=10$, $p_2-s_2=10$, with $n=1$ and $\beta_i^\star=4$ for all $i\in S$. \color{black}
%Consider a somewhat extreme example where $n = 1$, $p = 100$, %$s = 10$  and $\beta_i^\star=4$ for all $i\in S$. There are $b=2$ blocks: a large sparse block with $s_1=0$, $p_1-s_1=80$ and a smaller non-sparse block with $s_2=10$, $p_2-s_2=10$. 
In the left panel of Figure~\ref{fig:starter}, a standard threshold selector is used with $\tau=\sqrt{2 \ln(p-s)/ n}$, the minimal threshold such that the FWER vanishes (dashed line). The large number of truly zero parameters $p-s=90$ forces a high threshold, which leads to a false negative. In the right panel, block thresholds are set to $\tau_j=\sqrt{2 \ln(p_j-s_j)/ n}$.
Adjusting $\tau_2$ to the lower number of truly zero means in block 2 ($p_2-s_2$) allows recovering the true model. 
\begin{figure}[ht]
\begin{minipage}{1\columnwidth}
\begin{tabular}{M{62mm}M{62mm}}
        Standard thresholding & Block thresholding \\
         \includegraphics[scale=1]{images/p.intro.std.pdf} & \includegraphics[scale=1]{images/p.intro.block.pdf} \\
\end{tabular}
\end{minipage}
\caption{In both panels, the 90 white dots are generated from $N(0, 1)$ and the 10 black dots are generated from $N(4, 1)$. Left panel: the dashed line is $\sqrt{2 \ln (p-s)/n}$, the minimal threshold such that the FWER vanishes. Right panel: observations are grouped in two blocks. The dashed lines are the minimal block thresholds such that the FWER vanishes}
\label{fig:starter}
\end{figure}
\bigskip

\noindent\textbf{Conditions for variable selection consistency.} In the example, the possibility to lower thresholds in the less sparse blocks with $\hat S^2$ allowed recoverying the smallest active signal. % and the true model. 
More precisely, assumptions A4--A5 give ranges of threshold valuesthat are sufficient and essentially necessary for asymptotic support recovery. For the standard threshold selector $\hat{S}$, the range for the single threshold $\tau$ is
\begin{equation}\label{eq:rangethreshex}
  \sqrt{\frac{2\ln(p-s)}{n}}\;\;\leq\;\;\tau\;\;\leq\;\;\beta^*_{\min}-\sqrt{\frac{2\ln(s)}{n}}.	
\end{equation}
For the block threshold selector $\hat{S}^{b}$, 
%assuming without loss of generality that the global minimum $\beta_{\min}^*$ is in block $b$, 
the ranges for the $\tau_j$'s are
    \begin{eqnarray}\label{eq:rangeblockthreshex}
  \sqrt{\frac{2\ln(p_j-s_j)}{n}}\;\;&\leq\;\;\tau_j\;\;\leq&\;\;\beta^*_{\min,j}-\sqrt{\frac{2\ln(s_j)}{n}}   .
%  \sqrt{\frac{2\ln(p_1-s_1)}{n}}\;\;&\leq\;\;\tau_1\;\;\leq&\;\;\beta^*_{\min,1}-\sqrt{\frac{2\ln(s_1)}{n}}\nonumber\\
%  &\;\;\vdots \\
%  \sqrt{\frac{2\ln(p_b-s_b)}{n}}\;\;&\leq\;\;\tau_b\;\;\leq&\;\;\beta^*_{\min}-\sqrt{\frac{2\ln(s_b)}{n}} \nonumber.       
    \end{eqnarray}

The ranges in \eqref{eq:rangeblockthreshex} imply that $\hat{S}^{b}$ requires milder conditions for selection consistency than $\hat{S}$. 
%By definition, the ranges in \eqref{eq:rangeblockthreshex} are wider than the range in \eqref{eq:rangethreshex}, making conditions to get selection consistency milder with $\hat S^b$. 
Intuitively, if for some $\tau_k$ and $\tau_l$ the ranges do not overlap, 
then a global threshold $\tau$ cannot possibly satisfy \eqref{eq:rangeblockthreshex} and consistent selection is essentially not possible.
For example, this occurs if $\beta_{\min}^* - \sqrt{(2/n) \ln(s_b)} < \sqrt{(2/n) \ln(p_1-s_1)}$.
More precisely, Corollary \ref{cor:selconstimpossible} gives %We next give precise 
conditions under which consistent selection is possible with $\hat{S}^b$ but not with $\hat{S}$.
%, that follow from Propositions~\ref{prop:suffnecIblockthresh} and \ref{prop:necIIblockthresh}.

\begin{cor}\label{cor:selconstimpossible}
Assume A1--A5, that $s/p<1$, and $\beta^*_i=\beta_{\min}^*$ for all $i \in S$. If
\begin{eqnarray*}
    &\lim_{n\to \infty}\frac{\beta_{\min}^*}{\sqrt{\tfrac{2\ln(p-s)}{n}}+\sqrt{\tfrac{\pi}{2}\tfrac{\ln(s)}{n}}}\;<\; 1 
\end{eqnarray*}
then 
$\lim_{n \rightarrow \infty} P(\hat{S} = S) < 1$ and $\lim_{n \rightarrow \infty} P(\hat{S}^{b} = S)=1$.
\end{cor}

% for $\hat S^b$ and $\hat S$ depends on the chosen thresholds and decrease as thresholds shrink. However, thresholds also need to be large enough to avoid false positives. To make betamin conditions comparable and understand the complete potential gains with $\hat S^b$,  $\tau_j=\sqrt{2\ln(p_j-s_j)/n}$ for all $j$ for $\hat S^b$, $\tau=\sqrt{2\ln(p-s)/n}$ for $\hat S$. This choice allows to precisely assess what signals can or cannot be recovered with the two selectors. 

% For $\hat S$, the betamin condition A5 constrains only the global $\beta_{\min}^*$ which solely determines whether signal strength is large enough for variable selection consistency. To under , we compare next the constraints that selection consistency with $\hat S$ and $\hat S^b$ impose on that $\beta_{\min}^*=\beta_{\min,b}^*$. For simplicity we assume A5 is met for $\hat S^b$ in blocks other than $b$. 



% \paul{You had asked the question below. I have added more explanation below. The idea of this section is to analyze the smallest signal recoverable, the constraints on the overall $\beta_{\min}^*$}
% \david{Why is it enough to look only at block $b$ ? I mean, even if $\beta_{\min}^\star$ is in block $b$, $(p_j,s_j)$ from other blocks could be such that \eqref{eq:betamincon.suffblockthresh} holds in block $b$, but not for some other block $j$ no?}\piotr{Yes, this should be clarified.}

% From \eqref{eq:betamin}, we obtain the sufficient betamin condition for $\hat{S}$:
% \begin{equation}\label{eq:betamincon.suffthresh}
%     \beta_{\min}^* \;>\; \sqrt{\frac{2\ln(p-s)}{n}}+\sqrt{\frac{2\ln(s)}{n}}
% \end{equation} 
% For the block threshold selector $\hat{S}^{b}$, since $\beta_{\min}^*$ is located in block $b$, assumption \eqref{eq:betamin} requires:

% \begin{equation}\label{eq:betamincon.suffblockthresh}
%     \beta_{\min}^* \;>\; \sqrt{\frac{2\ln(p_b-s_b)}{n}}+\sqrt{\frac{2\ln(s_b)}{n}}
% \end{equation} 
% We then set: 
% \begin{equation}\label{eq:betaminbardef}
%     \underline{\beta}_{\min,orth}^{*} := \sqrt{\frac{2\ln(p-s)}{n}} + \sqrt{\frac{2\ln(s)}{n}},\qquad
%     \underline{\beta}^{*,b}_{\min,orth} := \sqrt{\frac{2\ln(p_b-s_b)}{n}} + \sqrt{\frac{2\ln(s_b)}{n}}
% \end{equation} Since $p_b-s_b \leq p - s$ and $s_b \leq s$, $\underline{\beta}^{*,b}_{\min,orth} \leq \underline{\beta}_{\min,orth}^{*}$ and assumption \eqref{eq:betamincon.suffblockthresh} is always milder than \eqref{eq:betamincon.suffthresh}. Selector $\hat{S}^{b}$ is then able to recover smaller signals than $\hat{S}$. The magnitude of the difference depends on how $p-s$ and $s$ compare to $p_b-s_b$ and $s_b$. The gain thus depends on whether the block structure induced by the external information is indicative of varying sparsity. \david{I didn't find sub-section "Minimal betamin conditions" super-informative (given what we already said earlier). This is however used later for Figure Figure~\ref{fig:ratiobetamincond}, so maybe we can introduce $\underline{\beta}_{\min,orth}^{*}$ later, when presenting that example.}\piotr{agreed}

\bigskip

\noindent\textbf{Rates of convergence.}   We just saw that block penalties can attain consistent model selection where standard penalties cannot. We now discuss differences in the probability of correct selection, when consistent selection is possible for both procedures. We assume that block thresholds 
are set to their oracle values $\tau_j^\star$ 
and similarly that the single threshold
takes its oracle value $\tau^\star$.
Let ${OR}^b_{orth}$ be the oracle convergence rate for $\tau_j^\star$  in \eqref{eq:optimrateconvgsm} and ${OR}_{orth}$ that for $\tau^\star$.
Then
\begin{equation}\label{eq:ratioCR}
    \frac{{OR}^b_{orth}}{{OR}_{orth}}\;=\;\sum_{j=1}^be^{-\frac{n}{8}\left({\beta^\star_{\min,j}}^2-{\beta_{\min}^\star}^2\right)-(\max\{\ln(p-s),\ln(s)\}-\max\{\ln(p_j-s_j),\ln(s_j)\})}
\end{equation}
Since $\beta_{\min,j}^\star \geq \beta_{\min}^\star$ and $\max\{\ln(p-s),\ln(s)\}\geq\max\{\ln(p_j-s_j),\ln(s_j)\}\big)$, equation \eqref{eq:ratioCR} shows that for $n$ large enough ${OR}^b_{orth} \leq {OR}_{orth}$ for any partition in blocks, i.e. the convergence rate for $\tau_j^\star$ is never worse than for $\tau^\star$. The magnitude of the gain depends on the discriminative power of the blocks.
For any sparse setting where $s_j<p_j-s_j$ for every $j$, 
%\paul{You added "mildly" sparse, I removed it. It does not matter if it is highly or midly sparse. By sparse I just mean $s_j<p_j-s_j$ in every block to know what is $\max\{\ln(p_j-s_j),\ln(s_j)\}$ and simplify notation}, 
assuming without loss of generality that $\beta_{\min,b}^\star=\beta_{\min}^\star$, we have
\begin{equation}\label{eq:goodorthratio}
    \frac{{OR}^b_{orth}}{{OR}_{orth}}\;=\; \frac{p_b-s_b}{p-s} \,+\,\sum_{j<b}\frac{p_j-s_j}{p-s}e^{-\frac{n}{8}\left({\beta^\star_{\min,j}}^2-{\beta_{\min}^\star}^2\right)}
\end{equation}
%\paul{this ratio goes to 1 if the betamins are equal across blocks and for some $j$, $p_j-s_j/p-s\to1$ (happens in our Examples). I guess this is due to simplifications in the bounding.}
There are then two sources of improvement in convergence rates: in any block $j$, a smaller number of truly zero means, that is $p_j-s_j< p-s$, or a larger smallest signal, $\beta^\star_{\min,j} > \beta_{\min}^\star$. Depending on the sparsity of the block, one of the two effects usually prevails. In blocks where $\beta^\star_{\min,j} \approx \beta_{\min}^\star$, improvement comes mainly from having a smaller number of truly zero means. In those blocks
one sets $\tau^\star_j < \tau^\star$, so that one may detect smaller signals. 
If
$\beta^\star_{\min,j} \gg \beta_{\min}^\star$,
one sets $\tau^\star_j > \tau^\star$
and the probability of false positives is reduced.

On the contrary, there is no improvement 
(${OR}^b_{orth}={OR}_{orth}$) 
when, for example, %as before $p_j-s_j>s_j$, but 
$p_j-s_j=(p-s)/b$ 
 and $\beta^\star_{\min,j}=\beta_{\min}^\star$ for all $j$. That is, one does not expect benefits when the number of non-zero signals and their magnitudes are similar across blocks. \color{black}
% the external information fails to identify blocks that are substantially differently sparse. 

\subsection{Examples with two blocks}\label{sec:examples}

We illustrate our results with four concrete examples contemplating different sparsity regimes and informativeness of the blocks, summarized in Table~\ref{tab:examples}. We focus on a sparse setting where $\ln(s)=o(\ln(p-s))$, with two blocks ($b=2$). We assume that $\beta_{\min}^*$, the smallest non-zero $|\beta_i^*|$ is located in block $2$. 

\begin{table*}[ht]
\caption{Sparsity and block assumptions in Examples 1--4}
\label{tab:examples}
\begin{tabular}{@{}cccccccc@{}}
\hline \\[-1em]
Example & $p-s$ & $p_1-s_1$ & $p_2-s_2$ & $s_1$ 
& $s_2$\\[3pt]
\hline\\[-1em]
$1$    & $\frac{3}{2}\,n$ & $\frac{3}{2}\,n-\sqrt{n}$ & $\sqrt{n}$  & $3\ln(n)/2$  & $3\ln(n)/2$ \\[3pt]
$2$    & $e^{n/20}$ & $e^{n/20}-n^2$ & $n^2$  & $3\ln(n)/2$  & $3\ln(n)/2$   \\[3pt]
$3$    & $n$ & $n-\ln(n)$ & $\ln(n)$  & $3\ln(n)/2$  & $3\ln(n)/2$   \\[3pt]
$4$    & $n$ & $n/2$ & $n/2$  & $3\ln(n)/2$  & $3\ln(n)/2$ \\[3pt]
\hline
\end{tabular}
\end{table*}

In Examples 1 to 3, external information is discriminative as it singles out a block $B_1$ that is %such that $\bs\beta^*_{B_2}$ is significantly 
sparser than block $B_2$. 
%than the global $\bs\beta^*$. 
In Example 1 the blocks are %external information is 
moderately discriminative, in that $(p-s)/(p_2-s_2)$ is a power of $n$. In Example 2 they are highly discriminative, since $(p-s)/(p_2-s_2)$ is exponential in $n$ and in Example 3 it is highly discriminative in that $\bs\beta^*$ is sparse overall but it is non-sparse within block 2 ($s_2 > p_2 - s_2)$. %but $\bs\beta^*_{B_2}$ is non-sparse ($\ln(p_2-s_2)=O(\ln(s_2))$ and $s_2>p_2-s_2$). 
Example 3 also differs from 1 and 2 in that it is less sparse. %set in only very mildly sparse settings. 
Example 4 is a non-discriminative \color{black}
%In Example 4, external information is non-discriminative, the block partition is an average
random guess where each block has half of the inactive parameters. %variables. 
%\paul{It is not super important but I think that what matters in terms of usefulness of the external information is $p-s/p_2-s_2$ but David has been stressing $p_1-s_1/p_2-s_2$ in his edits. My point of view is that one the main gains with $\hat S^b$ is that you can recover smaller signals in the blocks that are less sparse than the global set of variables. If you have a three block situation such: $p-s=e^n$, $p_1-s_1=e^n-n^2-n$, $p_2-s_2=n^2$ and $p_3-s_3=n$, $p_2-s_2/p_3-s_3$ is not telling of the gain with $\hat S^b$. To illustrate how much you can \emph{increase} penalties in some blocks, to me, the relevant ratio would then be $\beta_{\min,j}/\beta_{\min}$. }

%We then illustrate the benefit of block penalties in terms of rate of convergence with a plot of the ratio ${OR}^b_{orth}/{OR}_{orth}$.

\bigskip
\noindent\textbf{Selection consistency.}
%\noindent\textbf{Valid thresholds for $\hat S$ and $\hat S^b$.}
%To illustrate the benefits in %terms of conditions for selection consistency, 
Figure \ref{fig:boundthresh} plots the range of threshold values $\tau$ and $(\tau_1,\tau_2)$ ensuring selection consistency with $\hat{S}$ and $\hat{S}^{b}$ given in \eqref{eq:rangethreshex} and \eqref{eq:rangeblockthreshex},
%Figure \ref{fig:boundthresh} plots the upper and lower bounds for $\tau$ as in \eqref{eq:rangethreshex} and $\tau_1$ and $\tau_2$ as in \eqref{eq:rangeblockthreshex} 
in the four examples assuming $\beta^*_{\min,1}=2/3$  and $\beta^*_{\min,2}=\beta^*_{\min}=1/10$. 
In Example 2, asymptotic recovery with $\hat{S}$ is not possible in the range of values of $n$ considered while it is with $\hat{S}^{b}$.
In the other examples, asymptotic recovery with $\hat{S}$ is possible but it requires larger $n$ than with $\hat{S}^b$.
Example 1 shows that the gain can be large even when blocks are moderately discriminative, and Example 3 when some blocks are non-sparse. 
\color{black}
%Observe that, for $\tau$, the sparser $\bs \beta^*$ is, the larger the $n$ needed to make the upper bound to exceed the lower bound, i.e. to essentially make support recovery possible. 
%The block threshold $\tau_2$ requires smaller $n$, with the gap with respect to $\tau$ increasing as the difference in sparsity between block 2 and $V$ increases. 
In Example 4 when external information is non-discriminative, the required $n$ is the about the same for $\hat{S}^{b}$ and $\hat{S}$. 
%In Example 2 on the contrary, asymptotic recovery with $\hat{S}$ is not possible in the range of values of $n$ considered while it is with $\hat{S}^{b}$. Indeed, in that example, $\sqrt{\frac{2\ln(p-s)}{n}}+\sqrt{\frac{\ln(s)}{n}} \to \sqrt{1/10} \approx 0.32$ and $\sqrt{\frac{2\ln(p_2-s_2)}{n}}+\sqrt{\frac{\ln(s_2)}{n}} \to 0$. 
%Example 3 shows that the gain with $\hat{S}^{b}$ can be large even in settings where $V$ is only somewhat sparse.

\begin{figure}[ht]
\begin{minipage}{1\columnwidth}
\setlength{\tabcolsep}{1pt}
    \begin{tabular}{M{4mm}M{63mm}M{63mm}}
       \toprule
         & Example 1 & Example 2 \\
        \midrule
        $\tau$ & \includegraphics[scale=1]{images/taus/p60mm.scen1tau.pdf} & \includegraphics[scale=1]{images/taus/p60mm.scen2tau.pdf} \\
        $\tau_1$ & \includegraphics[scale=1]{images/taus/p60mm.scen1tau1.pdf} & \includegraphics[scale=1]{images/taus/p60mm.scen2tau1.pdf} \\
        $\tau_2$ & \includegraphics[scale=1]{images/taus/p60mm.scen1tau2.pdf} & \includegraphics[scale=1]{images/taus/p60mm.scen2tau2.pdf} \\
        \midrule
        & Example 3 & Example 4 \\
        \midrule
        $\tau$ & \includegraphics[scale=1]{images/taus/p60mm.scen3tau.pdf} & \includegraphics[scale=1]{images/taus/p60mm.scen4tau.pdf} \\
        $\tau_1$ &\includegraphics[scale=1]{images/taus/p60mm.scen3tau1.pdf} & \includegraphics[scale=1]{images/taus/p60mm.scen4tau1.pdf} \\
        $\tau_2$ &\includegraphics[scale=1]{images/taus/p60mm.scen3tau2.pdf} & \includegraphics[scale=1]{images/taus/p60mm.scen4tau2.pdf}\\
        \bottomrule
    \end{tabular}
\setlength{\tabcolsep}{6pt}
\end{minipage}
    \caption{Valid upper bound (solid) and lower bound (dashed) on $\tau$, $\tau_1$ and $\tau_2$ in Examples 1 to 4}
 \label{fig:boundthresh}
\end{figure}

\bigskip

\noindent\textbf{Smallest signal recoverable.} 
We compare the \emph{smallest signal recoverable} by $\hat S$ and $\hat S^b$ while still being selection consistent. We denote those $\underline{\beta}_{\min,orth}^{*}$ and $\underline{\beta}_{\min,b}^{*}$ respectively. In assumption A5, we set the threshold(s) to the lowest value such that the FWER vanishes (per Proposition~\ref{prop:suffnecIblockthresh}). Since we assumed that the global minimum $\beta_{\min}^*$ is in block $2$, we get:
\begin{equation}\label{eq:betaminbardef}
    \underline{\beta}_{\min,orth}^{*} := \sqrt{\frac{2\ln(p-s)}{n}} + \sqrt{\frac{2\ln(s)}{n}},\qquad
    \underline{\beta}^{*,b}_{\min,orth} := \sqrt{\frac{2\ln(p_2-s_2)}{n}} + \sqrt{\frac{2\ln(s_2)}{n}}
\end{equation} 
%We discuss the ratio $\underline{\beta}^{*,b}_{\min,orth}/\underline{\beta}_{\min,orth}^{*}$ in our four examples. 

The left panel in Figure~\ref{fig:ratios} plots the ratio %of smallest signal recoverable 
$\underline{\beta}^{*,b}_{\min,orth}/\underline{\beta}^{*}_{\min,orth}$. %in our four examples. 
In Example 1, with discriminative blocks, the
smallest signal recoverable with $\hat{S}^{b}$ is asymptotically about 25\% smaller than with the standard selector $\hat{S}$. 
In Example 2 where blocks are even more discriminative, the ratio converges to 0. %: the smallest signal recoverable with $\hat S^b$ is infinitely smaller.
In Example 4, with non-discriminative blocks, the ratio converges to 1, i.e. the benefits in the recoverable signal fade as $n$ grows. 
Example 3 illustrates how highly discriminative blocks can also bring significant benefits in terms of signal recoverable in a regime that is somewhat sparse.
\begin{figure}[ht]
\begin{minipage}{1\columnwidth}
\begin{tabular}{M{63mm}M{63mm}}
        $\underline{\beta}^{*,b}_{\min,orth}/\underline{\beta}^{*}_{\min,orth}$ & ${OR}^b_{orth}/{OR}_{orth}$ \\
         \includegraphics[scale=1]{images/p.ratio.smallestbeta.pdf} & \includegraphics[scale=1]{images/p.ratio.orconvrate.pdf} \\
\end{tabular}
\end{minipage}
\caption{Ratio of smallest signals recoverable (left) and oracle convergence rates (right) with $\hat S^b$ and with $\hat S$ in Examples 1--4}
\label{fig:ratios}
\end{figure}

\bigskip
\noindent\textbf{Oracle convergence rates.} In the right panel of Figure~\ref{fig:ratios}, we plot the ratio ${OR}^b_{orth}/{OR}_{orth}$ in Examples 1 to 3 where blocks are discriminative. We set $\beta_{\min }^*=\beta_{\min,2}^*=1/5$ and $\beta_{\min,1}^*=1.3\beta_{\min }^*$ to guarantee the recovery is possible with both $\hat S$ and $\hat S^b$. We leave the non-discriminative Example 4 out since the ratio is simply stable at 1. Figure~\ref{fig:ratios} shows that in the three examples the convergence is much faster with $\hat S^b$. %The more discriminative the blocks the faster.

\section{High-dimensional linear regression}
\label{sec:properties_legression}

In Section~\ref{sec:properties_orth} we showed that, in the sequence model, block penalties allow consistent model recovery in settings where it is otherwise not possible (e.g. smaller signals), and improves oracle consistency rates. We now extend the results to linear regression. 
We present sufficient and necessary assumptions, 
betamin conditions and rates for the probability of correct selection for the block informed $\ell_0$-penalized selector $\hat S^b$ in \eqref{eq:Mhat}. We also compare the properties of $\hat S^b$ to those of the standard $\ell_0$ selector $\hat S$. 
Whereas the framework is more involved than in the sequence model and the required proof techniques are different, the results remain conceptually the same. 

%Here $V$ is the set of all variables and $p_j-s_j$ and $s_j$ are respectively the numbers of inactive and active variables in block $B_j$.
%\paul{I removed a sentence that said here "Here $V$ is the set of all variables and $p_j-s_j$ and $s_j$ are respectively the numbers of inactive and active variables in block $B_j$." I don't see why we need to repeat it.} 
We require Assumptions A1--A2, but the number of active variables $s_j$ may either be fixed or go to infinity. 
That is, we do not assume A3 anymore. For simplicity, we assume that $\hat{S}^b$ is unique and that one constrains attention to the set of models $\M$,
such that
$\bs X_M$ has full column rank for any $M \subset \M$, and that the data-generating true support $S$ lies in $\M$. Non-uniqueness of $\hat{S}^b$ and non-full rank models can be accommodated though, at the expense of a slightly more involved treatment.

Our analysis relies on a classical connection
between $\ell_0$ penalties and Bayesian variable selection. 
Section \ref{ssec:regression_proof_strategy} reviews this connection and the proof strategy for our results.
The technical nature of the proof precludes a detailed exposition, we instead only present the most important ideas.
In Section \ref{ssec:regression_mainresult} we state our main theorem
on sufficient conditions for variable selection consistency for $\hat{S}^{b}$ and oracle convergence rates. %as we did for the sequence model. 
To assess the tightness of said sufficient conditions,  Section~\ref{ssec:neccond_linreg} gives related {\it necessary} conditions. %assumptions and betamin conditions for variable selection consistency with $\hat{S}^{b}$. 
Section \ref{ssec:regression_gains} discusses the gains of block penalization in further detail. Section \ref{ssec:nobetaminhyp} gives a general convergence result that, when certain betamin conditions do not hold, one still has guarantees of discarding inactive parameters and detecting sufficiently large active parameters.
%for $\hat{S}^{b}$ under no betamin condition which is used in Section \ref{sec:informed_l0}.
The latter result plays an important role for our data-based procedures in Section \ref{sec:informed_l0}. \color{black}


\subsection{Proof strategy}
\label{ssec:regression_proof_strategy}


We state a well-known
reformulation of $\hat S^b$ in linear regression \eqref{eq:linearmodel}. 
\begin{lem}\label{lem:useful-reformulation} Let $\tilde{\bs \beta}^{(M)}= (\bs X_{\!\!M}^\top \bs X_{\!\!M})^{-1}\bs X_{\!\!M}^\top\bs y \in \R^{|M|}$ be the MLE under model $M\in \M$. Then $\hat S^b\;\in\;\arg\max_{M\in \M} {\rm C}(M)$, where $M=M_1\cup \ldots\cup M_b$ and
\begin{equation}\label{eq:CM}
    {\rm C}(M)\;: =\;\tfrac{1}{2}\|\bs X_M\tilde{\bs\beta}^{(M)}\|^2-\sum_{j=1}^b\kappa_j |M_j|,
\end{equation}
\end{lem}
By Lemma~\ref{lem:useful-reformulation}, $\hat S^b$ selects
$M\in \M$ with the largest score $C(M)$, or equivalently the largest normalized score
\begin{equation}\label{eq:normalizedscore}
    NC(M)\;:=\;\frac{e^{{\rm C}(M)}}{\underset{M' \in \M}{\sum} e^{{\rm C}(M')}}.
\end{equation}
This normalized score can be understood as a pseudo-posterior probability for model $M$ in a Bayesian variable selection framework \citep{BIC}.
Proposition 1 in \cite{rossell2022concentration}, reproduced below as Lemma \ref{lem:L0toL1convergence}, proves that the expectation of such posterior probabilities bounds the probability of an incorrect model selection $(\hat S^b\neq S)$. 

\begin{lem}\label{lem:L0toL1convergence}
\begin{longlist}
\item 
$P(\hat{S}^{b} \neq S) \,\,\leq \,\, 2 \sum_{M \in \M\setminus \{S\}} \E\left(NC(M)\right).$
\item  For any $M,M'\in\M$, such that $M\neq M'$, 
$NC(M)\;\leq \;\big(1+e^{C(M')-C(M)}\big)^{-1}$.
\end{longlist}
\end{lem}
We use  Lemma~\ref{lem:L0toL1convergence} (i) to show the variable selection consistency of $\hat S^b$ by guaranteeing the vanishing, in expectation, of the sum of the normalized scores. Such vanishing bears resemblance to what is known as \emph{strong selection consistency} property \citep{narisettyhe2014}: the concentration of pseudo-posterior model probabilities on the truth $S$. We prove however $L_1$ convergence instead of convergence in probability.
To bound $\E(NC(M))$, by Lemma~\ref{lem:L0toL1convergence} (ii) with $M'=S$, we may instead bound a simple function of $C(S)-C(M)$: a pairwise comparison between each $M$ and the data-generating $S$. \color{black}
This is achieved by noting that $C(S)-C(M)=\tfrac12 L_{SM}+\Delta_{MS}$, where for any two models $M,T\subseteq V$, we denote

\begin{align}\label{eq:DeltaM}
&\Delta_{MT}\;:=\;\sum_{j=1}^b \kappa_j(|M_j|-|T_j|)
\\
&L_{TM}\;:=\;\|\bs X_{T}\tilde{\bs\beta}^{(T)}\|^2-\|\bs X_{M}\tilde{\bs\beta}^{(M)}\|^2.
\nonumber
\end{align}
%and for any two models $Q,M\subseteq V$, we denote 
%$$L_{QM}\;:=\;\|\bs X_{Q}\tilde{\bs\beta}^{(Q)}\|^2-\|\bs X_{M}\tilde{\bs\beta}^{(M)}\|^2.$$
To lower bound $C(S)-C(M)$ we use that $L_{SM}$ can be expressed in terms of chi-squared variables.
For any over-fitted model $M \supset S$, we have $\Delta_{MS}>0$ and it is well-known that $-L_{SM} \sim \chi_{|M \setminus S|}^2$. Hence, if one sets large enough $\kappa_j$ (and thus $\Delta_{MS}$), then $NC(M)$ vanishes. 
For non-overfitted  models, the idea is to take the union model $Q_S=S\cup M$, and to note that $L_{SM}=L_{Q_SM}-L_{Q_SS}$. Since $M \subset Q_S$, we may use the next lemma from \cite{rossell2022concentration} to bound $L_{Q_SM}$ and hence also $L_{SM}$.

\begin{lem}\label{lemma:s7}
    Let $M,Q$ be any two nested models such that $M \subseteq Q$. Then
$$
L_{QM}\;=\;\|\bs X_{Q}\tilde{\bs\beta}^{(Q)}\|^2-\|\bs X_{M}\tilde{\bs\beta}^{(M)}\|^2\;\sim\; \chi_{|Q\setminus M|}^2\left(\mu_{Q M}\right), 
$$
where $\chi_{k}^2(\mu)$ denotes the noncentral chi-squared distribution with $k$ degrees of freedom and noncentrality parameter $\mu$. The noncentrality parameter $\mu_{QM}$ is given by \begin{equation}\label{eq:muQM}
    \mu_{QM}\;\;:=\;\;\|\big(I_n-P_M\big) \bs X_{Q\setminus M} \bs \beta^*_{Q\setminus M}\|^2
\end{equation}
where $P_M=\bs X_{M}\left(\bs X_{M}^{\top} \bs X_{M}\right)^{-1} \bs X_{M}^{\top}$.
\end{lem}

For example, by Lemma~\ref{lemma:s7}, if $M\subset S$ then $L_{SM}\sim \chi^2_{|S|-|M|}(\mu_{SM})$, which has expectation $|S|-|M|+\mu_{SM}$.
If the noncentrality parameter $\mu_{SM}$ is large enough, then $L_{SM}$ and $C(S)-C(M)$ are also large, and $NC(M)$ vanishes. Lemma \ref{lem:boundrho} shows that large $\mu_{Q_SM}$ can be achieved by setting a betamin condition and an eigenvalue condition on $\bs X$ involving the following quantity.  
 \begin{equation}\label{eq:rhoX}
        \rho(\bs X)\,=\,\min_{M \in \M:\, S \not \subseteq M} \lambda_{\min}\big(\tfrac{1}{n} \boldsymbol{X}_{S \setminus M}^{\top}\left(I_n- P_M\right) \boldsymbol{X}_{S \setminus M}\big) \geq 0.
    \end{equation}
The quantity $\rho(\bs X)$ is nonnegative and relates to how distinguishable the other models $M\in \M$ are from $S$ \citep{infotheowainwright}.  More specifically, $\tfrac{1}{n} \boldsymbol{X}_{S \setminus M}^{\top}\left(I_n- P_M\right) \boldsymbol{X}_{S \setminus M}$ is the sample covariance matrix of the residuals when regressing $\boldsymbol{X}_{S \setminus M}$ on $\boldsymbol{X}_M$.
%That is, $\rho(\bs X)$ measures how well one can approximate the $\boldsymbol{X}_{S \setminus M}$ with $\boldsymbol{X}_M$.
In an orthonormal case where $\bs X^{\top}\bs X=n\I_p$, then $\rho(\bs X)=1$.  

\begin{lem}\label{lem:boundrho}
 For any $M\in\M$, let $Q_S=M \cup S$. Then the non-centrality parameter  \eqref{eq:muQM}
    \begin{equation}\label{eq:musmbound}
    \mu_{Q_SM}\; \geq \;n\,\rho(\bs X)\,\sum_{j=1}^b|S_j\setminus M_j|\,{\beta_{\min,\,j}^*}^2
\end{equation}
\end{lem}



\subsection{Sufficient conditions for consistency with block $\ell_0$ penalties}
%\subsection{Selection based on block $\ell_0$ penalty}
\label{ssec:regression_mainresult}

We state two sufficient assumptions for asymptotically recovering $S$. 
Building on our previous discussion, the first requires the block penalties $\kappa_j$ to be large enough,  
and the second is a betamin condition:\\[5pt] 
(A6) \quad For each block $j$, there exists $f_j\to \infty$ (as $n\to \infty$) such that for $n$ sufficiently large,
\begin{equation*}
    \kappa_j \;=\;\ln(p_j-s_j) + f_j
\end{equation*} 
(A7) \quad For each block $j$, there exists $g_j\to \infty$ such that for $n$ sufficiently large,
    \begin{equation*}
        \sqrt{\frac{(1-\gamma) n\rho(\bs X)}{6}}{{\beta_{\min}^*}_j} \,- \,\sqrt{\kappa_j} \;=\; \sqrt{\ln(s_j)} + g_j .
    \end{equation*}
where $\gamma:=\frac12(1+\max_j\frac{\ln(p_j-s_j)}{\kappa_j})\in (\tfrac12,1)$.\\[5pt] 

These assumptions are similar to A4--A5 formulated for the sequence model.
 By Proposition~\ref{prop:isthresh}, in the orthonormal 
 case $\bs X^{\top}\bs X=n\I_p$, setting block penalties $\kappa_j$ 
 is equivalent to hard thresholding with block thresholds $\tau_j=\sqrt{2 \kappa_j/n}$.
 Assumptions A4--A5 can then be translated into assumptions on $\kappa_j$, taking $\rho(\bs X)=1$. In that case, A6 is of the same order as A4, up to a term $f_j$ that can grow at an arbitrarily slow rate. %An oracle choice for $f_j$ is presented in Theorem~\ref{theo:suffcondlinearmodel}. 
 Further, if one sets $f_j$ such that $\ln(p_j-s_j)= O(f_j)$ then $1-\gamma > 0$, and then A7 essentially requires $n \rho(\bs X) \beta_{\min,j}^*$ to be larger (up to constants) than $\sqrt{\kappa_j} + \sqrt{\ln(s_j)}$, analogously to A5 in the sequence model. \color{black}
% A natural choice stemming from a Bayesian perspective discussed in Section~\ref{sec:informed_l0}, is $f_j=0.5\ln(n)$. %for all $j$. 
%Then, if $p= O(n^k)$ for some $k>0$ we have that $1-\gamma>0$. \color{black}
% The betamin assumption~A7 is the same as A5 up to, again, a term $g_j$ increasing at arbitrarily slow rate, and the factor $\sqrt{(1-\gamma)/3}$. \david{In A7 I see $(1-\gamma)/6$ not $(1-\gamma)/3$.} This factor is smaller than $1/2$ and depends on the choice of $f_j$'s in A6. \david{"This factor..." is confusing, better state the factor.} It is bounded away from 0 as long as $f_j= \Omega(\ln(p_j-s_j))$ for all $j$. \david{Can we instead say $\ln(p_j-s_j)=O(f_j)$ in the previous sentence? And in the next sentence, say instead: For example, if $f_j=0.5 \ln(n)$ then $(1-\gamma)>0$ as long as $p=O(n^k)$.}
% This is the case for the earlier natural choice \david{"earlier natural choice is confusing", specify} whenever $n^k=\Omega(p)$ for some $k>0$. \paul{essentially NOT ultra high dimension} \david{Can we just say $p=O(n^k)$?} When $f_j= o(\ln(p_j-s_j))$ for some $j$, A7 requires $\sqrt{\kappa_j}+\sqrt{\ln(s_j)} = o(n\rho(\bs X)\beta_{min,j})$.
%\david{Earlier we said that $f_j$ can grow at an arbitrarily slow rate, but we now see that $f_j$ has to be at least of order $p_j-s_j$. How about saying: This factor depends on the choice of $f_j$ in A6. For example, by default we take $f_j=...$, in which case $\gamma$ is bounded below by ... Perhaps also relate the choice of $(f_j,g_j)$ to the oracle $\kappa_j^\star$'s discussed later in the theorem, like "an ideal choice of $(f_j,g_j)$ leads to oracle penalties $\kappa_j^\star$ given in Theorem 4.2 below" (and, if it helps, briefly explain how the oracle $(f_j,g_j)$ depend on $s_j,p_j,\beta_{\min}^\star$)} \paul{$f_j$ can grow at any rate.  For any growth rate you will discard asymptotically the overfitting models. When $f_j= o(\ln(p_j-s_j))$ for some $j$, the betamin requirement is stronger than $\Theta(\sqrt{\kappa_j} +\; \sqrt{\ln(s_j)} + g_j)$. To see it, the betamin requirement is essentially: $\sqrt{\frac{ n\rho(\bs X)}{2}}{{\beta_{\min}^*}_j}  \;=2(\max_j\frac{\ln(p_j-s_j)}{f_j}+1)(\sqrt{\kappa_j} +\; \sqrt{\ln(s_j)}) + g_j$. It is not very intuitive to have such change in betamin condition i know. It is linked to a technicality of the proof. I have tried hard to get rid of it but could not. I have explained more.}
%\david{I get confused by first seeing $f_j=\Omega(\ln(p_j-s_j))$ and later $f_j=o(\ln(p_j-s_j))$.}
Finally, we remark that in our proof one could take a betamin condition that is slightly less strict than A7, but we present A7 here to facilitate comparison to A5 in the sequence model.  


We can now state our main theorem on the strong variable selection consistency of $\hat{S}^{b}$:
\begin{thm}\label{theo:suffcondlinearmodel}

Under assumptions A1, A2, A6 and, A7, we have
$$\lim_{n \to \infty} \sum_{M \in \M\setminus \{S\}}\E\left(NC(M)\right)= 0 \qquad \mbox{and} \qquad \lim_{n \to \infty} P(\hat{S}^b = S)= 1.$$
\end{thm}


Theorem~\ref{theo:suffcondlinearmodel} considers a linear penalty across blocks
$\sum_{j=1}^b \kappa_j|M_j|$. This linearity assumption can be easily relaxed. In Supplement of Section~\ref{sec:appendixnonlinear} we give a consistency results for nonlinear penalties. Those apply for example
to the EBIC penalty \citep{EBIC} and improve on known requirements on signal strength to have variable selection consistency \cite{ebicdiverging}.

From the proof of Theorem~\ref{theo:suffcondlinearmodel}, we can bound the convergence rate of $P(\hat S^{b}\neq S)$, given in Theorem \ref{prop:fullrecovreg}. We also give oracle block penalties that approximately optimize the bound, and provide the resulting oracle rate of convergence. 
In the statement, $\delta<1$ and $r_b>1$ should be understood as being arbitrarily close to 1. 
\begin{thm}\label{prop:fullrecovreg}
Assume A6--A7.  Then, for all sufficiently large $n$, any $\delta \in (0,1)$, and some $r_b$ such that $\lim_{n \rightarrow \infty} r_b=1$, 
\begin{equation}\label{eq:rateconvreg}
 P(\hat{S}^{b} \neq S)\;\leq\; 6(2^{2b}-2b)r_b\,\sum_{j=1}^b\,\, e^{-\tfrac{\delta}{2}\big[\kappa_j-\ln(p_j-s_j)\big]} + 
    e^{-\tfrac{\delta}{2}\big[(\sqrt{\frac{(1-\gamma) n\rho(\bs X)}{6}}{{\beta_{\min}^*}_j} - \sqrt{\kappa_j})^2 -\ln(s_j) \big]}.
\end{equation}
%where $r_b\to 1$ as $n\to +\infty$.

Moreover, suppose that for all $j=1,\ldots,b$ $\lim_{n\to \infty}\frac{\sqrt{(1-\gamma) n\rho(\bs X)/3}\;{\beta_{\min}^*}_j}{\sqrt{2\ln(p_j-s_j)} + \sqrt{2\ln(s_j)}} >1$ and the $\kappa_j$ are set at the oracle values:
	\begin{equation} \label{eq:oraclepenreg}
\sqrt{\kappa^\star_j\,} \;=\; \frac12\sqrt{\frac{(1-\gamma) n\rho(\bs X)}{6}}{{\beta_{\min}^*}_j}+\frac12\sqrt{\frac{6}{(1-\gamma) n\rho(\bs X)}}\frac{1}{{\beta_{\min}^*}_j}\big(\ln({p_j-s_j})-\ln({s_j})\big).
\end{equation}
Then for all sufficiently large $n$, A6 and A7 hold and for any $\delta \in (0,1)$: 
\begin{equation}\label{eq:oraclerateconvreg}
    P(\hat{S}^{b} \neq S) \; \leq \; 12(2^{2b}-2b)r_b \sum_{j=1}^b e^{-\tfrac{\delta}{2}\big[\frac{(1-\gamma) n\rho(\bs X)}{24}{{\beta_{\min}^*}_j}^2-\max\{\ln(p_j-s_j),\ln(s_j)\}\big]}.
\end{equation}
%where $r_b\to 1$ as $n\to+\infty$.
\end{thm}


\subsection{Necessary assumptions for consistency with block $\ell_0$ penalties}\label{ssec:neccond_linreg}

We derive assumptions on $\kappa_j$ and $\beta_{min,j}$ that are necessary for consistent variable selection with $\hat{S}^b$. Recall that the strength of the block penalties $\kappa_j$ drives the vanishing of the pseudo-posterior probabilities of over-fitted models.
%The block penalties need then to be large enough for the pseudo-posterior probabilities of even the least penalized over-fitted models to vanish. 
For every $j=1,\ldots,b$, denote by $O_j$ the set of models that over-fit by only one variable from block $B_j$:
$$O_j \:=\,\{M\in \M\;|\; M=S\cup\{i\} \text{ where } i\in B_j\setminus S_j\}$$
The asymptotic recovery of $S$ implies that, $\max_{M\in O_j}NC(M)/NC(S)<1$ with probability going to 1 as $n$ grows, for every $j=1,\ldots,b$. For every $M \in O_j$ the ratio $NC(M)/NC(S)$ grows with $L_{MS}$
%$M_k\in O_j$, the ratio $NC(M_k)/NC(S)$ grows with $L_{M_kS}$ 
which, by Lemma~\ref{lemma:s7}, is $\chi^2_1$ distributed. There exists then, for every $M \in O_j$, $Z_M \sim N(0,1)$ such that $L_{MS}=Z_{M}^2$.
%$M_k\in O_j$, $Z_{M_k} \sim N(0,1)$ such that $L_{M_kS}=Z_{M_k}^2$. 
%A necessary assumption on $\kappa_j$ follows from properties of the maximum of correlated Gaussian variables.
Let
%Denote the smallest eigenvalue of their correlation matrix:
\begin{align}
    &\underline{\lambda}_j\;:=\;\lambda_{min}(R^j)\;\;\text{where} \;\;R^j_{k,l}=\operatorname{corr}\big(Z_{M_k},Z_{M_l}\big),\;\;\forall M_k,M_l \in O_j
    \nonumber \\
    &\underline{\lambda}\;:=\;\lambda_{min}(R)\;\;\text{where} \;\;R_{k,l}=\operatorname{corr}\big(Z_{M_k},Z_{M_l}\big),\;\;\forall M_k,M_l \in \cup_{j=1}^b O_j.
    \nonumber
\end{align}
%\paul{note it is not the correlation matrix of the $\sqrt{L_{MS}}=|Z_M|$. This result is based on \cite{hartigan} that gives a lower bound for maximum of correlated Gaussians and not maximum of absolute value of correlated Gaussians. The lower bound is probably the same but proving it would require additional lemmas.}
%We denote $\underline{\lambda}$, the smallest eigenvalue of the correlation matrix of the $Z_{M_k}$, $M_k\in\cup_{j=1}^b O_j$. 
Proposition \ref{prop:nectypeIreg} below describes how $\underline{\lambda}_j$ relates to a necessary condition for consistency.
In Section \ref{ssec:regression_gains} we discuss how $\underline{\lambda}$ describes settings where $\hat S$ is not consistent but $\hat S^b$ is. \color{black}
%A lower bound on the maximum of correlated Gaussians in \cite{hartigan} yields the following result.
\begin{prop}\label{prop:nectypeIreg}
Assume A1-A2.
     If for some $j=1,\ldots,b$,  $\lim_{n\to\infty}\frac{\kappa_j}{\underline{\lambda}_j^2\ln(p_j-s_j)} <1$, then  %under assumptions A1 and A2, 
     $$\lim_{n \to \infty} P\Big(\max_{M\in O_j}\frac{NC(M)}{NC(S)}<1\Big)= 0 \quad\text{and}\quad \lim_{n \to \infty} P(\hat{S}^b=S)= 0.$$
\end{prop}

%\paul{Note in passing: \cite{hartigan} nicely complements the result Piotr found on independence being the "worst case" for the maximum of Gaussians. Interestingly it suggests that we could set lower penalty when inactive variables are highly correlated between each other. I remember very well that someone in Toronto remarked that it was somewhat strange that sufficient assumptions for penalty strength did not depend on the level of correlation between variables. I could not answer well back then!}

Proposition~\ref{prop:nectypeIreg} shows that 
a necessary assumption to recover $S$ asymptotically
is: 
\begin{equation}\label{cond:necIallblockpen}
\lim_{n\to\infty}\frac{\kappa_j}{\underline{\lambda}_j^2\ln(p_j-s_j)}\;\geq\;1\qquad\mbox{for all }j=1,\ldots,b.
\end{equation}
When the inactive variables are orthogonal, for every $j=1,\ldots,b$ and $M\in O_j$, the $Z_M$ are independent and $\underline{\lambda}_j=1$. Then \eqref{cond:necIallblockpen} recovers %Under orthogonal design, \eqref{cond:necIallblockpen} then recovers 
the necessary assumption~\eqref{cond:necIallblockthresh} presented for the sequence model.

Consider now an under-fitted model $M\subset S$. %The recovery of $S$ implies that with probability going to 1 as $n$ grows, $NC(M)/NC(S)<1$. 
By definition the ratio $NC(M)/NC(S)$ %By definition of $NC$ the ratio 
grows with block penalties $\kappa_j$ and shrinks with $L_{SM}$, which is distributed $\chi^2_{|S\setminus M|}(\mu_{SM})$ by Lemma~\ref{lemma:s7}.
For the ratio to be small, $\mu_{SM}$ must grow
%Small values of the ratio are attained only if $\mu_{SM}$ grows 
fast enough compared to $\kappa_j$. Lemma~\ref{lem:upperboundnoncentral} shows that $\mu_{SM}$ is bounded by the largest active signals in $S$ that is not in $M$.
\begin{lem}\label{lem:upperboundnoncentral}
    For any $T\subseteq S$, $$\mu_{ST} \leq n\,\bar{\lambda}\,\sum_{j=1}^b|S_j\setminus T_j|\,\max_{i\in S_j\setminus T_j}{\beta_{i}^*}^2 \quad \text{where}\quad \bar{\lambda}\,:=\,\lambda_{\max}\Big(\frac1n \bs X_{S}^\top\bs X_{S}\Big)$$
\end{lem}

It follows that a necessary condition is that the active $\beta_i^*$ that are missing in any underfitted $M$ are not too small. We next formally define what we mean by small and large signals, and also define a subset of intermediate signals 
%introduce a formal distinction between small and large signals in $S$. We also introduce a subset of signals that are not small or large 
that will be used in Section~\ref{ssec:nobetaminhyp}.  For fixed penalties $\kappa_j$'s and %and every $j=1,\ldots,b$, 
$\gamma$ and $g_j$ as defined in A7, let
\begin{eqnarray}
    &S_j^{S}(\kappa)\;&:=\;\Bigg\{\beta_i^* \in S_j \,\Big|\,\sqrt{n\bar{\lambda}}|\beta_{i}^*|=o\big(\sqrt{\kappa_j}\big) \Bigg\}\nonumber\\
    &S_j^{L}(\kappa)\;&:=\;\Bigg\{\beta_i^* \in S_j \,\Big|\,\sqrt{\frac{(1-\gamma) n\rho(\bs X)}{6}}|\beta_{i}^*| \,- \,\sqrt{\kappa_j} \;=\; \sqrt{\ln(s_j)} + g_j\Bigg\}\label{eq:defslsi}\\
    &S_j^{I}(\kappa)\;&:=\;S_j \setminus \big(S_j^{L}(\kappa) \cup S_j^{S}(\kappa)\big).
    \nonumber 
\end{eqnarray} \\
The subset $S_j^{S}(\kappa)$ gathers signals in $S_j$ that are small with respect to the penalty $\kappa_j$, and $S_j^{L}(\kappa)$ those that are large in that they satisfy Condition A7, and $S_j^{I}(\kappa)$ those that are neither large nor small. 
Proposition \ref{prop:nectypeIIreg} below says that if the set of small signals $S^{S}(\kappa)=\cup_{j=1}^b S_j^{S}(\kappa)$ is not empty, then consistent model recovery is not possible. \color{black}
%Take $M=S\setminus S^{S}(\kappa)$ where $S^{S}(\kappa)=\cup_{j=1}^b S_j^{S}(\kappa)$. If $S^{S}(\kappa)\neq \emptyset$, the missing active signals in $M$ and thus $\mu_{SM}$ are all small compared to the $\kappa_j$'s. \david{Should it be that "some active signals are small"? (rather than all)} 
%We then establish the following result.
%Using the properties of the noncental $\chi^2$ distribution, we then establish the following convergence of the probabilities of $(NC(M)/NC(S)<1)$ and thus of $(\hat{S}^{b}=S)$.

\begin{prop}\label{prop:nectypeIIreg}
    If $S^{S}(\kappa)\neq\emptyset$ and for all $j=1,\ldots,b$, $\kappa_j\to \infty$, then, under assumption A1, 
    $$\lim_{n \to \infty} P\Big(\frac{NC(S\setminus S^{S}(\kappa))}{NC(S)}<1\Big)= 0 \quad\text{and}\quad \lim_{n \to \infty} P(\hat{S}^b=S)= 0.$$
\end{prop}
%Property~\ref{prop:nectypeIIreg} shows that for any increasing penalty, when $S$ contains signals that are too small with respect to that penalty, the smaller model $S\setminus S^{S}(\kappa)$ is favored at the expense of $S$ and recovery fails. 
It follows that a necessary assumption for variable selection consistency is
\begin{equation}\label{cond:necIIallblockpen}
\lim_{n\to\infty}\frac{\sqrt{n\bar{\lambda}}\beta_{\min,j}^*}{\sqrt{\kappa_j}}\;>\;0\qquad\mbox{for all }j=1,\ldots,b.
\end{equation}
When active variables are orthonormal, %orthogonal and normalized, 
that is $\bs X_S^{\top}\bs X_S = nI$, then $\bar{\lambda}=1$. 
Assumption~\eqref{cond:necIIallblockpen} is then milder and less tight than the corresponding assumption~\eqref{cond:necIallblockthresh} for the sequence model. 

An immediate consequence of Proposition~\ref{prop:nectypeIreg} and Proposition~\ref{prop:nectypeIIreg} is the next result.
\begin{cor}\label{cor:necrecovreg}
    If for some $j\in\{1,\ldots,b\}$   $\lim_{n\to\infty}\frac{\sqrt{n\bar{\lambda}}\beta_{\min,j}^*}{\underline{\lambda}_j\sqrt{\ln(p_j-s_j)}} = 0$, then, under assumption A1 and A2,  $\lim_{n \to \infty} P(\hat{S}^b=S) < 1$.
\end{cor}
Corollary~\ref{cor:necrecovreg} yields a necessary assumption for consistency:
\begin{equation}\label{eq:betaminnecreg}
\lim_{n\to\infty}\frac{\sqrt{n\bar{\lambda}}\beta_{\min,j}^*}{\underline{\lambda}_j\sqrt{\ln(p_j-s_j)}} > 0\qquad\mbox{for all }j=1,\ldots,b.
\end{equation}
%In an orthonormal design where $\bs X^{\top} \bs X=nI$, %for all $j=1,\ldots,b$, 
%then $\underline{\lambda}_j=1$ and $\bar{\lambda}=1$. 
In the particular case of orthogonality and diverging $s$, the necessary assumption \eqref{eq:betamin} given for the sequence model is then stricter and tighter than \eqref{eq:betaminnecreg}. Assumption \eqref{eq:betaminnecreg} is however a more general necessary betamin condition that applies in all correlated, orthogonal, low and high dimensional settings, for $s$ fixed or diverging. %\paul{I am happy about this result because I believe there was no such tight necessary assumption for variable selection consistency in regression in the literature. The only ones I found were for LASSO (far less efficient/tight for selection in correlated settings) or compressed sensing}.

\subsection{Benefits of block penalties}
\label{ssec:regression_gains}

We discuss separately the benefits in terms of sufficient conditions for variable selection consistency and those in terms of convergence rate. The results are analogous to those in Section~\ref{sec:properties_orth} for the sequence model. %main message is that benefits discussed in the sequence model in Section~\ref{sec:properties_orth} extend to linear regression. 

\bigskip

\noindent Assumptions A6--A7 give ranges of penalties that are sufficient (but not necessary) for asymptotic support recovery. For simplicity, we restrict our discussion to choices of penalty $\kappa_j$ such that, in A6, $f_j=\Omega(\ln(p_j-s_j))$ and $\gamma$ is bounded in $(0,1/2)$. For the standard threshold selector $\hat{S}$, the single penalty $\kappa$ is essentially required to satisfy, %up to constants,
\begin{equation}\label{eq:rangepenex}
  \sqrt{\ln(p-s)}\;\;<\;\;\sqrt{\kappa}\;\;<\;\;\sqrt{n\rho(\bs X)}\beta^*_{\min}+\sqrt{\ln(s)}.	
\end{equation}
For a block threshold selector $\hat{S}^{b}$, assuming without loss of generality that the global minimum $\beta_{\min}^*$ is in block $b$, the ranges for the $\kappa_j$'s essentially are, %up to constant and bounded factors:
    \begin{eqnarray}\label{eq:rangeblockpenex}
      \sqrt{\ln(p_j-s_j)}\;\;&<\;\;\sqrt{\kappa_j}\;\;<&\;\;\sqrt{n\rho(\bs X)}\beta^*_{\min,j}-\sqrt{\ln(s_j)}
      %
%  \sqrt{\ln(p_1-s_1)}\;\;&<\;\;\sqrt{\kappa_1}\;\;<&\;\;\sqrt{n\rho(\bs X)}\beta^*_{\min,1}-\sqrt{\ln(s_1)}\nonumber\\
%  &\;\;\vdots \\
%  \sqrt{\ln(p_b-s_b)}\;\;&<\;\;\sqrt{\kappa_b}\;\;<&\;\;\sqrt{n\rho(\bs X)}\beta^*_{\min}-\sqrt{\ln(s_b)} \nonumber.       
    \end{eqnarray}
As for the sequence model, by definition, the ranges for $\hat S^b$ are wider than that for $\hat S$. Sufficient conditions to have variable selection consistency are then milder with block penalties. The results of Section~\ref{ssec:neccond_linreg} allow to precisely describe settings in which $\hat S^b$ is consistent and $\hat S$ is not.
\begin{cor}\label{cor:selconstimpossiblereg}
Assume A1, A2, A5, A6. If
\begin{eqnarray*}
&\lim_{n\to\infty}\frac{\sqrt{n\bar{\lambda}}\beta_{\min}^*}{\underline{\lambda}\sqrt{\ln(p-s)}} = 0
\end{eqnarray*}
then 
$P(\hat{S} = S)\not \to 1$ and $P(\hat{S}^{b} = S)\to 1$.
\end{cor}

Observe that \eqref{eq:rangepenex} and \eqref{eq:rangeblockpenex} are analogous to \eqref{eq:rangethreshex} and \eqref{eq:rangeblockthreshex} for the sequence model, up to rescaling by $\sqrt{2/n}$ and factor $\rho(\bs X)$. The gains in terms of valid thresholds for consistency discussed in the examples of Section \ref{sec:examples} remain applicable %are essentially the same for 
to linear regression. %The main difference lies in that $\rho(\bs X)$ may be smaller than 1 making the range for $\hat S$ narrower and $\hat S^b$ even more beneficial. 

Let $\underline{\beta}^{*,b}_{\min,reg}$ and $\underline{\beta}^{*}_{\min,reg}$ be the smallest signal recoverable by $\hat{S}^b$ and $\hat{S}$ respectively. \color{black}
%Consider the smallest signal recoverable by $\hat{S}^b$ and $\hat{S}$ while still being variable selection consistent that we denote $\underline{\beta}^{*,b}_{\min,reg}$ and $\underline{\beta}^{*}_{\min,reg}$. 
Assuming $\beta_{min}$ is in block $b$, A6-A7 require that they essentially satisfy \color{black}
%they are essentially required by A6 and A7 to satisfy, up to bounded factors,
\begin{eqnarray*}
    \underline{\beta}^{*,b}_{\min,reg}&>& \sqrt{\frac{\ln(p_b-s_b)}{ n\rho(\bs X)}} + \sqrt{\frac{\ln(s_b)}{ n\rho(\bs X)}} \\
    \underline{\beta}^{*}_{\min,reg}&>& \sqrt{\frac{\ln(p-s)}{ n\rho(\bs X)}} + \sqrt{\frac{\ln(s)}{ n\rho(\bs X)}} 
\end{eqnarray*}
These lower bounds are the same as \eqref{eq:betaminbardef} for the sequence model, up to rescaling by $\sqrt{2}$ and a factor $\rho(\bs X)$ that do not affect the ratio $\underline{\beta}^{*,b}_{\min,reg}/\underline{\beta}^{*}_{\min,reg}$. The discussion and examples of the benefits in the smallest recoverable signals in Section~\ref{sec:examples} hence extend to linear regression.

Let ${OR}^b_{reg}$ be the oracle convergence rate for $\hat{S}^{b}$ in Theorem~\ref{prop:fullrecovreg}, and ${OR}_{reg}$ that for $\hat{S}$. Then we have:
\begin{equation*}
    \frac{{OR}^b_{reg}}{{OR}_{reg}}=(2^{2b-1}-b)\frac{r_b}{r_1}\sum_{j=1}^be^{-\frac{\delta}{2}\left[\frac{ n\rho(\bs X)}{24}\left((1-\gamma){\beta^\star_{\min,j}}^2-(1-\gamma'){\beta_{\min}^\star}^2\right)+\max\{\ln(p-s),\ln(s)\}-\max\{\ln(p_j-s_j),\ln(s_j)\}\right]}
\end{equation*}
where $\gamma=\frac12(1+\max_j\frac{\ln(p_j-s_j)}{\kappa_j^*})$ and $\gamma'=\frac12(1+\frac{\ln(p-s)}{\kappa^*})$. The ratio above is essentially the same as \eqref{eq:ratioCR} for the sequence model, up to certain factors.  Specifically, $r_b/r_1$, $\delta$, $(1-\gamma)$ and $(1-\gamma')$ are close to 1, whereas $ \rho(\bs X)/24$ slows the convergence rate gains but does not alter the essence of the ratio. % as long as $\rho(\bs X)$ remains bounded above 0. 
The factor $2^{2b-1}$ highlights however a potential limitation: guarantees of gains with block $\ell_0$ penalties deteriorate when one considers a large number of blocks $b$.
We remark that such deterioration may be a consequence of our proof strategy, rather than an inherent limitation of block penalties. Studying cases with %large $b$, including 
$b \to \infty$ is left as future research.

\subsection{Convergence with no betamin condition}\label{ssec:nobetaminhyp}

We now derive a convergence result for pseudo-posterior probabilities $NC(M)$, $M\in\mathcal{M}$ under no assumption on the minimal signal strength. The result generalizes Theorem~\ref{theo:suffcondlinearmodel} and is key to the proofs of Section~\ref{sec:informed_l0}.

%In \eqref{eq:defslsi}, we defined, for a choice of penalties $\kappa_j$, the set of large and intermediate signals. 
Let $\T$ be the set of models that contain all large signals $S^L(\kappa)=\cup_{j=1}^b S^{L}_j(\kappa)$ defined in \eqref{eq:defslsi}, some or none of the intermediate signals in $S^{I}(\kappa)=\cup_{j=1}^b S_j^{I}(\kappa)$, and no truly inactive parameters. \color{black}
%\begin{equation}
%    \T \;:=\;\Big\{M \in \M \,\big|\, M\,=\, S^L(\kappa)\,\cup\,R,\; R\in \mathcal{P}\big(S^{I}(\kappa)\big)  \Big\}
%\end{equation}
%where for any set $A$, $ \mathcal{P}(A)$ denotes the power set of $A$.
Theorem \ref{thm:convtoT} below shows,
%Following a strategy similar to that of the proof of Theorem~\ref{theo:suffcondlinearmodel}, we demonstrate, 
assuming only sufficiently large penalties (A6), 
that posterior model pseudo-probabilities concentrate on $\T$. \color{black}
%the vanishing of the sum of pseudo-probabilities of all models but those in $\T$. 
%The main difference with respect to the proof of Theorem~\ref{theo:suffcondlinearmodel} is that, for any $M\not\in \T$, we bound $\E(NC(M))$ by a function of $C(M)-C(T)$ for a well chosen model $T$ in $\T$. 
\begin{thm}\label{thm:convtoT}
    Assume A1, A2, A6 and $|S^{S}_j(\kappa)|=O(p_j-s_j)$ for every $j=1,\ldots,b$:
    \begin{longlist}
        \item $\lim_{n \to \infty} \sum_{M \in \M\setminus \T} \E\left(NC(M)\right)= 0$,
        \item If $|S^{I}(\kappa)|=O(1)$, then $\lim_{n \to \infty} P(\hat{S}^b \in \T)= 1$.
    \end{longlist}
\end{thm}
%Theorem~\ref{thm:convtoT} shows that, asymptotically, with only assuming
That is, if one sets
sufficiently large penalties, all inactive and all small active signals are discarded. On the contrary, all large signals (relative to the specified penalties) are retained. Intermediate signals in $S^{I}(\kappa)$ may or may not be retained. Assumption $|S^{I}(\kappa)|=O(1)$ in $(ii)$ was made for simplicity, in fact $|S^{I}(\kappa)|$ can be allowed to increase moderately. 

\section{Data analysis with block $\ell_0$ penalization}\label{sec:informed_l0}

\subsection{A Bayesian-inspired estimator of sparsity}

In Sections~\ref{sec:properties_orth} and~\ref{sec:properties_legression}, we derived properties for block penalties where one sets the penalties to oracle values.
We now
propose two data analysis methods anchored in a Bayesian perspective
that do not require oracle values, and adapt to the unknown sparsity in the data-generating truth.
A Bayesian approach provides a natural estimator of the number of truly active variables in each block and thus, by using empirical Bayes, a straightforward way to adapt penalties to sparsity in each block.
\\
\\
Consider a Bayesian framework for \eqref{eq:linearmodel} where
the model $M$ is a $p$-dimensional vector of variable inclusion indicators $m_i=I(\beta_i \neq 0)$. Consider a joint prior on parameters and models 
\begin{equation}\label{eq:bayesianframework}
    p(\boldsymbol{\beta},M \mid \boldsymbol{\theta})=p( \boldsymbol{\beta} \mid  M) p(M \mid \boldsymbol{\theta})
\end{equation}
where $p(M \mid \boldsymbol{\theta})$ is the prior probability of model $M$, which depends on hyperparameters $\bs \theta$ described below,
and $p( \boldsymbol{\beta} \mid  M)$ is a prior on regression coefficients given the model. 

Posterior model probabilities are $p(M \mid \bs y, \bs \theta) \propto p(\bs y \mid M) p(M \mid \bs \theta)$, where $p(\bs y \mid M)$ is the so-called marginal likelihood of model $M$. The BIC approximation \citep{BIC} gives
\begin{eqnarray*}
 p( \bs y  \mid M)\approx p(\bs y \mid \tilde{\bs \beta}^{(M)})n^{-|M|/2}
\end{eqnarray*}
where $\tilde{\bs \beta}^{(M)}$ is the MLE under model $M$,
for a wide family of priors $p(\bs \beta \mid M)$, and hence
%It follows that the posterior model probabilities satisfy
\begin{eqnarray*}
\ln p(M \mid\bs y) \approx \ln p(\bs y \mid \tilde{\bs \beta}^{(M)}) - \frac{|M|}{2}\ln(n) + \ln p(M\mid\bs \theta) + \text{constant}.
\end{eqnarray*}
Thus, any block or standard $\ell_0$ penalty selector discussed in earlier sections can be written as $\hat{M}= \arg\max_M p(M \mid \bs y)$,
for some $p(M \mid \bs \theta)$. 
In particular, 
if variable inclusions are independent a priori,
with constant inclusion probability $\theta^{(j)}$ within each block $j$, then
\begin{eqnarray} \label{eq:block_modelprior}
p( M \mid \boldsymbol{\theta})=
        \prod_{i=1}^p \operatorname{Bern}\left(m_i ; \theta_i\right) I(M \in \M) \mbox{, where } \forall \, i\in B_j, \: \theta_i=\theta^{(j)}
\end{eqnarray}
where $\bs \theta=(\theta^{(1)},\ldots,\theta^{(b)})$.
Then $\ln p(M \mid \bs \theta)$ defines the block penalties studied in Section \ref{sec:properties_legression} with
\begin{equation}\label{eq:blockpenpriorinclu}
    \kappa_j=\tfrac{1}{2}\ln(n)+\ln\big({\theta^{(j)}}^{-1}-1\big),
\end{equation}
and $P(M \mid \bs y) \approx NC(M)$, the normalized score in Section \ref{sec:properties_legression}.
This connection motivates estimating the number of truly active variables in block $j$ by %for any block $\ell_0$ procedure:
\begin{equation}\label{eq:shat}
    \widehat{s}_j \,:=\, \sum_{i\in B_j}\,\,\sum_{M\in \M | i \in M}NC(M).
\end{equation}
%The normalized score $NC(M)$ can be seen as the posterior probability for model $M$ in Bayesian framework \eqref{eq:bayesianframework}. 
Replacing $NC(M)$ by $P(M \mid \bs y)$ in \eqref{eq:shat} gives
\begin{equation*}
    \sum_{i\in B_j}\,\,\sum_{M\in \M | i \in M}P(M\mid \bs y) = \sum_{i\in B_j}P(\beta_i \neq 0 \mid \bs y) = E\big(s_j \mid \bs y\big).
\end{equation*}
From a Bayesian perspective, $\widehat{s}_j$ is then %the sum of posterior inclusion probabilities and 
the posterior mean of $s_j$. The estimator of the proportion of truly active variables $\hat{\theta}^{(j)}= \widehat{s}_j/p_j$ can also be justified as an empirical Bayes estimator %procedure. The empirical Bayes estimator of $\bs \theta$ (cf \eqref{eq:block_modelprior}) 
maximizing the marginal likelihood of $\bs y$ given $\bs \theta$, i.e.
\begin{eqnarray}
\hat{\bs \theta}
%\hat{\bs \theta}^{EB}
\,:=\, \arg\max_{\bs \theta} p(\bs y \mid \bs \theta)=
\arg\max_{\bs \theta} \int p(\bs y \mid \bs \beta, M) d P(\bs \beta, M \mid \bs \theta).
\nonumber
\end{eqnarray}
% In the Bayesian framework \eqref{eq:block_modelprior}, $\widehat{s}_j/p_j$ corresponds to the empirical Bayes estimate of $\bs \theta$.

\begin{lem}\label{lem:empiricalbayesinterpret} For $p(M \mid \bs \theta)$ in \eqref{eq:block_modelprior}, the empirical Bayes estimator satisfies 
$$\hat{\theta}^{(j)} = \frac{1}{p_j}\sum_{i\in B_j}P(\beta_i \neq 0 \mid \bs y).$$
\end{lem}
The estimator $\widehat{s}_j/p_j$ is hence well-founded from a Bayesian perspective and, as shown below, also has attractive frequentist properties under mild assumptions. In Section~\ref{sec:properties_legression} we defined $S^{S}_j(\kappa)$ to be the set of small signals in block $j$ and  $S^{L}_j(\kappa)$ the set of larger signals (\emph{cf} \eqref{eq:defslsi}). In Theorem~\ref{thm:convtoT} we also derived a convergence result on pseudo posterior probabilities that yields the following asymptotic bounds on the frequentist expectation of $\widehat{s}_j/p_j$.
\begin{prop}
    \label{prop:shatconsistence}
Assume A1, A2, A6 and $|S^{S}_j(\kappa)|=O(p_j-s_j)$ for every $j=1,\ldots,b$. Then
$$ \,\, \frac{|S^{L}_j(\kappa)|}{p_j}\leq \lim_{n\to\infty}\E\bigg(\frac{\widehat{s}_j}{p_j}\bigg)\leq \frac{s_j-|S^{S}_j(\kappa)|}{p_j}\qquad\text{for all }j=1,\ldots,b$$
\end{prop}
An immediate consequence is that, if the betamin condition A7 also holds, then $\widehat{s}_j/p_j \stackrel{L_1}{\longrightarrow} s_j/p_j$
because $S_j=S_j^{L}(\kappa)$ and $S_j^{S}(\kappa)=\emptyset$ for all $j$. The proposition also shows that, when A7 is not met, $\widehat{s}_j/p_j$ is asymptotically downward biased by at least $|S^{S}_j(\kappa)|/p_j$, but it is guaranteed to be larger than the proportion of signals satisfying A7. 

\subsection{Data-based block selection method}

The oracle block penalties of Section~\ref{sec:properties_legression} have varying strength depending on the unknown $\beta^*_{\min,j}$ and number of active signals in each block $s_j$, and we just saw that
%While we do not have a reliable estimator of $\beta^*_{\min,j}$, we have just seen 
the latter can be reliably estimated. We propose a method that uses $\widehat{s}_j$ and the connection in \eqref{eq:blockpenpriorinclu} to adapt the $\ell_0$ penalty in each block. It has two variants. The first one directly uses the empirical Bayes approach of Lemma~\ref{lem:empiricalbayesinterpret}. % to set block penalties. 
The second variant is motivated by Theorem~\ref{theo:suffcondlinearmodel}. Considering jointly \eqref{eq:blockpenpriorinclu} and assumption A6, a natural choice is setting $\theta^{(j)}=(p_j-s_j+1)^{-1}$. This choice yields block penalties that are sufficiently large for Theorem~\ref{theo:suffcondlinearmodel} to hold, where $f_j$ in A6 takes value $\frac{1}{2}\ln(n)$ for all $j$. The second variant approximates this choice.\\

\textbf{Algorithm 1}
\begin{enumerate}
    \item Set $\kappa_j=\kappa^{\circ}=\ln(p)+\frac{1}{2}\ln(n)$ 
    for $j=1,\ldots,b$. Compute $\widehat{s}_j/p_j$ in \eqref{eq:shat} for $j=1,\ldots,b$.

    \item Obtain $\hat{S}^{EB,b}$ solving \eqref{eq:Mhat} with $\kappa_j^{EB}= \ln(p_j/\widehat{s}_j -1)+ \tfrac{1}{2}\ln(n)$.
    Alternatively, obtain $\hat{S}^{A,b}$ solving \eqref{eq:Mhat} with
    $\kappa_j^{A} = \ln(p_j-\widehat{s}_j)+ \tfrac{1}{2}\ln(n)$.
\end{enumerate}

We refer to $\hat{S}^{EB,b}$ as the block empirical Bayes selector, and to $\hat{S}^{A,b}$ as the block adaptive selector.
Step 1 approximates posterior model probabilities under equal prior inclusion probabilities $\theta^{(j)}=(p+1)^{-1}$ in \eqref{eq:block_modelprior} across blocks, and estimates the proportion of truly active coefficients with $\hat{s}_j/p_j$. Step 2 selects a model using the block penalties $\kappa_j^{EB}$ induced by $\theta^{(j)}=\hat{s}_j/p_j$,
or alternatively setting $\kappa_j^{A}= \kappa_j^{EB} + \ln(\hat{s}_j)$.
Under mild assumptions both $\kappa_j^{EB}$ and $\kappa_j^{A}$ lead to variable selection consistency.
The main difference is that in non-sparse settings where $s$ grows faster than $\sqrt{n}$, the penalty $\kappa_j^{EB}$ can be insufficient and there $\kappa_j^{A}$ might be preferred.

Theorem \ref{theo:empbayesselconsist} shows
the consistency of $\hat{S}^{EB,b}$ and $\hat{S}^{A,b}$, under
assumptions:\\[5pt]
(A8)\quad $s=o(\sqrt{n})$\\[5pt]
(A9)\quad For each block $j$, there exists $a_j\to \infty$ such that for $n$ sufficiently large, 
\begin{equation*}\label{cond:suffIIreg2}
        \sqrt{\frac{ (1-\psi)n\rho(\bs X)}{6}}{{\beta_{\min}^*}_j} \,- \,\sqrt{\ln\bigg(\frac{p_j}{|S^{L}_j(\kappa^{\circ})|}-1\bigg)+\frac12\ln(n)} \;=\; \sqrt{\ln(s_j)} + a_j .
    \end{equation*}
where $\psi=\frac12\big(1+\max_j \frac{\ln(p_j-s_j)}{\ln(p_j/s_j-1)+0.5\ln(n)}\big)$.
%Using Theorem~\ref{thm:convtoT} and Proposition~\ref{prop:shatconsistence}, we show that under A8 and A9, the $\kappa^{EB}_j$'s satisfy A6 and A7 with probability going to 1 and conclude on the selection consistency  of $\hat{S}^{EB,b}$.

\begin{thm}\label{theo:empbayesselconsist}
    Assume A1, A2, A8, A9, and $|S^{S}_j(\kappa)|=O(p_j-s_j)$ for every $j=1,\ldots,b$. Then
$\lim_{n \to \infty} P(\hat{S}^{EB,b} = S)= 1$.
\end{thm}

Theorem \ref{theo:informedl0consist} shows
the consistency of $\hat{S}^{A,b}$. It no longer requires A8, and A9 is replaced by the slightly more stringent A10. \\[5pt]
(A10)\quad For each block $j$, there exists $c_j\to +\infty$ such that for $n$ sufficiently large,
\begin{equation*}
        \sqrt{\frac{ (1-\xi)n\rho(\bs X)}{2}}{\beta_{\min}^*}_j \,- \,\sqrt{\ln\big(p_j-|S^{L}_j(\kappa^{\circ})|\big)+\frac12\ln(n)} \;=\; \sqrt{\ln(s_j)} + c_j .
    \end{equation*}
where  $\xi=\frac12(1+\max_j \frac{\ln(p_j-s_j)}{\ln(p_j-s_j)+0.5\ln(n)})$.

\begin{thm}\label{theo:informedl0consist}
    Assume A1, A2, A10, and $|S^{S}_j(\kappa)|=O(p_j-s_j)$ for every $j=1,\ldots,b$. Then 
$\lim_{n \to \infty} P(\hat{S}^{A,b} = S) = 1$.
\end{thm}

\subsection{Benefits of data-based block selection}

By Theorem~\ref{theo:empbayesselconsist} a standard (non block-based) empirical Bayes $\hat S^{EB}$ selector that sets a single step 2 penalty $\kappa^{EB}=\ln(p/\widehat{s} -1)+ \tfrac{1}{2}\ln(n)$ is variable selection consistent under the betamin assumption:
\begin{equation}\label{eq:betaminEB}
        \sqrt{\frac{ (1-\psi')n\rho(\bs X)}{6}}{\beta_{\min}^*} \,- \,\sqrt{\ln\bigg(\frac{p}{|S^{L}(\kappa^{\circ})|}-1\bigg)+\frac12\ln(n)} \;=\; \sqrt{\ln(s)} + a'_j .
    \end{equation}
where $\psi'=\frac12\big(1+ \frac{\ln(p-s)}{\ln(p/s-1)+0.5\ln(n)}\big)$ and $a'_j \to \infty$. We have:
\begin{eqnarray*}
&\sqrt{\ln\bigg(\frac{p}{|S^{L}(\kappa^{\circ})|}-1\bigg)+\frac12\ln(n)} + \sqrt{\ln(s)} \geq \sqrt{\ln\bigg(\frac{p_j}{|S^{L}_j(\kappa^{\circ})|}-1\bigg)+\frac12\ln(n)} + \sqrt{\ln(s_j)}
\end{eqnarray*}
and also $\psi > \psi'$. Sufficient conditions for consistency and the smallest signal recoverable are then milder for $\hat S^{EB,b}$ than for $\hat S^{EB}$. %The smaller $p_j/|S^{L}_j(\kappa^{\circ})|$ is compared to $p/|S^{L}(\kappa^{\circ})|$, the milder is A9  compared to \eqref{eq:betaminEB}. % in block $B_j$. 
%$\hat S^{EB,b}$ is then most beneficial when the external information identifies blocks with low proportion of inactive or small signals.

Similarly, a standard (non block-based) adaptive selector $\hat S^{A}$ setting a common penalty $\kappa^{A}=\ln(p-\widehat{s})+ \tfrac{1}{2}\ln(n)$ in Step 2 is consistent under stricter assumptions than $\hat S^{A,b}$. The smallest signal recoverable is also smaller with $\hat S^{A,b}$. %In that case, the gains are driven by the ability of external information to identify blocks with low \emph{numbers} of inactive or small signals.


%Our lack of convergence rates for $\hat S^{EB,b}$ and $\hat S^{A,b}$ precludes a detailed discussion of their benefits in this aspect. However, we remark that the $\kappa^{EB}_j$ and $\kappa^{A}_j$ are constantly lower than the single $ \kappa^{EB}$ and $\kappa^{A}$ respectively. This is in contrast to the oracle block penalties $\kappa^*_j$ that are larger in blocks with large signals. As a result, $\hat S^{EB,b}$ and $\hat S^{A,b}$ have higher probability of false positives and lower probability of false negatives than $\hat S^{EB}$ and $\hat S^{A}$ respectively.  $\hat S^{EB,b}$ and $\hat S^{A,b}$ underperform in settings with large signals and large number of inactive variables and outperform in settings with small signals. We show next with empirical results that the loss in probability of support recovery with $\hat S^{EB,b}$ and $\hat S^{A,b}$ is minor while the benefit potentially large.  

%\paul{Theorem on rates in Section~\ref{sec:properties_legression} hold when A6 and A7 hold. We can only prove A6 and A7 hold with probability going to 1. To have a rate for the EB we would need to have a rate at which A6 and A7 hold}

\section{Numerical illustrations}\label{sec:simulation}

We illustrate the performance of Algorithm 1 on simulated data. We run our proposed method in linear regression under the asymptotic regimes and block assumptions of Examples 1, 3 and 4 in Section \ref{sec:examples}. We also consider an additional setting, Example 5, that highlights differences between $\hat S^{EB,b}$ and $\hat S^{A,b}$. In that example, we set $p=n/2$, $s=3\ln(n)$, and two blocks such that $p_1-s_1=(n-\sqrt{n})/2$, $p_2-s_2=\sqrt{n}/2$ and $s_1 = s_2=3\ln(n)/2$. In step 1 of Algorithm 1 we take $\kappa^{\circ}=\frac12\ln(n)+\ln(p)$. 
To search over models, we rely on \color{black}
%We do not perform an exhaustive enumeration and scoring all of the models. We rely instead on 
the connection between $\ell_0$ penalties and Bayesian variable selection and %search over the model space through 
use the MCMC algorithm in function \texttt{bestIC} in R package \texttt{mombf}. Each visited model is 
%The models explored are 
scored with the BIC approximation \eqref{eq:blockpenpriorinclu}. In Step 2 of Algorithm 1, we obtain $\hat S^{EB,b}$, $\hat S^{A,b}$, $\hat S^{EB}$, $\hat S^{A}$ by scoring all the models visited by the MCMC in Step 1.

We simulate data with $n \in \{20,700\}$ and Gaussian covariates with unit variance and all pairwise correlations equal to 0.5. \color{black}
%under Gaussian designs  with $n$ ranging from 20 to 700 and population covariance matrix:
%$$\Sigma=\begin{pmatrix}
%1\,     &\, 0.5\,  &\ldots & \,0.5\\
%0.5\,   & \,1\,    &\ldots & \,0.5\\
%\vdots  &          &       & \vdots\\
%0.5\,   &\ldots    &\,1 \, & \,0.5 \\ 
%0.5\,   &\ldots    &\,0.5\,& \,1 
%\end{pmatrix}$$
Table~\ref{tab:betaminsim} summarizes the $(\beta_{\min,1}^*,\beta_{\min,2}^*)$ %betamin assumptions 
used in our simulations. %The $\beta_{\min}^*=\beta_{\min,2}^*$ are chosen so the probability of recovery with block and standard selectors approximately spans the interval $(0,1)$ for the range of $n$ considered. 
In Example 4, we set $\beta_{\min,1}^*=\beta_{\min,2}^*$ to represent a setting with non-discriminative blocks. %We use MCMC search and BIC approximation too.

\begin{table*}[ht]
%\caption{Betamin assumptions in simulations}
\label{tab:betaminsim}
\begin{tabular}{@{}cccc@{}}
\hline \\[-1em]
Example & $\beta_{\min,1}^*$ & $\beta_{\min,2}^*$  \\[3pt]
\hline\\[-1em]
$1$    & 0.8 & 0.33  \\[3pt]
$3$    & 0.8 & 0.33  \\[3pt]
$4$    & 0.33 & 0.33 \\[3pt]
$5$    & 0.8 & 0.2   \\[3pt]
\hline
\end{tabular}
\end{table*}

\begin{figure}[ht]
\begin{minipage}{1\columnwidth}
\begin{tabular}{M{63mm}M{63mm}}
        Example 1 - Lin. reg & Example 3 - Lin. reg\\
         \includegraphics[scale=1]{images/p.probrecov.ex1.slogn.pdf} & \includegraphics[scale=1]{images/p.probrecov.ex3.slogn.pdf} \\ 
         Example 4 - Lin. reg& Example 5 - Lin. reg \\
         \includegraphics[scale=1]{images/p.probrecov.ex4.slogn.pdf}  & 
         \includegraphics[scale=1]{images/p.probrecov.ex5.slogn.pdf} 
\end{tabular}
\end{minipage}
\caption{Probability of correct selection with $\hat S^{EB,b}$, $\hat S^{A,b}$, $\hat S^{EB}$, $\hat S^{A}$, and the EBIC penalty in Example 1 (top left), 2 (top right), 4 (bottom left) and 5 (bottom right)}
\label{fig:linregslogn}
\end{figure}

Figure~\ref{fig:linregslogn} plots the empirical probabilities of correct recovery for $\hat S^{EB,b}$, $\hat S^{A,b}$, $\hat S^{EB}$, $\hat S^{A}$ and the EBIC penalty, for Examples 1, 2 and 4. The probabilities are computed over 100 simulations for each $n$.
In Examples 1 and 3 where blocks are discriminative,
%In the examples with discriminative blocks, 
$\hat S^{EB,b}$ and $\hat S^{A,b}$ outperform $\hat S^{EB}$, $\hat S^{A}$ and the EBIC, particularly for small $n$. % in the harder conditions when $n$ is small. 
%When $n$ becomes large, $\hat S^{EB,b}$ and $\hat S^{A,b}$ lose their edge over $\hat S^{EB}$ and $\hat S^{A}$ respectively. This is expected because the penalties in $\hat S^{EB,b}$ and $\hat S^{A,b}$ are smaller than in $\hat S^{EB}$ and $\hat S^{A}$. When $n$ is large, false negatives are less of an issue and there is no more benefit in smaller penalties. 
%Importantly, the underperformance of $\hat S^{EB,b}$ and $\hat S^{A,b}$ is null or minor: the probabilities of the block selectors and their standard counterparts converge to same levels in our examples. $\hat S^{EB,b}$ and $\hat S^{A,b}$ then control the probability of false positives as well as $\hat S^{EB}$ and $\hat S^{A}$ with smaller penalties: they achieve a better balancing of the probabilities of false negatives and false positives. 
In the nondiscriminative block setting, Example 4, $\hat S^{EB,b}$ and $\hat S^{A,b}$ perform very similarly to  $\hat S^{EB}$ and $\hat S^{A}$ respectively.

In Examples 1, 3 and 4, the adaptive block selector $\hat S^{A,b}$ outperforms the empirical Bayes selector $\hat S^{EB,b}$ for large $n$. This occurs because $\hat{S}^{A,b}$ uses larger penalties
%is explained by the larger strength of $\hat S^{A,b}$ penalties 
that control better false positives, whereas false negatives are essentially not an issue anymore for such large $n$. 

In contrast, Example 5 is a setting where smaller penalties are advantageous, because the number of of inactive variables and $\beta_{\min}^*$ are small. In that case,  $\hat S^{EB,b}$ outperforms $\hat S^{A,b}$.

\section{Discussion}\label{sec:discussion}

We studied %provided an understanding of 
how external information can ease variable selection in the sequence model and high-dimensional linear regression. We studied the case where 
variables are partitioned %external information helps partition variables 
into blocks %with potentially distinct characteristics 
and introduced corresponding block-based %externally-informed 
$\ell_0$ selectors. % that let the penalty depend on those blocks. 
We showed that an oracle externally-informed selector converges faster and has less strict betamin condition than the standard $\ell_0$ oracle. We also provided concrete data analysis methods that incorporate external information to improve variable selection properties without requiring oracle knowledge.

A question %that immediately arise 
for future work is how much the assumption of fixed number of blocks $b$ can be relaxed. Our current proof strategies %consistency results 
are robust to moderate increases in the number of blocks but do not work when $b=p$ for example. Another interesting research direction is understanding the benefit of external information for parameter estimation and prediction error.
For example, it is possible to obtain estimation error bounds for the sequence model, but the results depend on the chosen estimator (e.g. $\ell_0$, $\ell_1$ or $\ell_2$) and ensuring their tightness requires separate work elsewhere. By focusing on model selection, we obtained results that apply to essentially all penalties / Bayesian methods in the sequence model.
%\paul{Adapting a proof in Rigollet's high-dim statistics notes I could easily give here a bound on estimation error in the Gaussian sequence model. I don't how tight it is though. For linear regresison, there is proof for BIC in Rigollet's notes and I need to check more closely if it can be adapted. Do you know David a proof for the estimation error with $\ell_0$ selection in linear regression? Would it be interesting/relevant to have those bounds here?} \paul{We could also mention externally informed $\ell_1$ in correlated deigns as future direction.}
%\david{Castillo, Schmidt-Hieber and van der Vaart have work on estimation error for the Sequence Model and linear regression, using posterior model probabilities that are analogs to our pseudo-posterior probabilities. I suppose that one could adapt their proof strategy. Miquel Torrens also studied MSE's in the Sequence Model, his proof strategy was different and more along the lines of what you did here. This said, I'd probably leave this for future work. It might be an easy-ish paper, perhaps.} The study of non-exchangeable penalties in a more general context is also a promising area of research. In particular whether a discriminative partition in blocks can be inferred instead of assuming it is given by external information.  

We also observe that our results bear similarities with conclusions from on related problems. In Section \ref{sec:examples}, we discussed how the benefits of externally-informed selectors depend on how discriminative is the external information. \cite{scarlett2012compressed} similarly showed that support recovery in compressed sensing becomes easier as the entropy of the distribution of inclusion probabilities across blocks shrinks. In hypothesis testing, \cite{genovese2006} demonstrated that the power of their informed weighted testing procedure increases when assignment of weights is positively associated with the null hypotheses being false. We expect analogous results to be achievable in other related problems like structural learning in graphical models or variable selection over partially ordered sets as introduced in \cite{taeb}, which are further interesting future research avenues. 


%%%%%%%%%%%%%%%%%%%%%%%%%%%%%%%%%%%%%%%%%%%%%%
%% Single Appendix:                         %%
%%%%%%%%%%%%%%%%%%%%%%%%%%%%%%%%%%%%%%%%%%%%%%
%\begin{appendix}
%\section*{???}%% if no title is needed, leave empty \section*{}.
%\end{appendix}
%%%%%%%%%%%%%%%%%%%%%%%%%%%%%%%%%%%%%%%%%%%%%%
%% Multiple Appendixes:                     %%
%%%%%%%%%%%%%%%%%%%%%%%%%%%%%%%%%%%%%%%%%%%%%%
%\begin{appendix}
%\section{???}
%
%\section{???}
%
%\end{appendix}

%%%%%%%%%%%%%%%%%%%%%%%%%%%%%%%%%%%%%%%%%%%%%%
%% Support information, if any,             %%
%% should be provided in the                %%
%% Acknowledgements section.                %%
%%%%%%%%%%%%%%%%%%%%%%%%%%%%%%%%%%%%%%%%%%%%%%
\begin{acks}[Acknowledgments]
%The authors would like to thank the anonymous referees, an XXXXX and the XXXXX for their constructive comments. 
Paul Rognon-Vael is also affiliated with the Department of Statistics and Operational Research, Universitat Politècnica de Catalunya. David Rossell is also affiliated to the Data Science Center, Barcelona School of Economics. Piotr Zwiernik is also affiliated with the Department of Economics and Business, Universitat Pompeu Fabra and Barcelona School of
Economics.

\end{acks}
%%%%%%%%%%%%%%%%%%%%%%%%%%%%%%%%%%%%%%%%%%%%%%
%% Funding information, if any,             %%
%% should be provided in the                %%
%% funding section.                         %%
%%%%%%%%%%%%%%%%%%%%%%%%%%%%%%%%%%%%%%%%%%%%%%
\begin{funding}
PRV acknowledges the support of Departament de Recerca i Universitats de la Generalitat de Catalunya.  DR was funded by grant Consolidaci\'on investigadora CNS2022-135963 by the AEI, and PID2022-138268NB-I00 by MCIN/AEI/10.13039/501100011033 /FEDER.
%Europa Excelencia EUR2020-112096 from the AEI/10.13039/501100011033 and European Union NextGenerationEU/PRT.
PZ was supported by NSERC grant RGPIN-2023-03481.

\end{funding}

%%%%%%%%%%%%%%%%%%%%%%%%%%%%%%%%%%%%%%%%%%%%%%
%% Supplementary Material, including data   %%
%% sets and code, should be provided in     %%
%% {supplement} environment with title      %%
%% and short description. It cannot be      %%
%% available exclusively as external link.  %%
%% All Supplementary Material must be       %%
%% available to the reader on Project       %%
%% Euclid with the published article.       %%
%%%%%%%%%%%%%%%%%%%%%%%%%%%%%%%%%%%%%%%%%%%%%%
%\begin{supplement}
%\stitle{???}
%\sdescription{???.}
%\end{supplement}

%%%%%%%%%%%%%%%%%%%%%%%%%%%%%%%%%%%%%%%%%%%%%%%%%%%%%%%%%%%%%
%%                  The Bibliography                       %%
%%                                                         %%
%%  imsart-???.bst  will be used to                        %%
%%  create a .BBL file for submission.                     %%
%%                                                         %%
%%  Note that the displayed Bibliography will not          %%
%%  necessarily be rendered by Latex exactly as specified  %%
%%  in the online Instructions for Authors.                %%
%%                                                         %%
%%  MR numbers will be added by VTeX.                      %%
%%                                                         %%
%%  Use \cite{...} to cite references in text.             %%
%%                                                         %%
%%%%%%%%%%%%%%%%%%%%%%%%%%%%%%%%%%%%%%%%%%%%%%%%%%%%%%%%%%%%%

%% if your bibliography is in bibtex format, uncomment commands:
\bibliographystyle{imsart-number} % Style BST file (imsart-number.bst or imsart-nameyear.bst)
\bibliography{bibliography.bib}       % Bibliography file (usually '*.bib')

%% or include bibliography directly:
% \begin{thebibliography}{}
% \bibitem{b1}
% \end{thebibliography}
\section{Auxiliary results}\label{auxiliary}

In this section we collect some technical results. Their proofs are gathered in section~\ref{sec:proofauxi}. 

\subsection{Solutions to penalized likelihood problems}

\begin{lem}\label{lem:l1isthreshold}
	In the sequence model \eqref{eq:GSM}, selecting the non-zero entries  of the solution to 
\begin{equation}\label{eq:genl1}
	\underset{\bs \beta \in \R^p}{\rm minimize} \;\;\tfrac12\|\bs y-\sqrt{n}\bs\beta\|^2+\sum_{j=1}^b \lambda_j \sum_{i \in B_j}|\beta_{i}|
\end{equation}
for non-negative $\lambda_1,\ldots,\lambda_b$  is equivalent to taking the selector $\hat S^{b}$ in \eqref{prop:isthresh} with $\bs \tau=(\lambda_1/n,\ldots,\lambda_b/n)$. Let $\bs\beta^\circ$ be either the MLE or the LASSO estimator with penalty $\lambda^\circ$. Selecting the non-zero entries of the solution to the problem
\begin{equation}\label{eq:genl1ada}
	\underset{\bs \beta \in \R^p}{\rm minimize} \;\;\tfrac12\|\bs y-\sqrt{n}\bs\beta\|^2+\sum_{j=1}^b \lambda_j \sum_{i \in B_j}\tfrac{|\beta_{i}|}{|\beta_i^\circ|}.
\end{equation}
is equivalent to taking the selector $\hat S^{b}$ with $\bs \tau=(\sqrt{\lambda_1/n},\ldots,\sqrt{\lambda_b/n})$ if $\bs\beta^\circ$ is the MLE and with $\bs \tau=(\lambda^\circ/2n+\sqrt{(\lambda^\circ/2n)^2+\lambda_1/n},\ldots,\lambda^\circ/2n+\sqrt{(\lambda^\circ/2n)^2+\lambda_b/n})$ if $\bs\beta^\circ$ is a LASSO estimate.
\end{lem}

\subsection{Tail bounds}

\begin{lem}\label{lemma:3tailbounds}\textbf{Tail bounds on the maximum and minimum of folded Gaussians} 
\begin{longlist}
        \item If $\bs y\sim N_p\left(0,\frac{1}{n}I_p \right)$, $p>1$, and $a \geq \sqrt{\frac{2\ln(p)}{n}}$, $$P\Big(\underset{i\in\{1,\ldots,p\}}{\max}|y_i|> a \Big) \leq \frac{e^{-\frac{n}{2}\left(a^2-\frac{2\ln (p)}{n}\right)}}{\sqrt{\pi\ln (p)}}.$$
    \item If $\bs y\sim N_s\left(\bs\mu,\sigma^2I_s\right)$, $s>1$, and $a \leq \underset{i\in\{1,\ldots,s\}}{\min}|\mu_i|$,
        $$P\Big(\underset{i\in\{1,\ldots,s\}}{\min}|y_i| > a\Big) \geq P\Big(\underset{i\in\{1,\ldots,s\}}{\min}|\mu_i|-\underset{i\in\{1,\ldots,s\}}{\max}|y_i-\mu_i| > a \Big).$$
    \item If $y_i\sim N\left(\mu,\frac{1}{n}\right)$ for $i=1,\ldots,s$ and $\mu \in \mathbb{R}$, and $a<|\mu|$, then
        $$P\Big(\underset{i\in\{1,\ldots,s\}}{\min}|y_i| > a \Big) \leq \exp\left\{-\frac{s}{2}\frac{e^{-2n(|\mu|-a)^2/\pi}-e^{-n(a+|\mu|)^2/2}}{\left(1-e^{-2n(|\mu|-a)^2/\pi}\right)^{\frac{1}{2}}+\left(1-e^{-n(a+|\mu|)^2/2}\right)^{\frac{1}{2}}}\right\}.$$
\end{longlist}
\end{lem}

\begin{lem}\label{lem:boundrhogen}
 For any $T\subseteq S$ and $M\in\M$ such that $T\not\subseteq M$, let $Q_T=M \cup T$. The non-centrality parameter defined in \eqref{eq:muQM} satisfies: 
\begin{equation}
    \mu_{Q_TM}\; \geq \;n\,\rho(\bs X)\,\sum_{j=1}^b|T_j\setminus M_j|\,\min_{i\in T_j\setminus M_j}{\beta_{i}^*}^2.
\end{equation}
\end{lem}

\begin{lem}\label{lemma:s1s3}
Let $W\sim\chi_{\nu}^2(\mu)$ with $\mu \geq 0$, then for any $w>\mu+\nu$
$$P(W > w) \leq e^{-\big(\frac{w+\mu}{2}  - \sqrt{2w( 2 \mu+ \nu)-2 \mu \nu - \nu^2 } \big)}. $$
Moreover, if $w$ is increasing such that $\nu=o(w)$ and $\mu=o(w)$, for any $\phi\in(0,1)$ and $w$ large enough
$$P(W > w)  \leq  e^{-\phi\frac{w}{2}}.$$
\end{lem}

\begin{lem}\label{lemma:s2} Let $W \sim \chi_\nu^2(\mu)$ with $\mu>0$. For any $w<\mu$, 
$$
P(W<w) \leq \frac{e^{-\frac{1}{2}(\sqrt{\mu}-\sqrt{w})^2}}{(\mu / w)^{\nu / 4}}.$$
\end{lem}

\begin{lem}\label{lemma:s18} 
    Let $W\sim \chi_\nu^2(\mu)$ with $\mu \geq 0$. Assume $g$ positive and increasing such that $\nu=o(\ln(g))$, $\mu=o(\ln(g))$. Let $\bar{u},\underline{u}$ in $(0,1)$ such that $1>\bar{u}>\underline{u}\geq\left(1+g^{\phi}e^{-(\nu+\mu)/2}\right)^{-1}$ where $\phi\in(0,1)$, then for $g$ large enough, we have 
$$
\int_{\underline{u}}^{\bar{u}} P\left(W>2\ln \left(\frac{g}{1 / u-1}\right)\right) d u \leq  \frac{1}{g^{\phi}}\Big(\bar{u}-\underline{u}+\ln\Big(\frac{\bar{u}}{\underline{u}}\Big)\Big).
$$

\end{lem}

\begin{lem}\label{lemma:s20}
 Let $W\sim \chi_\nu^2(\mu)$ with $\mu \geq 0$. Assume $g$ positive and increasing such that $\nu=o(\ln(g))$, $\mu=o(\ln(g))$, then for any $\alpha\in(0,1)$ and $g$ large enough, we have
$$
\int_0^1 P\left(W>2 \ln \left(\frac{g}{1 / u-1}\right)\right) d u\;\; \leq\;\;  g^{-\alpha}$$
\end{lem}


\subsection{Bounds related to model normalized scores}

Let $NC(M)$ be the normalized score for model $M$ defined in \eqref{eq:normalizedscore}, $\M$ the set of models under consideration, and $\mu_{Q M}$ be the noncentrality parameter for any two nested models $Q \supseteq M$ defined in Lemma~\ref{lemma:s7}. Lemma~\ref{lem:L0toL1convergencegen} generalizes Lemma~\ref{lem:L0toL1convergence} and shows that the probability of not selecting a set of models is bounded above by the expected sum of the normalized scores of the models outside the set. Lemma~\ref{lem:boundexpnc} provides a bound on the expected normalized score of any $M \in \M$. Lemma~\ref{lem:upperboundnoncentralgen} gives, for any $T\subseteq S$ and $M\in \M$, upper and lower bounds on $\mu_{Q_TT}$ where $Q_T=M\cup T$.

\begin{lem}\label{lem:L0toL1convergencegen} For $\hat S^b$ as in \eqref{eq:Mhat} and any $k$ models $M_1,\ldots,M_k$
$$P(\hat{S}^b\not\in \{M_1,\ldots,M_k\})\;\leq\; (k+1) \sum_{M \in \M\setminus \{M_1,\ldots,M_k\}}\E\left( NC(M)\right).$$
\end{lem}

\begin{lem}\label{lem:boundexpnc} For any $T\subseteq S$ and $M\in \M\setminus\{T\}$, denote $Q_T= M\cup T$ and $A_T:=\gamma\Delta_{MT}+\tfrac{1-\gamma}{6}\mu_{Q_TM}$ (cf \eqref{eq:DeltaM} and \eqref{eq:muQM}). Suppose that, for some $\gamma \in (1/2,1]$, $A_T>0$, $|M\setminus T|=o(A_T)$, and $\mu_{Q_T T}=o(A_T)$. For any $\psi\in(0,1)$ and $n$ large enough:
$$\E\left(NC(M)\right)  \;\;\leq\;\; e^{-\psi A_T}.$$
\end{lem}

\begin{lem}\label{lem:upperboundnoncentralgen}
    For any $T\subseteq S$ and $M\in \M$, let $Q_T=M\cup T$, $$\mu_{Q_TT} \leq n\,\bar{\lambda}\,\sum_{j=1}^b|(S_j\cap M_j )\setminus T_j|\,\max_{i\in (S_j\cap M_j )\setminus T_j}{\beta_{i}^*}^2 \quad \text{where}\quad \bar{\lambda}\,:=\,\lambda_{\max}\Big(\frac1n \bs X_{S}^\top\bs X_{S}\Big)$$
    and
    $$\mu_{Q_TT} \geq n\,\lambda_{\min}\Big(\frac1n \bs X_{S}^\top\bs X_{S}\Big)\,\sum_{j=1}^b|(S_j\cap M_j )\setminus T_j|\,\min_{i\in (S_j\cap M_j )\setminus T_j}{\beta_{i}^*}^2$$
\end{lem}

\section{Proofs of Section~\ref{sec:properties_orth}}
 
 

\subsection{Proof of Proposition~\ref{prop:isthresh}}
First note that by discarding constant term 
 $$\arg\max_{M\in \M} \left\{\max_{\beta\in \mathcal{L}_M}\ell(\bs y; \bs \beta)-\sum_{j=1}^b\kappa_j |M_j| \right\}=\arg\min_{M\in \M} \left\{\min_{\beta\in \mathcal{L}_M}\bigg\{ \tfrac{1}{2}\|\bs y-\sqrt{n}\bs\beta\|^2\right\}+\sum_{j=1}^b\kappa_j |M_j| \bigg\}$$

 We also have $\min_{\beta\in \mathcal{L}_M} \tfrac{1}{2}\|\bs y-\sqrt{n}\bs\beta\|^2=\tfrac{1}{2}\|\bs y-\sqrt{n}\tilde{\bs \beta}_M\|^2=\tfrac{1}{2}\|\bs y\|^2-\tfrac{1}{2}\|\sqrt{n}\tilde{\bs \beta}_M\|^2$. The maximization in \eqref{eq:Mhat} can be then replaced with:
     \begin{equation}\label{eq:MhatGSM}
  \arg\max_{M\in \M} \bigg\{\tfrac{1}{2}\| \sqrt{n}\tilde{\bs \beta}_M\|^2-\sum_{j=1}^b\kappa_j |M_j|\bigg\}, \qquad M=M_1\cup \ldots\cup M_b.	
\end{equation}
Under the assumptions of \eqref{eq:GSM}, we can write 
	$$
	\tfrac{1}{2}\|\sqrt{n}\tilde{\bs \beta}_M\|^2-\sum_{j=1}^b\kappa_j |M_j|\;=\;\sum_{j=1}^b\sum_{i\in M_j}(\tfrac{n}{2}{\tilde{ \beta}_i}^2-\kappa_j).
	$$
	Then \eqref{eq:MhatGSM} can be maximized with respect to each $M_j$ by including $i\in \hat S_j$ whenever $n\tilde\beta_i^2\geq 2\kappa_j$.

\subsection{Proof of Proposition~\ref{prop:suffnecIblockthresh}} 

\subsubsection{Part (i)} 
By the union bound and by Lemma~\ref{lemma:3tailbounds} (i),
\begin{equation}\label{eq:type1convrate}
P\left(\hat{S}^{b} \not\subseteq S\right)\;\leq\; \sum_{j=1}^b P\Big(\max_{i\in B_j\setminus S_j}|y_i/\sqrt{n}|>\tau_{j}\Big)\;\leq\;\sum_{j=1}^b\frac{e^{-\tfrac{n}{2}\left(\tau_j^2-\tfrac{2\ln(p_j-s_j)}{n}\right)}}{\sqrt{\pi\ln(p_j-s_j)}} 	.
 \end{equation}
 By A4, the numerator on the right is bounded above by $1$ for all sufficiently large $n$. It follows that 
	$$P\left(\hat{S}^{b} \not\subseteq S\right) \;\leq\; \sum_{j=1}^b (\pi\ln (p_j-s_j))^{-1/2} \to 0\qquad \mbox{as } n\to \infty.$$
\subsubsection{Part (ii)}  By independence, we have 
$$P(\hat{S}^{b} \subseteq S)\;=\;\prod_{j=1}^b P(\max_{i \in B_j\setminus S_j}|y_i/\sqrt{n}| \leq \tau_j).$$ 
Consider $j$ such that $\lim_{n\to \infty} \frac{\sqrt{n}\tau_j}{\sqrt{2\ln(p_j-s_j)}} \;<\;1$. In particular, there exists $c<1$ such that, for sufficiently large $n$, we have that $\tau_j\leq c\sqrt{2\ln(p_j-s_j)/n}$. Then, for any such $n$,
$$
P\Big(\max_{i \in B_j\setminus S_j}|y_i/\sqrt{n}| \leq \tau_j\Big)\;\leq\; P\Big(\max_{i \in B_j\setminus S_j}|y_i| \leq c \sqrt{2\ln(p_j-s_j)}\Big)\;\leq\;P\Big(\tfrac{1}{\sqrt{2\ln(p_j-s_j)}}\max_{i \in B_j\setminus S_j}y_i \leq c \Big).
$$
On the other hand, results from extreme value theory\cite{galambos} Example 4.4.1, show that 
 $$\frac{1}{\sqrt{2 \ln(p_j-s_j)}}\max_{i\in B_j\setminus S_j}y_i \;\overset{P}{\longrightarrow} 1$$
 and so 
 $$
 P\Big(\frac{1}{\sqrt{2\ln(p_j-s_j)}}\max_{i \in B_j\setminus S_j}y_i \leq c \Big)\;\to\;0,
 $$ 
 which implies that $P(\hat{S}^{b} \subseteq S)\to 0$. 

\subsection{Proof of Proposition~\ref{prop:suffnecIIblockthresh}}
	 By the union bound, 
	$P\big(\hat{S}^{b} \not\supseteq S\big)\leq \sum_{j=1}^b  P\left(\min_{i\in S_j}|y_i/\sqrt{n}|\leq\tau_j\right)$. 
	By Lemma~\ref{lemma:3tailbounds} (ii), for each $j$,	
	$$P\left(\min_{i\in S_j}|y_i/\sqrt{n}|\leq\tau_j\right)\;\leq\; P\left(\max_{i\in S_j}|y_i/\sqrt{n}-\beta_i| \geq\beta_{\min,j}^*-\tau_j \right).$$ By Lemma~\ref{lemma:3tailbounds} (i)  
	\begin{equation}\label{eq:type2convrate}
P\left(\hat{S}^{b}\not\supseteq S \right) \;\leq\;\sum_{j=1}^b\frac{e^{-\tfrac{n}{2}\left((\beta^\star_{\min,j}-\tau_j)^2-\tfrac{2\ln(s_j)}{n}\right)}}{\sqrt{\pi\ln(s_j)}}.		
	\end{equation}
 By A5, the nominator on the right is bounded above by $1$ for all sufficiently large $n$. It follows that, by A3 
	$$P\left(\hat{S}^{b} \not\supseteq S \right) \;\leq\;  \sum_{j=1}^b (\pi\ln (s_j))^{-1/2} \to 0\qquad \mbox{as } n\to \infty.$$ 

\subsection{Proof of Proposition~\ref{prop:necIIblockthresh}} 

\subsubsection{Part (i)} By independence we have that
$$P(\hat{S}^{b}\supseteq S)=\prod_{j=1}^b P(\min_{i \in S_j}|y_i/\sqrt{n}| > \tau_j).$$ 
where
$$
P(\min_{i \in S_j}|y_i/\sqrt{n}| > \tau_j)=\prod_{i \in B_j} P(|y_i/\sqrt{n}| > \tau_j). 
$$
Take any $j$, satisfying \eqref{cond:necIIallblockthresh2}. 
Denote by $i_{\circ}\in B_j$ an entry such that $|\beta_{i_{\circ}}^*|=\beta^*_{\min,\,j}$. In the proof of Lemma \ref{lemma:3tailbounds} (iii), we showed that if $y\sim N(\mu, \sigma^2)$, then for any $a<|\mu|$ , $P(|y|>a)  = P(z>\frac{a-|\mu|}{\sigma}) + P(z<-\frac{a+|\mu|}{\sigma})$ where $z\sim N(0,1)$. We have then
$$P(|y_{i_{\circ}}/\sqrt{n}|>\tau_j) 
    = P\left(z>\sqrt{n}\left(\tau_j-\beta_{\min,\,j}^*\right)\right) + P\left(z<-\sqrt{n}\left(\tau_j+\beta_{\min,\,j}^*\right)\right)$$ where $z\sim N(0,1)$. Since $\tau_j<\beta^\star_{\min,j}$ satisfies condition \eqref{cond:necIallblockthresh}, $\sqrt{n}\left(\tau_j+\beta_{\min,\,j}^*\right)\to\infty$ and we have that $P(z<-\sqrt{n}\left(\tau_j+\beta_{\min,\,j}^*\right)) \to 0$.
    
    Further,  by \eqref{cond:necIIallblockthresh2}  we have that $\underset{n \to +\infty}{\lim}P\left(z>\sqrt{n}\left(\tau_j-\beta_{\min,\,j}^*\right)\right)<1$ and hence that $\underset{n \to +\infty}{\lim}P\left(S \subseteq \hat{S}^{b} \right)<1$, as we wished to prove.
    
\subsubsection{Part (ii)}	
Take any $j=1,\ldots,b$ satisfying \eqref{cond:necIIallblockthresh}. By Lemma~\ref{lemma:3tailbounds} (iii), since $\beta^\star_{\min,j}>\tau_j$:
$$P\left(\min_{i \in S_j}\Big|\frac{y_i}{\sqrt{n}}\Big| > \tau_j \right)\;\leq\;
\exp\left\{-\frac{s_j}{2}\frac{e^{-\tfrac{2n}{\pi}(\beta^\star_{\min,j}-\tau_j)^2}-e^{-\tfrac{n}{2}(\tau_j+\beta^\star_{\min,j})^2}}{\left(1-e^{-2n(\beta^\star_{\min,j}-\tau_j)^2/\pi}\right)^{\frac{1}{2}}+\left(1-e^{-n(\tau_j+\beta^\star_{\min,j})^2/2}\right)^{\frac{1}{2}}}\right\}. $$
Let $a_{n}=2\left(1-e^{-2n(\beta^\star_{\min,j}-\tau_j)^2/\pi}\right)^{\frac{1}{2}}+2\left(1-e^{-n(\tau_j+\beta^\star_{\min,j})^2/2}\right)^{\frac{1}{2}}$ and note that $a_n\in (0, 4]$.  Thus, to show that $P\left(\min_{i \in S_j}|\tfrac{y_i}{\sqrt{n}}| > \tau_j \right)$ is bounded away from $1$, it is enough to show that 
$$
\lim_{n\to \infty} \frac{s_j}{a_n}({e^{-2n(\beta^\star_{\min,j}-\tau_j)^2/\pi}-e^{-n(\tau_j+\beta^\star_{\min,j})^2/2}})
\geq \lim_{n\to \infty} \frac{s_j}{4}({e^{-2n(\beta^\star_{\min,j}-\tau_j)^2/\pi}-e^{-n(\tau_j+\beta^\star_{\min,j})^2/2}})
\;>\;0.
$$

By \eqref{cond:necIIallblockthresh}, we have that $\lim_{n\to\infty} {s_j} e^{-2n(\beta^\star_{\min,j}-\tau_j)^2/\pi}\geq 1$. To conclude, we show that
$$
\lim_{n\to\infty} {s_j} e^{-n(\tau_j+\beta^\star_{\min,j})^2/2}\;\leq\;c\;<\;1
$$ 
for some $c\in(0,1)$ such that $\tfrac{s_j}{p_j}\leq c$. Such $c$ exists by our assumption $\tfrac{s_j}{p_j}<1$.
We equivalently need
$\sqrt{n}(\tau_j+\beta^\star_{\min,j})\geq\sqrt{{2\ln (s_j/c)}}$ for all $n$ sufficiently large. To show that, note that by assumption $\tau_j< \beta^\star_{\min,j}$ and $s_j/c\leq p_j$, which gives
$$
\frac{\tau_j+\beta^\star_{\min,j}}{\sqrt{\tfrac{2\ln(s_j/c)}{n}}}\;\geq\;\frac{2\tau_j}{\sqrt{\tfrac{2\ln(p_j)}{n}}}\;=\;\frac{\sqrt{\tfrac{2\ln(p_j-s_j)}{n}}}{\sqrt{\tfrac{2\ln(p_j)}{n}}}\frac{2\tau_j}{\sqrt{\tfrac{2\ln(p_j-s_j)}{n}}}\;=\;\sqrt{\frac{\ln(p_j-s_j)}{\ln(p_j)}}\frac{2\tau_j}{\sqrt{\tfrac{2\ln(p_j-s_j)}{n}}}.
$$
Since by assumption $\tau_j$ satisfies \eqref{cond:necIallblockthresh}, the second term on the right converges to something $\geq 2$, and it is enough to show that the first term converges to something $>1/2$. Using that $\tfrac{s_j}{p_j}<c$, we get
$$
\sqrt{\frac{\ln(p_j-s_j)}{\ln(p_j)}}
= \sqrt{\frac{\ln(p_j) + \ln(1-\frac{s_j}{p_j})}{\ln(p_j)}}
\;\geq\;\sqrt{\frac{\ln(p_j)+\ln(1-c)}{\ln(p_j)}}\;\to\;1,
$$
which concludes the proof.

\subsection{Proof of Lemma~\ref{lem:necbetaminblockthresh}}
    We have $\frac{\beta_{\min,j}^*}{\sqrt{\tfrac{2\ln(p_j-s_j)}{n}}+\sqrt{\tfrac{\pi}{2}\tfrac{\ln(s_j)}{n}}}= \frac{\beta_{\min,j}^*-\tau_j+\tau_j}{\sqrt{\tfrac{2\ln(p_j-s_j)}{n}}+\sqrt{\tfrac{\pi}{2}\tfrac{\ln(s_j)}{n}}}$ and \eqref{eq:necbetaminblockthresh} implies that either $\lim_{n\to \infty}\frac{\tau_j}{\sqrt{\tfrac{2\ln(p_j-s_j)}{n}}} <1$ or $\lim_{n\to \infty}\frac{\beta_{\min,j}^*-\tau_j}{\sqrt{\tfrac{\pi}{2}\tfrac{\ln(s_j)}{n}}} <1$. 
    If $\lim_{n\to \infty}\frac{\tau_j}{\sqrt{\tfrac{2\ln(p_j-s_j)}{n}}} <1$ by Proposition~\ref{prop:suffnecIblockthresh}, $\lim_{n \to \infty} P(\hat{S}^{b} = S) < 1$. If $\beta_{\min,j}<\tau_j$ recovery is not possible either by definition.
    The last case to examine is when $\beta_{\min,j}\geq\tau_j$, $\lim_{n\to \infty}\frac{\tau_j}{\sqrt{\tfrac{2\ln(p_j-s_j)}{n}}} \geq 1$ and $\lim_{n\to \infty}\frac{\beta_{\min,j}^*-\tau_j}{\sqrt{\tfrac{\pi}{2}\tfrac{\ln(s_j)}{n}}} <1$. By Proposition~\ref{prop:necIIblockthresh}, in that case $\lim_{n \to \infty} P(\hat{S}^{b} \supseteq S ) < 1$ and recovery is not possible asymptotically.

\subsection{Proof of Theorem~\ref{theo:fullrecovGS}}
By the union bound and the bounds in \eqref{eq:type1convrate} and \eqref{eq:type2convrate}, under A4 and A5, we have:
\begin{equation}\label{eq:fhfhfhf}
P(S\neq \hat{S}^b) \; \leq\;\sum_{j=1}^b\frac{e^{-\tfrac{n}{2}\left(\tau_j^2-\tfrac{2\ln(p_j-s_j)}{n}\right)}}{\sqrt{\pi\ln(p_j-s_j)}}+\sum_{j=1}^b\frac{e^{-\tfrac{n}{2}\left((\beta^\star_{\min,j}-\tau_j)^2-\tfrac{2\ln(s_j)}{n}\right)}}{\sqrt{\pi\ln(s_j)}}.
\end{equation}
We continue with showing that if \eqref{eq:betamin} holds then $\tau^\star_j$ satisfies A4 and A5. Under \eqref{eq:betamin}, there exists $c>1$ such that $\beta^\star_{\min,j}=c(\sqrt{\tfrac{2\ln(p_j-s_j)}{n}} + \sqrt{\tfrac{2\ln(s_j)}{n}})$. Denote $a=\sqrt{\tfrac{2\ln(p_j-s_j)}{n}}$ and $b=\sqrt{\tfrac{2\ln(s_j)}{n}}$. We have
$$
\tau_j^\star\;=\;\frac{c}{2} (a+b)+\frac{(a^2-b^2)}{2c(a+b)}\;=\;\frac{c}{2} (a+b)+\frac{a-b}{2c}\;=\;\frac{c^2+1}{2c}a+\frac{c^2-1}{2c}b.
$$
The fact that $\tau^\star_j>a$ (the lower bound in A4) is now clear because $\tfrac{c^2+1}{2c}> 1$, $\tfrac{c^2-1}{2c}> 0$ and $b \geq 0$. To show that the upper bound in A5 holds observe that
$$
\beta_{\min,j}^* - \tau^\star_j \;=\; c(a+b)- \frac{c^2+1}{2c}a-\frac{c^2-1}{2c}b \;= \; \frac{c^2-1}{2c}a+\frac{c^2+1}{2c}b\;>\;  b.
$$
Then for $n$ sufficiently large \eqref{eq:rateconvgsm} holds. For convenience, denote $d = \frac{1}{n {\beta_{\min}^*}_j} \ln (p_j / s_j-1)$ such that $\tau^*_j= {\beta_{\min}^*}_j/2 + d$. Then
$$e^{-\frac{n}{2}\left({\tau^*_j}^2-\frac{2\ln(p_j-s_j)}{n}\right)}
=e^{-\left[\frac{n}{8}{\beta_{\min}^*}_j^2+\frac{n}{2}d^2 + \frac{1}{2}\left(\ln (p_j- s_j)-\ln (s_j)\right) -\ln (p_j- s_j)\right]}$$
Then, by considering separately the two possible maxima in $\max\{\ln(p_j-s_j),\ln(s_j)\}$, we get that
$$\frac{e^{-\frac{n}{2}\left({\tau^*_j}^2-\frac{2\ln(p_j-s_j)}{n}\right)}}{e^{-[\frac{n}{8}{{\beta_{\min}^*}_j}^2-\max\{\ln(p_j-s_j),\ln(s_j)\}]}}= 
e^{-\left[\frac{n}{2}d^2 + \frac{1}{2}\left|\ln (p_j- s_j)-\ln (s_j)\right|\right]} <1.$$
Note that $e^{-\tfrac{n}{2}\big[{\tau_j^*}^2-\tfrac{2\ln(p_j-s_j)}{n}\big]}=e^{-\tfrac{n}{2}\big[(\beta^\star_{\min,j}-{\tau_j^*})^2-\tfrac{2\ln(s_j)}{n}\big]}$,  $\sqrt{\pi\ln(p_j-s_j)}>1$, and $\sqrt{\pi\ln(s_j)}>1$. Then from \eqref{eq:fhfhfhf} we have:
$$P(\hat{S}^{b} \neq S) \; \leq \; 2 \sum_{j=1}^be^{-[\frac{n}{8}{\beta^\star_{\min,j}}^2-\max\{\ln(p_j-s_j),\ln(s_j)\}]}.$$

\subsection{Proof of Corollary~\ref{cor:selconstimpossible}}

Since $\hat{S}$ is $\hat{S}^b$ with $b=1$, the assumptions of Lemma~\ref{lem:necbetaminblockthresh} for $\hat{S}$ are met and $\lim_{n \to \infty} P(\hat{S} = S) < 1$. Since A4 and A5 hold, by Proposition~\ref{prop:suffnecIblockthresh} and Proposition~\ref{prop:suffnecIIblockthresh}, $\lim_{n \to \infty} P(\hat{S}^b \not\subseteq S)= 0$ and $\lim_{n \to \infty} P(\hat{S}^b \not\supseteq S)= 0$, and then $\lim_{n \to \infty} P(\hat{S}^b = S)= 1$.

\section{Proofs of Section~\ref{sec:properties_legression}}
\subsection{Proof of Lemma~\ref{lem:useful-reformulation}}
	First note that by discarding constant term 
 $$\arg\max_{M\in \M} \left\{\max_{\beta\in \mathcal{L}_M}\ell(\bs y; \bs \beta)-\sum_{j=1}^b\kappa_j |M_j| \right\}=\arg\min_{M\in \M} \left\{\min_{\beta\in \mathcal{L}_M}\bigg\{ \tfrac{1}{2}\|\bs y-\bs X\bs\beta\|^2\right\}+\sum_{j=1}^b\kappa_j |M_j| \bigg\}$$
 
 We also have $\min_{\beta\in \mathcal{L}_M} \tfrac{1}{2}\|\bs y-\bs X\bs\beta\|^2=\tfrac{1}{2}\|\bs y-\bs X_{\!\!M}\tilde{\bs \beta}^{(M)}\|^2=\tfrac{1}{2}\|\bs y\|^2-\tfrac{1}{2}\|\bs X_{\!\!M}\tilde{\bs\beta}^{(M)}\|^2$. The maximization in \eqref{eq:Mhat} can be then replaced with maximization of $\tfrac{1}{2}
 \|\bs X_{\!\!M}\tilde{\bs\beta}^{(M)}\|^2-\sum_{j=1}^b \kappa_j|M_j|$.

\subsection{Proof of Lemma~\ref{lem:L0toL1convergence}}
The first statement follows directly from Lemma~\ref{lem:L0toL1convergencegen} by taking $\{M_1,\ldots,M_k\}=\{S\}$. For the second statement, note that  $NC(M)=(1+\sum_{N \neq M} e^{C(N)}/ e^{C(M)})^{-1}$. Since for any $M,N\in \M$ and $M\neq N$, $C(N) / C(M)>0$, $NC(M) \leq (1+e^{C(M')-C(M)})^{-1}$.

\subsection{Proof of Lemma~\ref{lemma:s7}}
    This result follows directly from Lemma~S7 in \cite{rossell2022concentration} taking $\phi^*=1$.

\subsection{Proof of Lemma~\ref{lem:boundrho}}
This result follows directly Lemma~\ref{lem:boundrhogen} by taking $T=S$, and observing that $\min_{i\in S_j\setminus M_j}{\beta^*_i}^2\geq {\beta^*_{\min,j}}^2$.

\subsection{Proof of Theorem~\ref{theo:suffcondlinearmodel}} We start by showing that, for any $M\in\mathcal{M}\setminus \{S\}$, $A_S$ as defined in Lemma~\ref{lem:boundexpnc} for $T=S$ satisfies the conditions of the lemma. Observe that $\Delta_{MS}$, defined in \eqref{eq:DeltaM}, can be rewritten as $\Delta_{MS}=\sum_{j=1}^b(|M_j\setminus S_j|-|S_j\setminus M_j|)\kappa_j$. By  Lemma~\ref{lem:boundrho}, for every $n\in \N$ we have 
\begin{equation}\label{eq:nonoverfittingcond}
   A_S\;=\; \gamma \Delta_{MS}+\frac{1-\gamma}{6}\mu_{Q_SM}\;\geq\; \gamma\sum_{j=1}^b|M_j \setminus S_j|\kappa_j + \sum_{j=1}^b |S_j\setminus M_j |\left(\tfrac{1-\gamma}{6} n \rho(\bs X) {\beta^*_{\min,j}}^{2} - \gamma\kappa_j\right) 
\end{equation} 
Since $M\neq S$, $|M\setminus S|\neq0$ or $|S\setminus M|\neq0$, then by assumptions A6 and A7, for $n$ large enough $A_S>0$. We immediately have $\mu_{Q_S S}=o(A_S)$ because $\beta^*_{Q_S\setminus S}=\beta^*_{M\setminus S}=0$ and $\mu_{Q_S S}=0$. If $|M\setminus S|=0$, $|M\setminus S|=o(A_S)$ also immediately. Consider now the case $|M\setminus S|\neq 0$. By assumption A7, the last term in \eqref{eq:nonoverfittingcond} is nonnegative and:
$$\frac{|M\setminus S|}{A_S}\leq \bigg[ \gamma\sum_{j=1}^b\frac{|M_j \setminus S_j|}{|M\setminus S|}\kappa_j \bigg]^{-1} \leq \Big[ \gamma\min_{j=1,\ldots,b}\kappa_j \Big]^{-1} $$
where the last inequality follows from $\sum_{j=1}^b\frac{|M_j \setminus S_j|}{|M\setminus S|}=1$. Then by assumption A6, when $|M\setminus S|=o(A_S)$ when $|M\setminus S|\neq 0$ too.
We conclude that for every $M\in \mathcal M\setminus \{S\}$, by Lemma~\ref{lem:boundexpnc}, for any $\psi\in(0,1)$ and $n$ large enough:
$$
\E(NC(M))\;\leq\;e^{-\psi A_S}.
$$
Let 
$$A_S^*\;:=\;\gamma\sum_{j=1}^b|M_j \setminus S_j|\kappa_j + \sum_{j=1}^b |S_j\setminus M_j |\left(\tfrac{1-\gamma}{6} n \rho(\bs X) {\beta^*_{\min,j}}^{2} - \gamma\kappa_j\right).$$
By \eqref{eq:nonoverfittingcond}, we have for $n$ large enough:
\begin{equation}\label{eq:boundexp}
    \E(NC(M))\;\leq\;e^{-\psi A_S^*}.
\end{equation}
Assumption~A7 implies there exist $g'_j \to +\infty$ such that:
$$\frac{(1-\gamma) n\rho(\bs X)}{6}{{\beta_{\min}^*}_j}^2 \,- \,\kappa_j \;=\; \ln(s_j) + g^{\prime}_{j}.$$
Let $\delta \in (0,1)$ and denote $\bar{m}_j=\max\big\{\frac{2\ln(p_j-s_j)}{f_j},\frac{2\ln(s_j)}{g'_{j}}\big\}$. Take $\psi=\max_{j=1,\ldots,b}\frac{\xi + \delta + \bar{m}_j}{1+\bar{m}_j}$ for some $\xi \in(0, 1- \delta)$ then $\psi \in (0,1)$ and we have for every $j=1,\ldots,b$:
$$\delta \leq \psi-(1-\psi)\bar{m}_j.$$ 
By definition of $\gamma$ we have: 
$$
\gamma\kappa_j\;\geq\;\frac12 \Big(1+\frac{\ln(p_j-s_j)}{\kappa_j}\Big)\kappa_j\;=\;\ln(p_j-s_j)+\frac12(\kappa_j-\ln(p_j-s_j)) \;=\; \ln(p_j-s_j)+ \frac12f_j.
$$
From our choice of $\psi$ and $\delta$, we get:
$$\psi\gamma\kappa_j \;\geq\; \psi\big(\ln(p_j-s_j)+ \frac12f_j\big) \;\geq\;  \ln(p_j-s_j)+ \delta\frac12f_j.$$
By A7 and by our choice of $\psi$ and $\delta$,
$$\psi\left(\tfrac{1-\gamma}{6} n \rho(\bs X) {\beta^*_{\min,j}}^{2} - \gamma\kappa_j\right) \;\geq\; \psi\left(\ln(s_j)+g_j'\right) \;\geq\; \ln(s_j)+\delta \frac{1}{2}g'_j.
$$
Using the expression for $A_S^*$, the bound in \eqref{eq:boundexp}, and the above inequalities, we get
$$
\E(NC(M))\;\leq\;\exp\left\{- \sum_{j=1}^b |M_j\setminus S_j|(\ln(p_j-s_j)+\delta \tfrac{f_j}{2})-\sum_{j=1}^b |S_j\setminus M_j|(\ln(s_j)+\delta \tfrac{g'_j}{2})\right\}.
$$

Denote $\mathcal{S}= \sum_{M \in \M\setminus \{S\}} \E\left(NC(M)\right)$ for convenience. We can split $\mathcal{S}$ into sums over the models that have the same number of spurious variables and same number of missing variables in every block, that is the models $M$ such that for all $j$, $|M_j\setminus S_j|=u_j$ and $|S_j\setminus M_j|=w_j$ with $u_j \in \{0,\ldots,p_j-s_j\}$ and $w_j \in \{0,\ldots,s_j\}$, not all simultaneously zero:
$$\mathcal{S} =\sum_{w_1=0}^{s_1}\cdots\sum_{w_b=0}^{s_b}\sum_{u_1=0}^{p_1-s_1}\cdots\sum_{u_b=0}^{p_b-s_b}\sum_{M \neq S:\forall \,j |M_j\setminus S_j|=u_j, |S_j\setminus M_j|=w_j } \E\left(NC(M)\right)$$
The number of models with given collection of $u_j$'s and $w_j$'s is $\prod_{j=1}^b {\binom{p_j-s_j}{u_j}}{\binom{s_j}{w_j}}$. We thus have:
\begin{equation*}
    \mathcal{S} \leq \prod_{j=1}^b\left(1+\sum_{u_j=1}^{p_j-s_j}{\binom{p_j-s_j}{u_j}}e^{- u_j(\ln(p_j-s_j)+\delta \tfrac{f_j}{2})}\right)\left(1+
    \sum_{w_j=1}^{s_j}{\binom{s_j}{w_j}}
    e^{-w_j(\ln(s_j)+\delta \tfrac{g'_j}{2})}\right) - 1
\end{equation*}
where $-1$ stems from discarding $S$ from the product-sum, which corresponds to the case $u_j=0$ and $w_j=0$ for all $j$. A standard bound on binomial coefficient for $1\leq k\leq n$ is:
\begin{equation}\label{eq:upperboudbinom}
    \binom{n}{k} \leq \left(\frac{n\,e}{k}\right)^k \leq \left(n\,e\right)^k
\end{equation}
Then
\begin{equation*}
    \mathcal{S} \;\;\leq\;\; \prod_{j=1}^b\left(1+\sum_{u_j=1}^{p_j-s_j}e^{-u_j\left(\delta \tfrac{f_j}{2}-1\right)}\right)\left(1+
    \sum_{w_j=1}^{s_j}    e^{-w_j\left(\delta \tfrac{g'_j}{2}-1\right)}\right) - 1
\end{equation*}
Denote
$$
d_j\;=\;e^{1-\delta \tfrac{f_j}{2}},\qquad g_j\;=\;e^{1-\delta \tfrac{g'_j}{2}}
$$
where both expressions go to zero as $n$ increases.
For every $j$, we have:
\begin{eqnarray*}
    &1+\sum_{u_j=1}^{p_j-s_j}e^{-u_j\left(\delta \tfrac{f_j}{2}-1\right)}
    =\frac{1-d_j^{p_j-s_j+1}}{1-d_j} \\
    &1+\sum_{w_j=1}^{s_j}    e^{-w_j\left(\delta \tfrac{g'_j}{2}-1\right)} =\frac{1-g_j^{s_j+1}}{1-g_j}.
\end{eqnarray*}
Since both expressions converge to $1$ as $n$ grows, we conclude $\mathcal S\to 0$.

 \subsection{Proof of Theorem~\ref{prop:fullrecovreg}}
From the proof of Theorem~\ref{theo:suffcondlinearmodel} we have:
\begin{equation}\label{eq:prodconstreg}
    \mathcal{S} \;\;\leq\;\; \prod_{j=1}^b\left(1+\sum_{u_j=1}^{p_j-s_j}e^{-u_j\left(\delta \tfrac{f_j}{2}-1\right)}\right)\left(1+
    \sum_{w_j=1}^{s_j}    e^{-w_j\left(\delta \tfrac{g'_j}{2}-1\right)}\right) - 1
\end{equation}
For every $j$, we have:
\begin{eqnarray*}
    &\sum_{u_j=1}^{p_j-s_j}e^{-u_j\left(\delta \tfrac{f_j}{2}-1\right)}
    =g_j\,\frac{1-g_j^{p_j-s_j}}{1-g_j} \\
    &\sum_{w_j=1}^{s_j}    e^{-w_j\left(\delta \tfrac{g'_j}{2}-1\right)} =g_j\,\frac{1-g_j^{s_j}}{1-g_j}.
\end{eqnarray*}
Each of the sum above is smaller than 1 for $n$ sufficiently large. Then every term of degree larger than 1 in the product in \eqref{eq:prodconstreg} is bounded above by the sum of terms of degree 1. For $n$ large enough, we get: 
 \begin{equation*} \mathcal{S} \leq (2^{2b}-2b)\sum_{j=1}^b \bigg[g_j\,\frac{1-g_j^{p_j-s_j}}{1-g_j}+ g_j\,\frac{1-g_j^{s_j}}{1-g_j}\bigg ]. \end{equation*} 
 Take $$r_b=\max_{j=1,\ldots,b}\bigg\{\frac{1-g_j^{p_j-s_j}}{1-g_j} \,,\,\frac{1-g_j^{s_j}}{1-g_j} \bigg\}.$$
 By Lemma~\ref{lem:L0toL1convergence} and observing that $2e \leq 6$, we get: 
 \begin{eqnarray*}
     P(\hat{S}^{b} \neq S)\;&\leq&\;  2 \sum_{M \in \M\setminus \{S\}} \E\left(NC(M)\right) \\\;&\leq&\; (2^{2b}-2b)6r_b\,\,\sum_{j=1}^b\,\, e^{-\tfrac{\delta}{2}\big[\kappa_j-\ln(p_j-s_j)\big]} + e^{-\tfrac{\delta}{2}\big[\frac{(1-\gamma) n\rho(\bs X)}{2}{{\beta_{\min}^*}_j}^2 - \kappa_j -\ln(s_j) \big]}
 \end{eqnarray*}
and 
 $$ P(\hat{S}^{b} \neq S)\;\leq \;(2^{2b}-2b)6r_b\sum_{j=1}^b e^{-\tfrac{\delta}{2}\big[\kappa_j-\ln(p_j-s_j)\big]} + e^{-\tfrac{\delta}{2}\big[(\sqrt{\frac{(1-\gamma) n\rho(\bs X)}{6}}{{\beta_{\min}^*}_j} - \sqrt{\kappa_j})^2 -\ln(s_j) \big]}.$$

We continue with showing that if \eqref{betamincond:minimalreg} holds then $\kappa^\star_j$ satisfies A6 and A7. Under \eqref{betamincond:minimalreg}, there exists $c$ such that $\lim_{n\to\infty} c>1$ and $\sqrt{\frac{(1-\gamma)n}{6}}\beta^\star_{\min,j}=c\big(\sqrt{\ln(p_j-s_j)} + \sqrt{\ln(s_j)}\big)$. Denote $a=\sqrt{\ln(p_j-s_j)}$ and $b=\sqrt{\ln(s_j)}$. Proceeding as in the proof of Theorem~\ref{theo:fullrecovGS} shows that for $n$ large enough, we have $\sqrt{\kappa^*_j} \geq  \big(\frac{(c-1)^2}{2c}+1\big)\sqrt{\ln(p_j-s_j)}$ and $\frac{\sqrt{(1-\gamma) n\rho(\bs X)}}{7}{{\beta_{\min}^*}_j}-\sqrt{\kappa^*} \geq \big(\frac{(c-1)^2}{2c}+1\big)\sqrt{\ln(s_j)}$
    which implies A6 and A7.
Then there exists $r_b\to 1$ such that for $n$ large enough the bound \eqref{eq:rateconvreg} on $P(\hat S^b \neq S)$ holds. Proceeding as in the proof of Theorem~\ref{theo:fullrecovGS} also shows that for any $j$ 
    $$e^{-\frac{\delta}{2}\left[{\kappa^*_j}-\ln(p_j-s_j)\right]}+e^{-\frac{\delta}{2}\big[(\sqrt{\frac{(1-\gamma) n\rho(\bs X)}{2}}{{\beta_{\min}^*}_j} - \sqrt{\kappa_j^*})^2 -\ln(s_j) \big]} <2 e^{-\frac{\delta}{2}[\frac{(1-\gamma) n\rho(\bs X)}{24}{{\beta_{\min}^*}_j}^2-\max\{\ln(p_j-s_j),\ln(s_j)\}]}. $$
from which we conclude \eqref{eq:oraclerateconvreg}.

\subsection{Proof of Proposition~\ref{prop:nectypeIreg}}

Recovery of $S$ requires that $\max_{M\in O_j}\frac{NC(M)}{NC(S)}<1$ with prob-
ability going to 1 as n grows. By the definition of the normalized criterion $NC$:
$$P(\hat{S}^b=S) \leq P\Big(\max_{M\in O_j}\frac{NC(M)}{NC(S)}<1\Big)\leq  P\big(\max_{M\in O_j}\sqrt{L_{MS}}<\sqrt{2\kappa_j}\Big).$$
By Lemma~\ref{lemma:s7}, for every $M\in O_j$, $L_{MS}\in \chi^2_1$ and there exists $Z_M\sim N(0,1)$ such that $\sqrt{L_{MS}}=|Z_M|$. Since for every $M\in O_j$, $|Z_M|\geq Z_M$, we have:
\begin{equation}\label{eq:upperboundrecov}
    P(\hat{S}^b=S) \leq  1-P\big(\max_{M\in O_j}Z_M\geq\sqrt{2\kappa_j}\Big).
\end{equation}

The set $O_j$ has cardinality $p_j-s_j$, then by \cite{hartigan}, Theorem 3.4  and assumption A2, for any $\varepsilon>0$:
$$P\big(\max_{M\in O_j}Z_M \geq \underline{\lambda}_j \sqrt{2 \ln(p_j-s_j)}(1-\varepsilon)\big) \to 1.$$
If $\lim_{n\to\infty}\frac{\kappa_j}{\underline{\lambda}_j^2\ln(p_j-s_j)} <1$, then $P\big(\max_{M\in O_j}Z_M\geq\sqrt{2\kappa_j}\big) \to 1$ and by \eqref{eq:upperboundrecov} $P(\hat{S}^b=S) \to 0$.

\subsection{Proof of Lemma~\ref{lem:upperboundnoncentral}}
The result follows directly from Lemma~\ref{lem:upperboundnoncentralgen}, with $M=S$.

\subsection{Proof of Proposition~\ref{prop:nectypeIIreg}}

Let $M=S\setminus S^{S}(\kappa)$, since $S^{S}(\kappa)\neq \emptyset$, $M\neq S$. The event $(\hat{S}^b=S)$ implies $\big(\frac{NC(M)}{NC(S)}<1\big)$ then $P(\hat{S}^b=S)\leq P\big(\frac{NC(M)}{NC(S)}<1\big)$. By definition of $NC$, we have
$P\big(\frac{NC(M)}{NC(S)}<1\big) = P(L_{SM}> 2\Delta_{SM})$ where by Lemma~\ref{lemma:s7},  $L_{SM}\sim \chi_{|S\setminus M|}^2(\mu_{SM})$. We  bound next $P(L_{SM}> 2\Delta_{SM})$ with Lemma~\ref{lemma:s1s3}. We first show that $|S\setminus M|=o(2\Delta_{SM})$. and $\mu_{SM}=o(2\Delta_{SM})$. We have:
$$\frac{|S\setminus M|}{2\Delta_{SM}}\leq \bigg[\sum_{j=1}^b\frac{2|S_j \setminus M_j|}{|S\setminus M|}\kappa_j \bigg]^{-1} \leq \Big[ 2\min_{j=1,\ldots,b}\kappa_j \Big]^{-1} $$
where the last inequality follows from $\sum_{j=1}^b\frac{|S_j \setminus M_j|}{|S\setminus M|}=1$. By assumption, for all $j=1,\ldots,b$ $\kappa_j\to +\infty$ and the lef-hand side above goes to 0 as $n$ grows. By Lemma~\ref{lem:upperboundnoncentral}, we also have:
$$\frac{\mu_{SM}}{2\Delta_{SM}}\leq \frac{ \sum_{j=1}^b |S_j\setminus M_j|n\bar{\lambda}\max_{i\in S_j\setminus M_j}{\beta_i^*}^2}{\sum_{j=1}^b |S_j\setminus M_j|2 \kappa_j}.$$
Observe that:
\begin{eqnarray*}
    \sum_{j=1}^b\frac{n\bar{\lambda}\max_{i\in S_j\setminus M_j}{\beta_i^*}^2}{2\kappa_j}
    &\geq &\sum_{j=1,|S_j\setminus M_j|\neq 0}^b\frac{|S_j\setminus M_j|n\bar{\lambda}\max_{i\in S_j\setminus M_j}{\beta_i^*}^2}{|S_j\setminus M_j|2\kappa_j} \\
    &=& \frac{\sum_{j=1,|S_j\setminus M_j|\neq 0 }^b|S_j\setminus M_j|n\bar{\lambda}\max_{i\in S_j\setminus M_j}{\beta_i^*}^2\Big(1+\sum_{l\neq j}\frac{|S_l\setminus M_l|2\kappa_l}{|S_j\setminus M_j|2\kappa_j}\Big)}{\sum_{j=1}^b|S_j\setminus M_j|2\kappa_j} \\
    &\geq & \frac{\sum_{j=1}^b|S_j\setminus M_j|n\bar{\lambda}\max_{i\in S_j\setminus M_j}{\beta_i^*}^2}{\sum_{j=1}^b \big|S\setminus M_j\big|2\kappa_j}
\end{eqnarray*}
where last inequality follows from $\Big(1+\sum_{l\neq j}\frac{|S_l^S|\kappa_l}{|S_j^{S}(\kappa)|\kappa_j}\Big)\geq 1$ for all $j$ such that $|S_j^{S}(\kappa)|\neq 0$, and then
$$\frac{\mu_{SM}}{2\Delta_{SM}}\leq \sum_{j=1}^b\frac{n\bar{\lambda}\max_{i\in S_j\setminus M_j}{\beta_i^*}^2}{2\kappa_j}.$$
For every $j=1,\ldots,b$, $S_j\setminus M_j \subset S_j^{S}(\kappa)$ , then by definition of the $S_j^{S}(\kappa)$, the right-hand 
side above goes to 0 and so does $\frac{\mu_{SM}}{2\Delta_{SM}}$. By Lemma~\ref{lemma:s7}, for any $\phi \in (0,1)$, for $n$ large enough:
\begin{equation}
    P(L_{SM}> 2\Delta_{SM}) \leq e^{-\phi2\Delta_{SM}}
\end{equation} 
where the right hand-side goes to 0 since for all $j=1,\ldots,b$, $\kappa_j \to +\infty$. 

\subsection{Proof of Corollary~\ref{cor:necrecovreg}}
If for some $j\in\{1,\ldots,b\}$ $\lim_{n\to\infty}\frac{\sqrt{n\bar{\lambda}}\beta_{\min,j}^*}{\sqrt{\underline{\lambda}_j\ln(p_j-s_j)}} = 0$ then for that $j$ we have:
$$\lim_{n\to\infty}\frac{\sqrt{n\bar{\lambda}}\beta_{\min,j}^*}{\sqrt{\kappa_j}}\sqrt{\frac{\kappa_j}{\underline{\lambda}_j\ln(p_j-s_j)}} = 0$$
If $\lim_{n\to\infty}\frac{\sqrt{n\bar{\lambda}}\beta_{\min,j}^*}{\sqrt{\kappa_j}}>0$, then $\lim_{n\to\infty}\sqrt{\frac{\kappa_j}{\underline{\lambda}_j\ln(p_j-s_j)}} = 0$ and by Proposition~\ref{prop:nectypeIreg} $P(\hat{S}^b=S)\not\to1$. If $\lim_{n\to\infty}\frac{\sqrt{n\bar{\lambda}}\beta_{\min,j}^*}{\sqrt{\kappa_j}}=0$, then by Proposition~\ref{prop:nectypeIIreg} $P(\hat{S}^b=S)\not\to1$.

\subsection{Proof of Corollary~\ref{cor:selconstimpossiblereg}}

Since $\hat{S}$ is $\hat{S}^b$ with $b=1$, the assumptions of Corollary~\ref{cor:necrecovreg} for $\hat{S}$ are met and $P(\hat{S} = S)\not \to 1$. Since A6 and A7 hold, by Theorem~\ref{theo:suffcondlinearmodel} $P(\hat{S}^b = S)\to 1$.

\subsection{Proof of Theorem~\ref{thm:convtoT}}
By definition of $\T$, for any $M \not\in\T$, there exists a unique $T^M\in\T$ such that $M\cap S^{I}(\kappa)=T^M\cap S^{I}(\kappa)$ and $M\neq T^M$. We next show that $A_{T^M}$ as defined in Lemma~\ref{lem:boundexpnc} satisfies the conditions of the lemma. Observe that $\Delta_{MT^M}$, defined in \eqref{eq:DeltaM}, can be rewritten as $\Delta_{MT^M}=\sum_{j=1}^b(|M_j\setminus T^M_j|-|T^M_j\setminus M_j|)\kappa_j$. By  Lemma~\ref{lem:boundrhogen}, for every $n\in \N$ we have 
\begin{equation}\label{eq:nonoverfittingcond2}
   \gamma \Delta_{MT^M}+\frac{1-\gamma}{2}\mu_{Q_{T^M}M}\;\geq\; \gamma\sum_{j=1}^b|M_j \setminus T^M_j|\kappa_j + \sum_{j=1}^b |T^M_j\setminus M_j |\Big(\tfrac{1-\gamma}{6} n \rho(\bs X) \min_{i\in |T^M_j\setminus M_j|}{\beta^*_i}^{2} - \gamma\kappa_j\Big) 
\end{equation} 
Since $T^M \subseteq S^L(\kappa) \cup S^{I}(\kappa)$ and $M\cap S^{I}(\kappa)=T^M\cap S^{I}(\kappa)$, $T^M\setminus M \subseteq S^L(\kappa)$ and by definition of $S^L(\kappa)$, the rightmost component in \eqref{eq:nonoverfittingcond2} is nonnegative and positive if $|T^M\setminus M|\neq0$, for $n$ large enough. By assumption A6, component $\gamma\sum_{j=1}^b|M_j \setminus T^M_j|\kappa_j$ is nonnegative and positive if $|M\setminus T^M|\neq0$.  $M\neq T^M$ implies necessarily $|M\setminus T^M|\neq0$ or $|T^M\setminus M|\neq0$ and then, for $n$ large enough, $A_{T^M}>0$. If $|M\setminus T^M|=0$, we immediately have $|M\setminus T^M|=o(A_{T^M})$. If $|M\setminus T^M|\neq 0$, since the rightmost component in \eqref{eq:nonoverfittingcond2} is nonnegative, we have:
$$\frac{|M\setminus T^M|}{A_{T^M}}\leq \bigg[ \gamma\sum_{j=1}^b\frac{|M_j \setminus T^M_j|}{|M\setminus T^M|}\kappa_j \bigg]^{-1} \leq \bigg[ \gamma\min_{j=1,\ldots,b}\kappa_j \bigg]^{-1} $$
where the last inequality follows from $\sum_{j=1}^b\frac{|M_j \setminus T^M_j|}{|M\setminus T^M|}=1$. Then, by assumption A6, we have $|M\setminus T^M|=o(A_{T^M})$ when $|M\setminus T^M|\neq 0$ too. 
We have $M\setminus T^M\subseteq V\setminus (S^{I}(\kappa) \cup S^L(\kappa))$. If $M\cap S^{S}(\kappa)=\emptyset$, we immediately get $\mu_{Q_{T^M} T^M}=o(A_{T^M})$ because $\bs\beta^*_{Q_{T^M} \setminus T^M}=0$ and $\mu_{Q_{T^M} T^M}=0$. If $M\cap S^{S}(\kappa)\neq\emptyset$, by Lemma~\ref{lem:upperboundnoncentralgen}, we then have:
$$\frac{\mu_{Q_{T^M} T^M}}{ A_{T^M}} \leq \frac{n\,\bar{\lambda}\,\sum_{j=1}^b|(S_j\cap M_j )\setminus T^M_j|\,\max_{i\in (S_j\cap M_j )\setminus T^M_j}{\beta_{i}^*}^2}{ \gamma\sum_{j=1}^b|M_j \setminus T^M_j|\kappa_j}.$$
For all $j=1,\ldots,b$, $|M_j \setminus T^M_j| \geq |(S_j\cap M_j )\setminus T^M_j|$. Moreover, since  $M\setminus T^M\subseteq V\setminus (S^{I}(\kappa) \cup S^L(\kappa))$ and $T^M\cap S^{I}(\kappa)=M\cap S^{I}(\kappa)$, $(S_j\cap M_j )\setminus T^M_j \subseteq S^{S}_j(\kappa)$ for all $j=1,\ldots,b$. We then get:
$$\frac{\mu_{Q_{T^M} T^M}}{ A_{T^M}} \leq \frac{n\,\bar{\lambda}\,\sum_{j=1}^b|(S_j\cap M_j )\setminus T^M_j|\,\max_{i\in S^{S}_j(\kappa)}{\beta_{i}^*}^2}{ \gamma\sum_{j=1}^b|(S_j\cap M_j ) \setminus T^M_j|\kappa_j}.$$
Observe that:
\begin{eqnarray*}
    \sum_{j=1}^b\frac{n\bar{\lambda}\max_{i\in S^{S}_j(\kappa)}{\beta_i^*}^2}{\gamma\kappa_j}
    &\geq &\sum_{j=1,|(S_j\cap M_j )\setminus T^M_j|\neq 0}^b\frac{|(S_j\cap M_j )\setminus T^M_j|n\bar{\lambda}\max_{i\in S^{S}_j(\kappa)}{\beta_i^*}^2}{|(S_j\cap M_j )\setminus T^M_j|\gamma\kappa_j} \\
    &=& \frac{\sum_{j=1,|(S_j\cap M_j )\setminus T^M_j|\neq 0 }^b|(S_j\cap M_j )\setminus T^M_j|n\bar{\lambda}\max_{i\in S^{S}_j(\kappa)}{\beta_i^*}^2\Big(1+\sum_{l\neq j}\frac{|(S_l\cap M_l)\setminus T^M_l|\gamma\kappa_l}{|(S_j\cap M_j )\setminus T^M_j|\gamma\kappa_j}\Big)}{\sum_{j=1}^b|(S_j\cap M_j )\setminus T^M_j|\gamma\kappa_j} \\
    &\geq & \frac{\sum_{j=1}^b|(S_j\cap M_j )\setminus T^M_j|n\bar{\lambda}\max_{i\in S^{S}_j(\kappa)}{\beta_i^*}^2}{\sum_{j=1}^b \big|(S\cap M_j)\setminus T^M_j\big|\gamma\kappa_j} 
\end{eqnarray*}
where last inequality follows from $\Big(1+\sum_{l\neq j}\frac{|(S_l\cap M_l)\setminus T^M_l|\gamma\kappa_l}{|(S_j\cap M_j )\setminus T^M_j|\gamma\kappa_j}\Big)\geq 1$ for all $j$ and then,
$$\frac{\mu_{Q_{T^M} T^M}}{ A_{T^M}} \leq \sum_{j=1}^b\frac{n\bar{\lambda}\max_{i\in S^{S}_j(\kappa)}{\beta_i^*}^2}{\gamma\kappa_j}.$$
By definition of the $S_j^{S}(\kappa)$, the right-hand side above goes to 0 as $n$ grows and $\mu_{Q_{T^M} T^M}=o(A_{T^M})$ when $M\cap S^{S}(\kappa)\neq\emptyset$ too. 
We conclude that for every $M\in \mathcal M\setminus \T$, by Lemma~\ref{lem:boundexpnc}, for any $\psi\in(0,1)$ and $n$ large enough:
$$
\E(NC(M))\;\leq\;e^{-\psi A_{T^M}}.
$$
Let 
$$A_{T^M}^*\;:=\;\gamma\sum_{j=1}^b|M_j \setminus T^M_j|\kappa_j + \sum_{j=1}^b |T^M_j\setminus M_j |\left(\tfrac{1-\gamma}{6} n \rho(\bs X) \min_{i\in S^{L}_j(\kappa)}{\beta^*_i}^{2} - \gamma\kappa_j\right) $$
By \eqref{eq:nonoverfittingcond} and because $\min_{i\in |T^M_j\setminus M_j|}{\beta^*_i}^{2}\geq \min_{i\in S^{L}_j(\kappa)}{\beta^*_i}^{2}$, we have for $n$ large enough:
\begin{equation}\label{eq:boundexp2}
    \E(NC(M))\;\leq\;e^{-\psi A_{T^M}^*}.
\end{equation}
The definition of $S^{L}_j(\kappa)$ implies there exist $g'_j \to +\infty$ such that:
$$\frac{(1-\gamma) n\rho(\bs X)}{6}\min_{i\in S^{L}_j(\kappa)}{\beta^*_i}^{2}\,- \,\kappa_j \;=\; \ln(s_j) + g'_{j}.$$
Let $\delta \in (0,1)$ and denote $\bar{m}_j=\max\big\{\frac{2\ln(p_j-s_j)}{f_j},\frac{2\ln(s_j)}{g'_{j}}\big\}$. Take $\psi=\max_{j=1,\ldots,b}\frac{\xi + \delta + \bar{m}_j}{1+\bar{m}_j}$ for some $\xi \in(0, 1- \delta)$ then $\psi \in (0,1)$ and we have for every $j=1,\ldots,b$:
$$\delta \leq \psi-(1-\psi)\bar{m}_j.$$ 
By definition of $\gamma$ we have: 
$$
\gamma\kappa_j\;\geq\;\frac12 \Big(1+\frac{\ln(p_j-s_j)}{\kappa_j}\Big)\kappa_j\;=\;\ln(p_j-s_j)+\frac12(\kappa_j-\ln(p_j-s_j)) \;=\; \ln(p_j-s_j)+ \frac12f_j.
$$
From our choice of $\psi,\delta$, we get:
$$\psi\gamma\kappa_j \;\geq\; \psi\big(\ln(p_j-s_j)+ \frac12f_j\big) \;\geq\;  \ln(p_j-s_j)+ \delta\frac12f_j.$$
By A7 and by our choice of $\psi,\delta$,
$$\psi\left(\tfrac{1-\gamma}{6} n \rho(\bs X) \min_{i\in S^{L}_j(\kappa)}{\beta^*_i}^{2} - \gamma\kappa_j\right) \;\geq\; \psi\left(\ln(s_j)+g_j'\right) \;\geq\; \ln(s_j)+\delta \frac{g'_j}{2}.
$$
Using the expression for $A_{T^M}^*$, the bound in \eqref{eq:boundexp2}, and the above inequalities, we get
$$
\E(NC(M))\;\leq\;\exp\left\{- \sum_{j=1}^b |M_j\setminus T^M_j|(\ln(p_j-s_j)+\delta \tfrac{f_j}{2})-\sum_{j=1}^b |T^M_j\setminus M_j|(\ln(s_j)+\delta \tfrac{g'_j}{2})\right\}.
$$

Denote $\mathcal{S}= \sum_{M \in \M\setminus \T} \E\left(NC(M)\right)$ for convenience. Recall that for $M \in \M\setminus \T$ there exists a unique $T^M\in\T$ such that $M\cap S^{I}(\kappa)=T^M\cap S^{I}(\kappa)$.  We can split $\mathcal{S}$ into sums over the models $M$ such that, for all $j$, $|M_j\setminus T^M_j|=u_j$ and $|T^M_j\setminus M_j|=w_j$ with $u_j \in \{0,\ldots,p_j-s_j,\ldots,p_j-s_j+|S_j^{S}(\kappa)|\}$ and $w_j \in \{0,\ldots,|S^{L}_j(\kappa)|\}$, not all simultaneously zero:
$$\mathcal{S} =\sum_{w_1=0}^{|S^L(\kappa)_1|}\cdots\sum_{w_b=0}^{|S^L(\kappa)_b|}\sum_{u_1=0}^{p_1-s_1+|S_1^S|}\cdots\sum_{u_b=0}^{p_b-s_b+|S_b^S|}\sum_{M \neq T^M:\forall \,j |M_j\setminus T^M_j|=u_j, |T^M_j\setminus M_j|=w_j } \E\left(NC(M)\right)$$
The number of models with given collection of $u_j$'s and $w_j$'s is $\prod_{j=1}^b {\binom{p_j-s_j+|S_j^{S}(\kappa)|}{u_j}}{\binom{|S_j^{L}(\kappa)|}{ w_j}}$. We thus have:
\begin{equation*}
    \mathcal{S} \leq \prod_{j=1}^b\left(1+\sum_{u_j=1}^{p_j-s_j+|S_j^{S}(\kappa)|}{\binom{p_j-s_j +|S_j^{S}(\kappa)|}{u_j}}e^{- u_j(\ln(p_j-s_j)+\delta \tfrac{f_j}{2})}\right)\left(1+
    \sum_{w_j=1}^{|S_j^{L}(\kappa)|}{\binom{|S_j^{L}(\kappa)|}{ w_j}}
    e^{-w_j(\ln(s_j)+\delta \tfrac{g'_j}{2})}\right) - 1.
\end{equation*}
where $-1$ stems from discarding the cases $M=T^M$ ($u_j=0$ and $w_j=0$ for all $j$). A standard bound on binomial coefficient for $1\leq k\leq n$ is:
\begin{equation*}
    \binom{n}{k} \leq \left(\frac{n\,e}{k}\right)^k \leq \left(n\,e\right)^k.
\end{equation*}
Then
\begin{equation*}
    \mathcal{S} \;\;\leq\;\; \prod_{j=1}^b\left(1+\sum_{u_j=1}^{p_j-s_j+|S_j^{S}(\kappa)|}e^{-u_j\left(\delta \tfrac{f_j}{2}-\ln\big(1+\frac{|S^{S}_j(\kappa)|}{p_j-s_j}\big)-1\right)}\right)\left(1+
    \sum_{w_j=1}^{|S_j^{L}(\kappa)|}    e^{-w_j\left(\ln\big(\frac{s_j}{|S_j^{L}(\kappa)|}\big)+\delta \tfrac{g'_j}{2}-1\right)}\right) - 1.
\end{equation*}
Denote
$$
d_j\;=\;e^{1+\ln\big(1+\frac{|S^{S}_j(\kappa)|}{p_j-s_j}\big)-\delta \tfrac{f_j}{2}},\qquad g_j\;=\;e^{1-\ln\big(\frac{s_j}{|S_j^{L}(\kappa)|}\big)-\delta \tfrac{g'_j}{2}}.
$$
Since $|S^{S}_j(\kappa)|=O(p_j-s_j)$ and $s_j\geq|S_j^{L}(\kappa)|$, $d_j$ and $g_j$ go to zero as $n$ increases.
For every $j$, we have:
\begin{eqnarray*}
    &1+\sum_{u_j=1}^{p_j-s_j+|S_j^{S}(\kappa)|}e^{-u_j\left(\delta \tfrac{f_j}{2}-1\right)}
    =\frac{1-d_j^{p_j-s_j+|S_j^{S}(\kappa)|+1}}{1-d_j} \\
    &1+\sum_{w_j=1}^{|S_j^{L}(\kappa)|}    e^{-w_j\left(\delta \tfrac{g'_j}{2}-1\right)} =\frac{1-g_j^{|S_j^{L}(\kappa)|+1}}{1-g_j}.
\end{eqnarray*}
Since both expressions converge to $1$ as $n$ grows, we conclude $\mathcal S\to 0$ which proves (i). Finally, (ii) follows from (i) and observing that by Lemma~\ref{lem:L0toL1convergencegen}, 
$P(\hat{S}^{b} \not\in \T ) \,\,\leq \,\, (2^{|S^{I}(\kappa)|}+1) \sum_{M \in \M\setminus \{\T\}} \E\left(NC(M)\right) = (2^{|S^{I}(\kappa)|}+1) \mathcal{S}$.


\section{Proofs of Section~\ref{sec:informed_l0}}

\subsection{Proof of Lemma~\ref{lem:empiricalbayesinterpret}}
Denote $H(\bs \theta)=p(\bs y \mid \bs \theta)$. For every $j=1,\ldots,b$, the partial derivative with respect to $\theta^{(j)}$ is:
    \begin{equation}\label{eq:logderivative}
        \frac{\partial \ln H(\bs \theta)}{\partial \theta^{(j)}}=\frac{\partial H(\bs \theta)}{\partial \theta^{(j)}}\,H(\boldsymbol{\theta})^{-1}
    \end{equation}
    Observe that:
    \begin{equation}\label{eq:partialmarglikelihood}
    H(\bs \theta)=\sum_{M \in \M}p(\bs y \mid M)p(M \mid \bs \theta)
    \quad \text{and}\quad 
    \frac{\partial H(\bs \theta)}{\partial \theta^{(j)}}=\sum_{M\in \M}p(\bs y \mid M)\frac{\partial p(M \mid \bs \theta)}{\partial \theta^{(j)}}
    \end{equation}
    By our choice of prior \eqref{eq:block_modelprior}:
    \begin{equation}\label{eq:partialprobmodel}
    \frac{\partial p(M \mid \bs \theta)}{\partial \theta^{(j)}}= p(M \mid \bs \theta) \bigg(\frac{\sum_{i\in B_j}m_i}{\theta^{(j)}}-\frac{p_j-\sum_{i\in B_j}m_i}{1-\theta^{(j)}} \bigg)        
    \end{equation}
    Replacing \eqref{eq:partialprobmodel} into \eqref{eq:partialmarglikelihood}, we get:
    \begin{equation*}\label{eq:partialmarglikelihood2}
        \frac{\partial H(\bs \theta)}{\partial \theta^{(j)}}=\frac{1}{\theta^{(j)}}\sum_{i \in B_j}\sum_{M\in \M: m_i=1}p(\bs y \mid M)p(M \mid \bs \theta) - \frac{1}{1-\theta^{(j)}}\bigg[p_j H(\bs \theta)-\sum_{i \in B_j}\sum_{M\in \M: m_i=1}p(\bs y \mid M)p(M \mid \bs \theta)
        \bigg]
    \end{equation*}
    Note that:
    $$\frac{\sum_{M\in \M: m_i=1}p(\bs y \mid M)p(M \mid \bs \theta)}{\sum_{M\in \M}p(\bs y \mid M)p(M \mid \bs \theta)} = P(m_i = 1\mid \bs y)$$
    By \eqref{eq:logderivative}, we then get the following expression of the partial derivative for every $j=1,\ldots,b$:
    $$\frac{\partial \ln H(\bs \theta)}{\partial \theta^{(j)}}=\frac{1}{\theta^{(j)}}\sum_{i \in B_j}P(m_i = 1\mid \bs y) - \frac{1}{1-\theta^{(j)}}\bigg[p_j-\sum_{i \in B_j}P(m_i = 1\mid \bs y)
        \bigg]$$
    Setting the partial derivatives to 0 and solving for $\theta^{(j)}$ gives the desired result.

\subsection{Proof of Proposition~\ref{prop:shatconsistence}}
We start by showing the lower bound. For every $j=1,\ldots,b$, we have the decomposition:
\begin{gather}\label{eq:decompshat}
    \frac{\widehat{s}_j}{p_j}
   =A_j\,+\, C_j \\
       \text{where} \quad A_j=\frac{1}{p_j}\sum_{i \in B_j}\sum_{M \in \T|i\in M}NC(M) \quad\text{and}\quad  C_j=\frac{1}{p_j}\sum_{i \in B_j}\sum_{M \in \M\setminus \T|i\in M}NC(M). \nonumber
\end{gather}
For every $j=1,\ldots,b$, $A_j$ can be furthered decomposed:
\begin{eqnarray*}
  A_j
   &=& \sum_{M \in \T}NC(M)\sum_{i \in B_j}\frac{I(i\in M_j)}{p_j}\\
   &=& \sum_{M \in \T}NC(M)\sum_{i \in S^{L}_j(\kappa)}\frac{I(i\in M_j)}{p_j}+\sum_{M \in \T}NC(M)\sum_{i \in B_j\setminus S^{L}_j(\kappa)}\frac{I(i\in M_j)}{p_j}\\
   &=& \frac{|S^{L}_j(\kappa)|}{p_j}\sum_{M \in \T}NC(M)+\sum_{M \in \T}NC(M)\sum_{i \in B_j\setminus S^{L}_j(\kappa)}\frac{I(i\in M_j)}{p_j}.
\end{eqnarray*}
where the last equality follows from $I(i\in M_j)=1$ for all $i\in S^L(\kappa)$ since $M\in\T$. The righmost term above and $C_j$ are nonnegative, we then get:
$$\E\Big(\frac{\widehat{s_j}}{p_j}\Big)\geq\frac{|S^{L}_j(\kappa)|}{p_j}\sum_{M \in \T}\E(NC(M))$$
and by Theorem~\ref{thm:convtoT} (i), $\lim_{n\to\infty}\E\Big(\frac{\widehat{s_j}}{p_j}\Big)\geq \frac{|S^{L}_j(\kappa)|}{p_j}$ for every $j=1,\ldots,b$.

We continue with the upper bound. Since that for all $M\in\T$, $I(i\in M_j)=0$ for all $i\in B_j\setminus(S^{L}_j(\kappa)\cup S^{I}(\kappa)_j)$, $A_j$ can be further written:
$$
  A_j
   = \frac{|S^{L}_j(\kappa)|}{p_j}\sum_{M \in \T}\E(NC(M))+\sum_{M \in \T}\E(NC(M))\sum_{i \in S^{I}(\kappa)_j}\frac{I(i\in M_j)}{p_j}.
$$
Using that $\sum_{i \in S^{I}(\kappa)_j}\frac{I(i\in M_j)}{p_j}\leq \frac{|S^{I}(\kappa)_j|}{p_j}$ for all $M$, we have:
\begin{eqnarray}\label{eq:uuuuu}
    \frac{\widehat{s}_j}{p_j}&\leq \frac{|S^{L}_j(\kappa)|+|S^{I}(\kappa)_j|}{p_j}\sum_{M \in \T}NC(M) + C_j.
\end{eqnarray}
A bound on $C_j$ for every $j=1,\ldots,b$ is:
\begin{equation}\label{eq:uuuuu2}
    C_j = \sum_{M \in \M\setminus \T}NC(M)\sum_{i \in B_j}\frac{I(i\in M_j)}{p_j}\leq \sum_{M \in \M\setminus \T}NC(M)
\end{equation}
where the inequality follows from $\sum_{i \in B_j}\frac{I(i\in M_j)}{p_j}\leq 1$ for all $M$.
Taking expectation in \eqref{eq:uuuuu} and \eqref{eq:uuuuu2} gives:
\begin{eqnarray*}
    \E\bigg(\frac{\widehat{s_j}}{p_j}\bigg) 
    &\leq \frac{|S^{L}_j(\kappa)|+|S^{I}(\kappa)_j|}{p_j}\sum_{M \in \T}\E(NC(M)) + \sum_{M \in \M\setminus \T}\E(NC(M)).
\end{eqnarray*}
 By Theorem~\ref{thm:convtoT} (i), $\sum_{M \in \T}\E(NC(M))\to 1$ and $\sum_{M \in \M\setminus\T}\E(NC(M))\to 0$, then $ \lim_{n\to+\infty}\E\big(\frac{\widehat{s_j}}{p_j}\big) \leq  \frac{|S^{L}_j(\kappa)|+|S^{I}(\kappa)_j|}{p_j}= \frac{s_j-|S^{S}_j(\kappa)|}{p_j}$ for every $j=1,\ldots,b$, which proves the upper bound.


\subsection{Proof of Theorem \ref{theo:empbayesselconsist}}
Denote for any $M\in\M$, $NC^{\circ}(M)$, the normalized criterion value for model $M$ under step 1 penalty $\kappa^{\circ}$. For this choice of penalty and every $j=1,\ldots,b$, we have:
\begin{equation*}
    \frac{\widehat{s_j}}{p_j}
   = \frac{1}{p_j}\sum_{M \in \M}NC^{\circ}(M)\sum_{i \in B_j}I(i\in M) =\sum_{M \in \M | M \neq S}\frac{|M_j|}{p_j}NC^{\circ}(M) +\frac{s_j}{p_j}NC^{\circ}(S) .
\end{equation*}
Using that $NC^{\circ}(S)=1-\sum_{M \in \M | M \neq S}NC(M)$, we get
\begin{equation*}
   \frac{\widehat{s}_j}{p_j}
   =\sum_{M \in \M | M \neq S}\frac{|M_j|-s_j}{p_j}NC^{\circ}(M) \,\,+\,\, \frac{s_j}{p_j}.
\end{equation*}
Splitting the sum in the right-hand side above between models $M$ that are overfitting in $B_j$ in the sense $|M_j|> s_j$ and those that are underfitting in $B_j$, in the sense $|M_j|<s_j$, we can write:
\begin{equation}\label{eq:rewriteshat}
  \frac{\widehat{s}_j}{p_j}  = \frac{s_j}{p_j}\,\,+\,\, O_j^{\circ}-\,\,U_j^{\circ}
\end{equation}
where:
$$
    O_j^{\circ} = \sum_{M \in \M | |M_j| > s_j}\frac{|M_j|-s_j}{p_j}NC^{\circ}(M)\quad\text{and}\quad
    U_j^{\circ} = \sum_{M \in \M | |M_j| < s_j}\frac{s_j-|M_j|}{p_j}NC^{\circ}(M).
$$
By \eqref{eq:rewriteshat}, we have the following decomposition of step 2 penalties :
\begin{eqnarray}\label{eq:decompstep2pen}
    \kappa^{EB}_j
    &=& \ln(p_j-s_j) + \ln\Big(\frac{\sqrt{n}}{s_j}\Big) + \ln\Big(\frac{p_j-\widehat{s}_j}{p_j-s_j}\Big) + \ln\Big(\frac{s_j}{\widehat{s}_j}\Big) \nonumber\\
    &=& \ln(p_j-s_j)+\ln\Big(\frac{\sqrt{n}}{s_j}\Big)+\ln\Big(1-\frac{p_j(O^{\circ}_j-U^{\circ}_j)}{p_j-s_j}\Big)+\ln\Big(\frac{s_j}{\widehat{s}_j}\Big).
\end{eqnarray}
We first show the $\kappa^{EB}_j$'s satisfy assumption~A6 with probability going to 1. Since $U_j^{\circ}$ is nonnegative, a lower bound on $\kappa^{EB}_j$ for every $j=1,\ldots,b$ is:
\begin{equation}\label{eq:decompstep2pen2}
    \kappa^{EB}_j
    \geq  \ln(p_j-s_j)+\ln\Big(\frac{\sqrt{n}}{s_j}\Big)+\ln\Big(1-\frac{p_j O^{\circ}_j}{p_j-s_j}\Big)+\ln\Big(\frac{s_j}{\widehat{s}_j}\Big).
\end{equation}
We have :
$$\frac{p_j O^{\circ}_j}{p_j-s_j} =  \sum_{M \in \M | |M_j| > s_j}\frac{|M_j|-s_j}{p_j-s_j}NC^{\circ}(M)\leq \sum_{M \in \M | |M_j| > s_j}NC^{\circ}(M).$$ 
If $M$ is such that $|M_j|>s_j$, then $M\not\in\T$ and by Theorem~\ref{thm:convtoT}, $O^{\circ}_j$ and $\ln\Big(1-\frac{p_j O^{\circ}_j}{p_j-s_j}\Big)$ vanish in probability. By assumption $\ln(\sqrt{n}s_j^{-1})\to+\infty$. Then, to show A6 holds it is enough to show that with probability going to 1 $\ln\Big(\frac{s_j}{\widehat{s}_j}\Big)$ is nonnegative. By \eqref{eq:uuuuu} and \eqref{eq:uuuuu2} in the proof of Proposition~\ref{prop:shatconsistence}, we have:
\begin{equation*}
    \frac{\widehat{s}_j}{p_j}\leq \frac{|S^{L}_j(\kappa)|+|S^{I}(\kappa)_j|}{p_j}\sum_{M \in \T}NC(M) + \sum_{M \in \M\setminus \T}NC(M).
\end{equation*}
By Theorem~\ref{thm:convtoT} (i), $\sum_{M \in \T}NC(M)$ and $\sum_{M \in \M\setminus\T}NC(M)$ converge in probability to 1 and 0 respectively. We then have with probability going to 1:
\begin{equation*}
    \frac{\widehat{s}_j}{p_j}\leq \frac{|S^{L}_j(\kappa)|+|S^{I}(\kappa)_j|}{p_j} \quad\text{and}\quad\ln\Big(\frac{s_j}{\widehat{s}_j}\Big)\geq \ln\bigg(\frac{s_j}{|S^{L}_j(\kappa)|+|S^{I}(\kappa)_j|} \bigg) \geq 0.
\end{equation*}
We conclude that we probability going to 1:
\begin{equation}\label{eq:eeeeeeeh}
    \kappa^{EB}_j
    \geq  \ln(p_j-s_j)+\ln\Big(\frac{\sqrt{n}}{s_j}\Big)
\end{equation}
and the $\kappa^{EB}_j$ satisfy assumption~A6.\medskip

We go on with showing that assumption~A9 implies A7 with probability going to 1. Observe:
$$\ln\Big(\frac{\widehat{s}_j}{s_j}\Big)=\ln\Big(1+\frac{p_j(O^{\circ}_j-U^{\circ}_j)}{s_j}\Big) .$$
Since $O_j^{\circ}$ is nonnegative, from \eqref{eq:decompstep2pen}, for every $j=1,\ldots,b$, an upper bound on $\kappa^{EB}_j$ is:
\begin{equation}\label{eq:gggggggg}
    \kappa^{EB}_j \leq \ln\big(p_j-s_j\big)+\ln\Big(\frac{\sqrt{n}}{s_j}\Big)+\ln\bigg(\frac{1+\frac{p_j}{p_j-s_j}U^{\circ}_j}{1-\frac{p_j}{s_j}U^{\circ}_j}\bigg)
\end{equation}
We split the sum in $U^{\circ}_j$ between models in $\T$ and those out of $\T$. Using that if $M\in \T$, $s_j-|M_j|\leq s_j-|S^{L}_j(\kappa^{\circ})|$ and if $M\not\in \T$, $s_j-|M_j|\leq s_j$, we have:
\begin{eqnarray*}
    U^{\circ}_j \leq \frac{s_j-|S^{L}_j(\kappa^{\circ})|}{p_j}\sum_{M \in \T | |M_j| < s_j}NC^{\circ}(M)+  \frac{s_j}{p_j}\sum_{M \in \M \setminus \T| |M_j| < s_j}NC^{\circ}(M)
\end{eqnarray*}
By Theorem~\ref{thm:convtoT}, $\sum_{M \in  \T| |M_j| < s_j}NC^{\circ}(M)$ and $\sum_{M \in \M \setminus \T| |M_j| < s_j}NC^{\circ}(M)$ converge in probability to 1 and 0 respectively. We then get than with probability going to 1:
\begin{eqnarray*}
    \frac{p_j}{p_j-s_j}U^{\circ}_j \leq \frac{s_j-|S^{L}_j(\kappa^{\circ})|}{p_j-s_j}\\
    \frac{p_j}{s_j}U^{\circ}_j \leq  \frac{s_j-|S^{L}_j(\kappa^{\circ})|}{s_j}
\end{eqnarray*}
By replacing $\frac{p_j}{p_j-s_j}U^{\circ}_j$ and $\frac{p_j}{s_j}U^{\circ}_j$ in \eqref{eq:gggggggg} by the bounds above, we get that  with probability going to 1:
\begin{equation}\label{eq:finalupperboundkappaj}
  - \,\sqrt{\kappa^{EB}_j} \geq  \,- \,\sqrt{\ln\big(p_j/|S^{L}_j(\kappa^{\circ})|-1\big)+\frac{1}{2}\ln\Big(n\Big) }
\end{equation}
By \eqref{eq:eeeeeeeh}, with probability going to 1:
\begin{equation}\label{eq:rrrrrrrrr}
    1-\frac12\Big(1+\max_{j=1, \ldots,b}\frac{\ln(p_j-s_j)}{\kappa^{EB}_j}\Big) \geq 1-\psi
\end{equation}
By \eqref{eq:finalupperboundkappaj} and \eqref{eq:rrrrrrrrr}, under assumption~A9, assumption~A7 holds with probability going to 1. Since A6 and A7 hold with probability going to 1, by Theorem \ref{theo:suffcondlinearmodel}, $P(\hat{S}^{EB,b}\neq S) \to 0$.

\subsection{Proof of Theorem \ref{theo:informedl0consist}} We have that $\kappa_j^{A}=\kappa_j^{EB}+\ln(\widehat s_j)$ and by \eqref{eq:decompstep2pen2} in the proof of Theorem \ref{theo:empbayesselconsist}:
\begin{equation*}
    \kappa_j^{A}
    \geq  \ln(p_j-s_j)+\ln(\sqrt{n})+\ln\Big(1-\frac{p_j O^{\circ}_j}{p_j-s_j}\Big).
\end{equation*}
From the proof of Theorem ~\ref{theo:empbayesselconsist}, by Theorem~\ref{thm:convtoT}, $\ln\Big(1-\frac{p_j O^{\circ}_j}{p_j-s_j}\Big)$ vanishes in probability. With probability going to 1, we then have:
\begin{equation}\label{eq:finallowerboundkappaj2}\kappa^{A}_j 
    \geq  \ln(p_j-s_j)+\ln(\sqrt{n})
\end{equation}
and the $\kappa^{A}_j$ satisfy assumption~A6.

Using that $\kappa_j^{A}=\kappa_j^{EB}+\ln(\widehat s_j)$ and $O_j^{\circ}$ in \eqref{eq:decompstep2pen} is nonnegative, for every $j=1,\ldots,b$, an upper bound on $\kappa^{A}_j$ is:
\begin{equation*}
    \kappa^{A}_j \leq \ln\big(p_j-s_j\big)+\ln(\sqrt{n})+\ln\bigg(1+\frac{p_j}{p_j-s_j}U^{\circ}_j\bigg)
\end{equation*}
As in the proof of Theorem ~\ref{theo:empbayesselconsist}, by Theorem~\ref{thm:convtoT}, with probability going to 1:
\begin{equation}\label{eq:finalupperboundkappaj2}
    \frac{p_j}{p_j-s_j}U^{\circ}_j \leq \frac{s_j-|S^{L}_j(\kappa^{\circ})|}{p_j-s_j}\quad\text{and}\quad -\sqrt{\kappa_j} \geq -\sqrt{\ln\big(p_j-|S^{L}_j(\kappa^{\circ})|\big)+\frac{1}{2}\ln(n)}
\end{equation}
Moreover by \eqref{eq:finallowerboundkappaj2}, we have with probability going to 1:
\begin{equation}\label{eq:rrrrrrrrr2}
    1-\frac12\Big(1+\max_{j=1, \ldots,b}\frac{\ln(p_j-s_j)}{\kappa_j}\Big) \geq 1-\xi
\end{equation}
By \eqref{eq:finalupperboundkappaj2} and \eqref{eq:rrrrrrrrr2}  under assumption~A10, assumption~A7 holds with probability going to 1. Since A6 and A7 hold with probability going to 1, by Theorem \ref{theo:suffcondlinearmodel}, $P(\hat{S}^{A,b}\neq S) \to 0$.

\section{Proofs of auxiliary results}\label{sec:proofauxi}

\subsection{Proof of Lemma~\ref{lem:l1isthreshold}}
    Denote $\hat{\bs \beta}$, the minimizer of the problem \eqref{eq:genl1}. The MLE under the full model is $\tilde{\bs\beta}=\tfrac{1}{\sqrt{n}} \bs y$ and the optimized function in \eqref{eq:genl1} can be then rewritten as
	$$
	\tfrac{1}{2}\bs y^{\top}\bs y +\tfrac{n}{2}\sum_{j=1}^b\sum_{i\in B_j}\left(\beta_i^2-2\tilde\beta_i\beta_i+2\tfrac{\lambda_j}{n}|\beta_i|\right)
	$$
	which we optimize with respect to each $\beta_i$ separately. For any $i\in B_j$ we have $\hat\beta_i=0$ if and only if $|\tilde \beta_i|\leq \frac{\lambda_j}{n}$. Similarly  for \eqref{eq:genl1ada}, denoting $\hat{\bs \beta}$ the minimizer of the problem, for any $i\in B_j$ we have $\hat\beta_i=0$ if and only if $|\tilde \beta_i|\leq \frac{\lambda_j}{|\beta^{\circ}_i|n}$. If $\bs \beta^{\circ}= \tilde{\bs\beta}$, then for any $i\in B_j$ $\hat{\beta}_i = 0$ if and only if $|\tilde \beta_i|\leq \sqrt{\tfrac{\lambda_j}{n}}$. If $\bs\beta^\circ$ is a LASSO estimate with penalization $\lambda^\circ$, then for any $i\in B_j$ 
 $\hat{\beta}_i = 0$ if and only if $\tilde{\beta}^2_i + \operatorname{sign}(\lambda^{\circ}/n-\tilde{\beta}_i) \frac{\lambda^\circ}{n}\tilde{\beta}_i-\frac{\lambda_j}{n} \leq 0$. Equivalently, for any $i\in B_j$ 
 $\hat{\beta}_i = 0$ if and only if $|\tilde{\beta}_ i| \leq  \frac{1}{2}\left(\lambda^\circ/n+\sqrt{(\lambda^\circ/n)^2+4 \lambda_j/n}\right)$.

\subsection{Proof of Lemma~\ref{lemma:3tailbounds}}

	(i) \,\, By union bound and the identical distribution of the $y_i$'s,
\begin{center}
    $P\Big(\underset{i\in\{1,\ldots,p\}}{\max}|y_i|> a \Big)     \leq \sum_{i=1}^{p}{ P\left(|\sqrt{n}y_i|>\sqrt{n}a \right)}
    = p\; P\left(|z|>\sqrt{n}a\right).$
\end{center}
    where $z\sim N(0,1)$. 
    By symmetry and the standard tail bound for standard normal, $P(z \geq \delta) \leq \frac{1}{\sqrt{2\pi}} \frac{1}{\delta} e^{-\delta^{2} / 2}$:
\begin{center}
  $P\Big(\underset{i\in\{1,\ldots,p\}}{\max}|y_i|> a \Big) 
     \leq  p\; 2 \, P\left(z>\sqrt{n}a\right) \leq \frac{1}{\sqrt{\pi}}\frac{\sqrt{2}}{\sqrt{n}a}p e^{-\frac{n}{2} a^{2}}.$
\end{center}
Since $a \geq \sqrt{\frac{2\ln(p)}{n}}$, $\frac{\sqrt{2}}{\sqrt{n}a}\leq \frac{1}{\sqrt{\ln(p)}}$. Taking $a^2=a^2-\frac{2\ln(p)}{n}+\frac{2\ln(p)}{n}$, we get \begin{center}
$ P\Big(\underset{i\in\{1,\ldots,p\}}{\max}|y_i|> a \Big) \leq \frac{1}{\sqrt{\pi\ln(p)}} e^{-\frac{n}{2} \left(a^{2}-\frac{2\ln(p)}{n}\right)}.$ \\\end{center}
(ii) $ \underset{i \in \{1,\ldots,s\}}{\min}|\mu_i|- \underset{i \in \{1,\ldots,s\}}{\max}|y_i-\mu_i| > a$ implies $\forall i \; \underset{i \in \{1,\ldots,s\}}{\min}|\mu_i| - |y_i-\mu_i| > a $. Since $\forall i \; |\mu_i|< |\mu_i-y_i|+|y_i|$ we have that $ \forall i \; \min_{i \in \{1,\ldots,s\}}|\mu_i| - |y_i-\mu_i| > a$ implies $\forall i \; |y_i| > a $ and $ \underset{i \in \{1,\ldots,s\}}{\min}|y_i| > a$. By the properties of the probability $P$: 
\begin{center}$ P\Big(\underset{i \in \{1,\ldots,s\}}{\min}|\mu_i|-\underset{i \in \{1,\ldots,s\}}{\max}|y_i-\mu_i| > a\Big) \leq P\Big(\underset{i \in \{1,\ldots,s\}}{\min}|y_i| > a\Big) $. \end{center}
(iii)\,\, We start by bounding the probability $P(|y_i|)>a$ for any $i$. We have that:
$$P(|y_i|>a) = P(y_i>a) + P(y_i<-a) = P(z>\frac{a-\mu}{\sigma}) + P(z<-\frac{a+\mu}{\sigma}) \quad \text{where }z\sim N(0,1).$$
If $\mu>0$, then $P(|y_i|>a) =  P(z>\frac{a-|\mu|}{\sigma}) + P(z<-\frac{a+|\mu|}{\sigma})$. If $\mu<0$, by symmetry, $P(z>\frac{a-\mu}{\sigma})=P(z>\frac{a+|\mu|}{\sigma})=P(z<-\frac{a+|\mu|}{\sigma})$ and $P(z<-\frac{a+\mu}{\sigma})=P(z<\frac{|\mu|-a}{\sigma})=P(z>\frac{a-|\mu|}{\sigma})$. Then for any $\mu$, $P(|y_i|>a)  = P(z>\frac{a-|\mu|}{\sigma}) + P(z<-\frac{a+|\mu|}{\sigma}) $ and using the symmetry of the standard Gaussian for any $a<|\mu|$: 
\begin{equation}\label{eq:random}
    P(|y_i|>a) = 1+P(0<z<\frac{|\mu|-a}{\sigma})-P(0<z<\frac{a+|\mu|}{\sigma}).
\end{equation}
From \cite{chu}, for any $\delta>0$, $\frac{1}{2}\left(1-e^{-\delta^2/2}\right)^{\frac{1}{2}}\leq P(0 <z < \delta)\leq \frac{1}{2}\left(1-e^{-2\delta^2/\pi}\right)^{\frac{1}{2}}$. Thus, for any $\gamma,\delta>0$:
$$ P(0<z \leq \gamma)-P(0<z \leq \delta)\leq
\frac{1}{2}\left(1-e^{-2\gamma^2/\pi}\right)^{\frac{1}{2}}-\frac{1}{2}\left(1-e^{-\delta^2/2}\right)^{\frac{1}{2}}.$$
Applying the above to \eqref{eq:random}, we get:
$$P(|y_i|>a) \leq 
1+ \frac{1}{2}\left(\left(1-e^{-2(|\mu|-a)^2/\pi\sigma^2}\right)^{\frac{1}{2}}-\left(1-e^{-(a+|\mu|)^2/2\sigma^2}\right)^{\frac{1}{2}}\right).\\$$

Using that $x^{\frac{1}{2}}-y^{\frac{1}{2}}=\frac{x-y}{x^{\frac{1}{2}}+y^{\frac{1}{2}}}$ for $x,y>0$, we get:
$$P(|y_i|>a) \leq 1- \frac{e^{-2(|\mu|-a)^2/\pi\sigma^2}-e^{-(a+|\mu|)^2/2\sigma^2}}{2\left(\left(1-e^{-2(|\mu|-a)^2/\pi\sigma^2}\right)^{\frac{1}{2}}+\left(1-e^{-(a+|\mu|)^2/2\sigma^2}\right)^{\frac{1}{2}}\right)}.$$

Finally, $P(\min_{i\in\{1,\ldots,s\}}|y_i|>a)=\exp\left(s\ln{P(|y_i|>a)}\right)$, using that $\ln\left(1+ x \right)< x$ for $x>-1$ we get the desired inequality.

\subsection{Proof of Lemma~\ref{lem:L0toL1convergencegen}}
Suppose $NC(M_1)+\ldots+NC(M_k)>\frac{k}{k+1}$ then for any $M\neq \{M_1,\ldots,M_k\}$, $NC(M)<\frac{1}{k+1}$. In addition, if $NC(M_1)+\ldots+NC(M_k)>\frac{k}{k+1}$ then necessarily $\max_{l=1,\ldots,k}NC(M_l)> \frac{1}{k+1}$ and $\hat{S}^{b}\in \{M_1,\ldots,M_k\}$. Consequently, 
$$
P(\hat{S}^b\not\in \{M_1,\ldots,M_k\})\; \leq\; P\Big(NC(M_1)+\ldots+NC(M_k)\leq\frac{k}{k+1}\Big).$$
Moreover, we have:
$$P\Big(NC(M_1)+\ldots+NC(M_k)\leq\frac{k}{k+1}\Big)\;=\;P\Big(\sum_{M \in \M\setminus \{M_1,\ldots,M_k\}} NC(M)\geq \frac{1}{k+1}\Big).$$
The result follows from the the Markov's inequality applied to the right-hand side above.

\subsection{Proof of Lemma~\ref{lem:boundrhogen}} 
The non-centrality parameter $\mu_{Q_TM}$, as defined in \eqref{eq:muQM}, satisfies
\begin{eqnarray*}
    \mu_{Q_TM}&=&\|\big(I_n-P_M\big) \bs X_{Q_T\setminus M} \bs \beta^*_{Q_T\setminus M}\|^2\\
    &=&\|\big(I_n-P_M\big) \bs X_{T\setminus M} \bs \beta^*_{T\setminus M}\|^2\\
    &= & n {\beta^*_{T\setminus M}}^\top \left(\tfrac1n \bs X_{T\setminus M}^\top(I_n-P_M)\bs X_{T\setminus M}\right)\beta^*_{T\setminus M}\\
    &\geq & n\lambda_{\min}\left(\tfrac1n \bs X_{T\setminus M}^\top(I_n-P_M)\bs X_{T\setminus M}\right)\|\beta^*_{T\setminus M}\|^2,
\end{eqnarray*}
where the second equality follows from observing that $Q_T\setminus M=T\setminus M$. Since $T$ is a susbet of $S$, by reordering columns, $X_{S\setminus M}=[X_{T\setminus M},X_{S\setminus (M\cup T)}]$ and $\tfrac1n \bs X_{T\setminus M}^\top(I_n-P_M)\bs X_{T\setminus M}$ is a diagonal block of $\tfrac1n \bs X_{S\setminus M}^\top(I_n-P_M)\bs X_{S\setminus M}$. By interlacing, $\lambda_{\min}\left(\tfrac1n \bs X_{T\setminus M}^\top(I_n-P_M)\bs X_{T\setminus M}\right) \geq \lambda_{\min}\left(\tfrac1n \bs X_{S\setminus M}^\top(I_n-P_M)\bs X_{S\setminus M}\right)$. Finally observing that 
$\lambda_{\min}\left(\tfrac1n \bs X_{T\setminus M}^\top(I_n-P_M)\bs X_{T\setminus M}\right) \geq \rho(\bs X)$ and $\|\beta^*_{T\setminus M}\|^2\geq \sum_{j=1}^b|T_j\setminus M_j|\,\min_{i\in T_j\setminus M_j}{\beta_{i}^*}^2$, gives the desired result.

\subsection{Proof of Lemma~\ref{lemma:s1s3}}
By \cite{birge}, Lemma 8.1 we have for any $x>0$
$$P(W> (\nu+\mu)+2 \sqrt{(\nu+2 \mu) x}+2 x) \leq e^{-x}.$$
The function $f:x\mapsto (\nu+\mu)+2 \sqrt{(\nu+2 \mu) x}+2 x$ is  one-to-one between $\R^+$ and $( \nu+
\mu,+\infty)$. 
We have for any $w>\mu+\nu$:
$$P(W > w) \leq e^{-f^{-1}(w)} \:=\; e^{-\big(\frac{w+\mu}{2}  - \sqrt{2w(2 \mu+\nu)-2 \mu \nu - \nu^2 }\big) } .$$
Observe that:
$$\frac{w+\mu}{2}  - \sqrt{2w( 2 \mu+ \nu)-2 \mu \nu - \nu^2 }= \frac{w}{2}\bigg(1+\frac{\mu}{w} -\sqrt{\frac{8(2 \mu+\nu)}{w}\bigg(1-\frac{\nu\mu}{2w\mu+w\nu}-\frac{\nu^2}{2(2w\mu+w\nu)} \bigg)}\bigg).$$
 If $\nu=o(w)$ and $\mu=o(w)$, we then have $\frac{w+\mu}{2}  - \sqrt{2w( 2 \mu+ \nu)-2 \mu \nu - \nu^2 }= \frac{w}{2}(1+o(1))$ and for any $\phi\in(0,1)$ and $n$ large enough:
$$P(W > w) \leq  e^{-\phi\frac{w}{2}}$$

\subsection{Proofs of Lemma~\ref{lemma:s2}}
The result follows directly from \cite{rossell2022concentration}, Lemma S2.

\subsection{Proof of Lemma~\ref{lemma:s18}}
For any $u\in [\underline{u},\bar{u}]$, we have $2\ln \left(\frac{g}{1 / u-1}\right) \geq 2\ln \left(\frac{g}{1 / \underline{u}-1}\right)$. Since $\underline{u}\geq (1+g^{\phi}e^{-(\nu+\mu)/2})^{-1}$,  we also have for any $u\in (\underline{u},\bar{u})$:
$$2\ln \left(\frac{g}{1 / u-1}\right) \geq 2(1-\phi)\ln(g)+\nu+\mu .$$
By assumption $\nu=o\left(\ln (g)\right)$ and $\mu=\left(\ln (g)\right)$, then for any $u\in (\underline{u},\bar{u})$, $\nu=o\left(2\ln \left(\frac{g}{1 / u-1}\right)\right)$ and $\mu=o\left(2\ln \left(\frac{g}{1 / u-1}\right)\right)$. By Lemma~\ref{lemma:s1s3}, for any $\phi\in(0,1)$ and $n$ large enough:
\begin{equation}\label{eq:randomm}
    \int_{\underline{u}}^{\bar{u}} P\left(W>2\ln \left(\frac{g}{1 / u-1}\right)\right) d u< \frac{1}{g^{\phi}} \int_{\underline{u}}^{\bar{u}}(1 / u-1)^{\phi} d u .
\end{equation}
Applying the change of variables $v=1 / u-1$ to the integral above gives
$$
\int_{\underline{u}}^{\bar{u}}(1 / u-1)^{\phi} d u=\int_{1 / \bar{u}-1}^{1 / \underline{u}-1} \frac{v^{\phi}}{(v+1)^2} d v.
$$
Note that $(1+x)^r \leq 1+r x$  for $x \geq-1$ and $r\in[0,1]$ and that $v^{\phi}=(v-1 + 1)^{\phi}$ with $v-1>-1$ since $v<1 / \bar{u}-1$ and $\bar{u}<1$. Then $v^{\phi}=(v-1 + 1)^{\phi}\leq 1 + \phi (v-1)\leq 1 + \phi(v+1)$. As a result:
\begin{equation}\label{eq:random2}
\int_{\underline{u}}^{\bar{u}}(1 / u-1)^{\phi} d u<\int_{1 / \bar{u}-1}^{1 / \underline{u}-1} \frac{1}{(v+1)^2} +\frac{\phi}{v+1} d v = \bar{u}-\underline{u}+\phi\ln\Big(\frac{\bar{u}}{\underline{u}}\Big)
\end{equation}
The results follows inputing the bound from \eqref{eq:random2} in \eqref{eq:randomm} and using that $\phi<1$.


\subsection{Proof of Lemma~\ref{lemma:s20}}
Since a probability is bounded by $1$, for any $a\in(0,1)$.
\begin{equation}\label{eq:basebound}
    \int_0^1 P\left(W>2 \ln \left(\frac{g}{1 / u-1}\right)\right) d u \;\leq\; 2 a+\int_{a}^{1-a} P\left(W>2\ln\left(\frac{g}{1 / u-1}\right)\right) d u.
\end{equation}

Take $a=\left(1+g^{\phi}e^{-(\nu+\mu)/ 2}\right)^{-1}$ for some $\phi\in(\alpha,1)$. By Lemma~\ref{lemma:s18} with $\underline{u}=a$ and $\bar{u}=1-a$:
$$
\int_0^1 P\left(W>2 \ln \left(\frac{g}{1 / u-1}\right)\right) d u \;\leq\; \frac{2}{1+g^{\phi}e^{-(\nu+\mu)/2}}+ 
\frac{1-2a+\ln(g^{\phi}e^{-(\nu+\mu)/2})}{g^{\phi}} $$
Since $1+g^{\phi}e^{-(\nu+\mu)/ 2}>g^{\phi}e^{-(\nu+\mu)/ 2}$ and $1-2a-\frac{\nu+\mu}{2}\leq \ln(g^\phi)$ for $n$ large enough, we have:
$$
\int_0^1 P\left(W>2 \ln \left(\frac{g}{1 / u-1}\right)\right) d u \;\leq\; g^{-\phi}\big(2e^{(\nu+\mu)/ 2}+2\ln(g^{\phi})\big) $$
We have:
\begin{equation}\label{eq:random24}
    \frac{g^{-\phi}2e^{(\nu+\mu)/ 2}}{g^{-\alpha}} \;=\;e^{-(\phi-\alpha)\ln(g)\big(1-\frac{\nu+\mu}{2(\phi-\alpha)\ln(g)}- \frac{\ln(2)}{(\phi-\alpha)\ln(g)}\big)}
\end{equation}
and 
\begin{equation}\label{eq:random25}
    \frac{g^{-\phi}2\ln(g^{\phi})}{g^{-\alpha}} \;=\;e^{-(\phi-\alpha)\ln(g)\big(1-\frac{\ln(\phi\ln(g))}{(\phi-\alpha)\ln(g)}- \frac{\ln(2)}{(\phi-\alpha)\ln(g)}\big)}
\end{equation}
Since $\alpha<\phi$, $g$ is increasing and $\nu=o(\ln(g))$ and $\mu=o(\ln(g))$, both expressions in \eqref{eq:random24} and \eqref{eq:random25} vanish as $n$ grows and for $n$ large enough:
$$
\int_0^1 P\left(W>2 \ln \left(\frac{g}{1 / u-1}\right)\right) d u \;\leq\; g^{-\alpha}$$

\subsection{Proof of Lemma~\ref{lem:boundexpnc}}
Since $M\in \mathcal M\setminus \{T\}$, by Lemma~\ref{lem:L0toL1convergence} (ii),  $NC(M)\leq (1+e^{C(T)-C(M)})^{-1}$.  By the standard formula for expectation of a nonnegative random variable in terms of its CDF, we get:
\begin{eqnarray*}
  \E\left(NC(M)\right)  &\leq& \int_0^1 P\left((1+e^{C(T)-C(M)})^{-1}\geq u\right){\rm d} u\\
  &=&\int_0^1 P\left(C(T)-C(M) \leq \ln{(\tfrac1u-1)}\right){\rm d}u\\
  &=&\int_0^1 P\Big(-\tfrac12 L_{TM} \geq \Delta_{MT}-\ln{(\tfrac1u-1)}\Big) {\rm d}u.
\end{eqnarray*}
Let $Q_T=M \cup T$ and recall that $L_{TM}=L_{Q_TM}-L_{Q_TT}$. For any $\gamma\in (\frac{1}{2},1]$
$$
-\tfrac12 L_{TM} -(\Delta_{MT}-\ln{(\tfrac1u-1)})\;=\;\left(\tfrac12 L_{Q_TT}-\ln{\Bigl(\frac{e^{\gamma\Delta_{MT}}}{\tfrac1u-1}\Bigl)}\right)-\Bigl(\tfrac12 L_{Q_TM}+(1-\gamma)\Delta_{MT}\Bigl).
$$
Observe that for any random variables $U$, $V$, and any $\epsilon,\gamma' \geq 0$, the event $
\{U-V \geq 0\}$ implies $\{U \geq \gamma' \epsilon\}\cup
\{V  < \gamma' \epsilon\}$. Let $U=\tfrac12 L_{Q_TT}-\ln{\Bigl(\frac{e^{\gamma\Delta_{MT}}}{\tfrac1u-1}\Bigl)}$ and $V=\tfrac12 L_{Q_TM}+(1-\gamma)\Delta_{MT}$. Take $\epsilon=\mu_{Q_TM}$ and $\gamma'=\tfrac16({1-\gamma})$. Recall that $A_T=\gamma\Delta_{MT}+\gamma'\mu_{Q_TM}$ and
$$
\{U\geq \gamma'\epsilon\}\;=\;\{\tfrac12 L_{Q_TT}\geq \ln{\Bigl(\frac{e^{\gamma\Delta_{MT}}}{\tfrac1u-1}}\Bigl)+\gamma'\mu_{Q_TM}\}\;=\;\{\tfrac12 L_{Q_TT}\geq \ln{\Bigl(\frac{e^{A_T}}{\tfrac1u-1}}\Bigl)\},
$$
$$
\{V<\gamma'\epsilon\}\;=\;\{\tfrac12 L_{Q_TM}<-(1-\gamma)\Delta_{MT}+\gamma'\mu_{Q_TM}\}\;=\;\{\tfrac12 L_{Q_TM}<\gamma'(\mu_{Q_TM}-6\Delta_{MT})\}.
$$
By the union bound we have:
\begin{equation}\label{eq:prooflemmas20nonspu}
    \E\left(NC(M)\right)  \leq \int_0^1P\Bigl(\tfrac12 L_{Q_TT}\geq \ln{\Bigl(\frac{e^{A_T}}{\tfrac1u-1}\Bigl)}\Bigl){\rm d}u+ P\Bigl(\tfrac12 L_{Q_TM}<\gamma'(\mu_{Q_TM}-6\Delta_{MT})\Bigl).
\end{equation}
We first bound the first term in \eqref{eq:prooflemmas20nonspu}. If $M \subset T$ then $Q_T=T$, $L_{Q_TT}=0$, and this term is zero. Suppose now that $T$ is a strict subset of $Q_T$. By Lemma~\ref{lemma:s7}, $L_{Q_TT}$ has a $\chi^2$-distribution. More precisely, $L_{Q_TT} \sim \mathcal{\chi}_{|Q_T\setminus T|}^2(\mu_{Q_TT})$ with $|Q_T\setminus T|=|M\setminus T|$. By assumption, $A_T>0$, $|M\setminus T|=o(\ln(e^{A_T}))$ and $\mu_{Q_TT}=o(\ln(e^{A_T}))$, then  by Lemma~\ref{lemma:s20}, for $\alpha\in(\psi,1)$, and $n$ large enough:
\begin{equation}\label{eq:boundLQS}
\int_0^1P\Bigl(L_{Q_TT}>2\ln{\Bigl(\frac{e^{A_T}}{1 / u-1}\Bigl)}\Bigl){\rm d}u \;<\; 
e^{-\alpha A_T}.
\end{equation}
We now bound the second term in \eqref{eq:prooflemmas20nonspu}. If $T \subset M$, then $Q_T=M$, $L_{Q_TM}=0$, $\mu_{Q_TM}=0$, and this term is zero. If $M$ is a strict subset of $Q_T$ then, by Lemma~\ref{lemma:s7}, $L_{Q_TM} \sim \chi_{|Q_T\setminus M|}^2(\mu_{Q_TM})$ with $|Q_T\setminus M|=|T\setminus M|$.
Clearly, when $\mu_{Q_TM}\leq 6\Delta_{MT}$, this term is also zero so suppose  $\mu_{Q_TM}> 6\Delta_{MT}$.  
We have  
$$P\left(L_{Q_TM}<2\gamma'(\mu_{Q_TM}-6\Delta_{MT})\right)\; \leq \; P\left(L_{Q_TM}<2\gamma'\mu_{Q_TM}\right)$$ 
and, by Lemma~\ref{lemma:s2}, as $\gamma'\in (0,\tfrac{1}{12})$
$$P\left(L_{Q_TM}<2\gamma'\mu_{Q_TM}\right) \;\;\leq\;\; (2\gamma')^{\frac{|T\setminus M|}{4}} e^{-\frac{1}{2}(1-\sqrt{2\gamma'})^2\mu_{Q_TM}}\;\;\leq\;\;  (\tfrac12)^{\tfrac34|T\setminus M|}e^{-\frac{1}{2}(1-\sqrt{2\gamma'})^2\mu_{Q_TM}}.$$
Since $\mu_{Q_TM}>6\Delta_{MT}$, we get
$$
A_T\;=\;\gamma'\mu_{Q_TM}+\gamma\Delta_{MT}\;<\;\frac{1-\gamma}{6}\mu_{Q_TM}+\frac{\gamma}{6}\mu_{Q_TM}\;=\;\frac16\mu_{Q_TM}\;\leq\;\frac12(1-\sqrt{2\gamma'})^2\mu_{Q_TM},
$$
where the last inequality follows from the fact that $\gamma'\in(0, \tfrac{1}{12})$. We conclude:
\begin{equation}\label{eq:boundLQM}
    P\left(L_{Q_TM}<2\gamma'(\mu_{Q_TM}-8\Delta_{MT})\right) \;\leq\; \left(\frac12\right)^{\tfrac34|T\setminus M|}e^{-A_T} \leq e^{-\alpha A_T}
\end{equation}
Summing the bounds in \eqref{eq:boundLQS} and \eqref{eq:boundLQM} gives that for $n$ large enough $\E\left(NC(M)\right) < 2e^{-\alpha A_T}$. Since $\psi<\alpha$, for $n$ large enough $\E\left(NC(M)\right) < e^{-\psi A_T}$.

\subsection{Proof of Lemma~\ref{lem:upperboundnoncentralgen}}
By Lemma~\ref{lemma:s7}, we have
\begin{eqnarray*}
    \mu_{Q_TT}&=&{\bs \beta^*_{Q_T\setminus T}}^{\top}\bs X_{Q_T\setminus T}^{\top}  \big(I_n-P_T\big) \bs X_{Q_T\setminus T} \bs \beta^*_{Q_T\setminus T}\\
    &=&{\bs \beta^*_{M\setminus T}}^{\top}\bs X_{M\setminus T}^{\top}  \big(I_n-P_T\big) \bs X_{M\setminus T} \bs \beta^*_{M\setminus T}\\
    &=&{\bs \beta^*_{(S\cap M)\setminus T}}^{\top}\bs X_{(S\cap M)\setminus T}^{\top}  \big(I_n-P_T\big) \bs X_{(S\cap M)\setminus T} \bs \beta^*_{(S\cap M)\setminus T}
\end{eqnarray*}
where the third equality follows from $\bs \beta^*_{M\setminus S}=0$.
We start with the upper bound. Denote for any square matrix $A$, its largest eigenvalue 
$\lambda_{\max}(A)$. We have:
$$\mu_{Q_TT}\leq n\lambda_{\max}\Big(\frac{1}{n}\bs X_{(S\cap M)\setminus T}^{\top}  \big(I_n-P_T\big) \bs X_{(S\cap M)\setminus T}\Big) \|\bs \beta^*_{(S\cap M)\setminus T}\|^2$$
Let $B=\tfrac{1}{n}\bs X_{(S\cap M)\setminus T}^{\top}  \big(I_n-P_T\big) \bs X_{(S\cap M)\setminus T}$ and $C=\tfrac{1}{n}\bs X_{(S\cap M)\cup T}^{\top}\bs X_{(S\cap M)\cup T}$. $B$ is the Schur complement $\big(C \,  / \, \tfrac{1}{n}\bs X_T^{\top} \bs X_T\big)$. The inverse $B^{-1}$ is then a diagonal block of $C^{-1}$. By interlacing. $\lambda_{min}(B^{-1})\geq \lambda_{min}(C^{-1})$ and then $\lambda_{\max}(B)\leq \lambda_{\max}(C)$. We also have $(S\cap M)\cup T \subseteq S$, then by interlacing again $\lambda_{\max}(C)\leq \lambda_{\max}\Big(\tfrac{1}{n}\bs X_S^{\top}\bs X_S\Big)=\bar{\lambda}$. The upper bound follows from also observing that $\|\bs \beta^*_{(S\cap M)\setminus T}\|_2^2\leq \sum_{j=1}^b |(S_j\cap M_j)\setminus T_j|\max_{i\in (S_j\cap M_j)\setminus T_j}{\beta_i^*}^2$.

A lower bound on $\mu_{Q_TT}$ is: 
$$\mu_{Q_TT}\geq n\lambda_{\min}\Big(\frac{1}{n}\bs X_{(S\cap M)\setminus T}^{\top}  \big(I_n-P_T\big) \bs X_{(S\cap M)\setminus T}\Big) \|\bs \beta^*_{(S\cap M)\setminus T}\|_2^2$$
The final bound is obtained by similar interlacing argument as for the upper bound and the bound $\|\bs \beta^*_{(S\cap M)\setminus T}\|_2^2\geq \sum_{j=1}^b |(S_j\cap M_j)\setminus T_j|\min_{i\in (S_j\cap M_j)\setminus T_j}{\beta_i^*}^2$

\section{Gaussian Sequence Model with finite number of active signals}\label{sec:appendixGSMsfinite}

We derive here selection properties of the block $\ell_0$ penalties in the Gaussian sequence model dropping assumption~A3 and assuming instead:\\[5pt]
(A11) \quad for all $j$, \;$s_j \,< \, \infty$.\\[5pt]
Changing A3 for A11 implies redeveloping results relative to the probability of false negatives and consequently sufficient and necessary betamin assumptions. Proposition~\ref{prop:isthresh} on the equivalence between block penalties and thresholding in the GSM as well as Property \ref{prop:suffnecIblockthresh} on the probability of false positives do not assume A3, do not depend on results assuming A3 and hold equally under A11.

\subsection{Selection based on block thresholds}

Consider the betamin assumption:\\[5pt]
(A12) \quad for all $j$, \; $\sqrt{n}(\beta_{\min,j}^*-\tau_j) \to \infty$.\\[5pt]
Under A11 and A12, $\hat S^b$ holds the following properties.
\begin{prop}\label{prop:suffnecIIblockthreshbis}
 In the sequence model \eqref{eq:GSM}, assume A1, A2 and A11.
 \begin{longlist}
  \item If A12 holds, $P(\hat{S}^{b}\supseteq S ) \to 1$.
  \item If A4 and A12 hold,  $P(\hat{S}^{b}=S  ) \to 1$.
 \end{longlist}
\end{prop}
Under A11, A12 is then sufficient for $\hat{S}^{b}$ to hold the screening property. When, in addition, assumption A4, that requires the block thresholds grow at least as fast as $\sqrt{2\ln(p_j-s_j)/n}$, holds, $\hat{S}^{b}$ is variable selection consistent.

In Proposition~\ref{prop:necIIblockthresh} (i) we showed that A12 is also necessary for $\hat{S}^{b}$ to hold the screening property, independently of assumptions on the $s_j$'s. It follows that, under A11, A12 is a necessary and sufficient assumption for $\hat{S}^{b}$ to hold the property. By Proposition~\ref{prop:suffnecIblockthresh} A4 is necessary and sufficient for the vanishing of the FWER. We then obtain the next proposition.

\begin{lem}\label{lem:necbetaminblockthresh2}
    In the sequence model \eqref{eq:GSM}, assume A1, A2, A11 and that there exists $j\in\{1,\ldots,b\}$ such that
\begin{equation}\label{eq:necbetaminblockthreshbis}
      \lim_{n\to \infty}\sqrt{n}\beta_{\min,j}^*-\sqrt{2\ln(p_j-s_j)} < +\infty  
    \end{equation}
then \;  $P(\hat{S}^{b} = S)\not \to 1$
\end{lem}
By Lemma~\ref{lem:necbetaminblockthresh2}, under A11, a necessary assumption for asymptotic support recovery is:
\begin{equation}\label{eq:nectypeIIGSMfinite}
    \sqrt{n}\beta_{\min,j}^*-\sqrt{2\ln(p_j-s_j)} \to \infty
\end{equation}

The earlier Theorem~\ref{theo:fullrecovGS} on rates of convergence with $\hat S^b$ holds under assumptions A4 and A5, independently of A3. Observe that under A11, A12 implies A5 for $n$ large enough. Then Theorem~\ref{theo:fullrecovGS} holds equally under A11, assuming A4 and A12.

It follows that the gains in terms of convergence rate with informed block thresholding under A11 are the same as those as under A3 shown in Sections \ref{sec:prosconsblock} and \ref{sec:examples}. To discuss gains in terms of conditions for consistency, we assume without loss of generality that the smallest active signal is located in block $b$, that is $\beta_{\min}^*=\beta_{\min,b}^*$. By Propositions~\ref{prop:suffnecIIblockthreshbis} and Lemma~\ref{lem:necbetaminblockthresh2}, sufficient and necessary conditions for the standard threshold selector $\hat S$ to be variable selection consistent essentially require:
\begin{equation}\label{eq:rangethreshexfins}
     \sqrt{\frac{2\ln(p-s)}{n}}\; \leq \;\tau\; < \;\beta_{min}^*
\end{equation}
While for $\hat{S}^b$, those conditions require:
\begin{eqnarray}\label{eq:rangeblockthreshexfins}
    \sqrt{\frac{2\ln(p_1-s_1)}{n}}\; &\leq& \;\tau_1\; < \;\beta_{min,1}^* \nonumber\\
    &\vdots&\\
    \sqrt{\frac{2\ln(p_b-s_b)}{n}}\; &\leq& \;\tau_b\; < \;\beta_{min}^* \nonumber\end{eqnarray}
By definition the ranges above are wider for $\hat S^b$ than for $\hat S$. Conditions to have selection consistency are then milder with $\hat S^b$. The next corollary gives precise conditions in which $\hat S^b$ is variable selection consistent while $\hat S$ is not.
\begin{cor}\label{cor:sfin}
 In the sequence model \eqref{eq:GSM}, assume A1, A2, A4, A11 and A12.  If
\begin{eqnarray*}
&\lim_{n\to\infty}\sqrt{n}\beta_{\min}^*-\sqrt{2\ln(p-s)} \;<\;+\infty
\end{eqnarray*}
then 
$P(\hat{S} = S)\not \to 1$ and $P(\hat{S}^{b} = S)\to 1$.  
\end{cor}
Observe that equations in \eqref{eq:rangethreshexfins} and \eqref{eq:rangeblockthreshexfins} are the same as under A3, up to a term logarithmic in the number of active signals. In Section~\ref{sec:examples}, we discuss the values of $n$ that make possible support recovery with $\hat S$ and $\hat S$ in different examples under A3. It becomes possible when $\beta^*_{\min} \geq \ln(p-s)+\ln(s)$ with $\hat S$, when $\beta^*_{\min,j} \geq \ln(p_j-s_j)+\ln(s_j)$ for all $j$ with $\hat S^b$. Since in the examples 1, 2 and 4, $\ln(s)=o(\ln(p-s))$ and for all $j$, $\ln(s_j)=o(\ln(p_j-s_j))$, the gains with $\hat S^b$ in those examples are essentially the same as under A3. In example 3, the gain is even larger than under A3: the value of $n$ that makes recovery possible with $\hat S^b$ in block 2 becomes $O(\ln(\ln(n)))$ against $O(\sqrt{\ln(n)}$ in Section~\ref{sec:examples}, while that with $\hat S$ remains essentially unchanged.

Setting $\tau=\sqrt{2\ln(p-s)/n}$ and $\tau_j=\sqrt{2\ln(p_j-s_j)/n}$ for all $j$, we get that the smallest signal recoverable by $\hat S$ and $\hat S^b$ while still being consistent, $\underline{\beta}^{*,b}_{\min,orth}$ and $\underline{\beta}^{*}_{\min,orth}$, must essentially satisfy:
$$\underline{\beta}^{*,b}_{\min,orth}>\sqrt{2\ln(p_b-s_b)/n}\quad \text{and}\quad \underline{\beta}^{*}_{\min,orth}>\sqrt{2\ln(p-s)/n}$$
The lower bounds above are the same as the expressions compared in Section~\ref{sec:examples} up to a term logarithmic in the number of active signals.  Since in the examples 1, 2 and 4, $\ln(s)=o(\ln(p-s))$ and for all $j$, $\ln(s_j)=o(\ln(p_j-s_j))$, the ratio of the bounds is essentially the same as the ratio  $\underline{\beta}^{*,b}_{\min,orth}/\underline{\beta}^{*}_{\min,orth}$ under A3. The gains in terms of smallest signal recoverable in those examples are the same as under A3.

\subsection{Proofs}

\subsubsection{Proof of Proposition~\ref{prop:suffnecIIblockthreshbis}}
By the proof of Proposition~\ref{prop:suffnecIIblockthresh}, we have: \begin{equation*}
P(\hat{S}^{b}\not\supseteq  S)\;\leq\; \sum_{j=1}^b\frac{e^{-\tfrac{n}{2}\left((\beta^\star_{\min,j}-\tau_j)^2-\tfrac{2\ln(s_j)}{n}\right)^2}}{\sqrt{\pi\ln(s_j)}}.
 \end{equation*}
By assumption A12, the right-hand side goes to 0 which shows (i). By Proposition~\ref{prop:suffnecIblockthresh}, since A4 is assumed to hold, we also have $P\big(\hat{S}^{b} \subseteq S\big)\to 1$. We conclude $P\big(\hat{S}^{b} = S\big)\to 1$.

\subsubsection{Proof of Lemma~\ref{lem:necbetaminblockthresh2}}
 We have $\sqrt{n}\beta_{\min,j}^*-\sqrt{2\ln(p_j-s_j)} \,=\, \sqrt{n}\beta_{\min,j}^*-\sqrt{n}\tau_j+\sqrt{2\ln(p_j-s_j)}\Big(\tfrac{\sqrt{n}\tau_j}{\sqrt{2\ln(p_j-s_j)}}-1\Big)$ and \eqref{eq:necbetaminblockthreshbis} implies that there exists $c\in\R^+$ such that:
  \begin{equation}\label{eq:proofnecbetaminbis}
      \lim_{n\to \infty}\sqrt{n}\beta_{\min,j}^*-\sqrt{n}\tau_j+\sqrt{2\ln(p_j-s_j)}\Big(\tfrac{\sqrt{n}\tau_j}{\sqrt{2\ln(p_j-s_j)}}-1\Big)\leq c
  \end{equation}

Consider the case $\lim_{n\to +\infty}\tfrac{\sqrt{n}\tau_j}{\sqrt{2\ln(p_j-s_j)}}-1<0$. Then by Proposition~\ref{prop:suffnecIblockthresh},
$P(\hat{S}^{b} \subseteq S)\not \to 1$ and $P(\hat{S}^{b} = S)\not \to 1$. 
Now consider the case $\lim_{n\to +\infty}\tfrac{\sqrt{n}\tau_j}{\sqrt{2\ln(p_j-s_j)}}-1 \geq0$. Then by \eqref{eq:proofnecbetaminbis},  $\lim_{n\to \infty}\sqrt{n}\beta_{\min,j}^*-\sqrt{n}\tau_j \leq c$ and by Proposition~\ref{prop:necIIblockthresh} i), $P(\hat{S}^{b} \supseteq S)\not \to 1$ and $P(\hat{S}^{b} = S)\not \to 1$.

\subsubsection{Proof of Corollary~\ref{cor:sfin}}
Since $\hat{S}$ is $\hat{S}^b$ with $b=1$, the assumptions of Lemma~\ref{lem:necbetaminblockthresh2} for $\hat{S}$ are met and $P(\hat{S} = S)\not \to 1$. Since A4 and A12 hold, by Proposition~\ref{prop:suffnecIblockthresh} and Proposition~\ref{prop:suffnecIIblockthreshbis}, $P(\hat{S}^b \subseteq 1)\to 1$ and $P( \hat{S}^b \supseteq S)\to 1$, and then $P(\hat{S}^b = S)\to 1$.

\section{Non linear block $\ell_0$ penalties in high-dimensional linear regression}\label{sec:appendixnonlinear}
\sdescription{Short description of Supplement C.}

\subsection{Selection properties of nonlinear block $\ell_0$ penalties}

In this section we show the variable selection consistency of the block $\ell_0$ penalties in the linear regression without assuming linearity of the penalty functions. We let, for all $j=1,\ldots,b$, $\eta_j$ be any non-negative and increasing penalty functions on the natural numbers. The selector based on those block penalties is:
\begin{equation}
  \hat S^b\;\in\;\arg\max_{M\in \M} \left\{\max_{\beta\in \mathcal{L}_M} \ell(\bs y;\bs \beta)-\sum_{j=1}^b\eta_j (|M_j|)\right\}.
\end{equation}

Rewrite the difference in penalty between any model $M$ and $S$:
\begin{equation}\label{eq:DeltaMbis}
\Delta_{MS} \; := \; \sum_{j=1}^b \eta_j(|M_j|)-\eta_j(|S_j|)
\end{equation}
We define the normalized difference in block penalty between any model $M$ and $S$ as:
 $$\tilde{\kappa}_{M, j}\; := \; \begin{cases}
     \frac{\eta_j(|M_j|)-\eta_j(|S_j|)}{|M_j|\,-\,|S_j|} \;\;\; \text{if}\;\;\; |M_j|\,\neq\,|S_j| \\
     0 \;\;\;\text{if}\;\;\; |M_j|\,=\,|S_j|.
 \end{cases}  $$
The quantity $\tilde{\kappa}_{M, j}$ plays a similar role as $\kappa_j$. It is the average penalty incurred for adding a variable from block $B_j$ to the model $M$. Unlike in the linear case, $\tilde{\kappa}_{M, j}$ is not constant but a function of $|M_j|$ and $|S_j|$. If the penalty is linear, $\tilde{\kappa}_{M, j}$ is exactly $\kappa_j$ for all $M\in\M$.

To show the consistency of $\hat S^b$, we replace assumption A6 by an assumption on the $\tilde{\kappa}_{M, j}$ and make a betamin assumption.\\[5pt]
(A13) \quad For each block $j$, there exists $f'_j\to \infty$ (as $n\to \infty$) such that for all $M \in \M$ such that $|M_j| \neq |S_j|$ and $|M_j \setminus S_j| >0$, for $n$ sufficiently large:
\begin{equation*}
    \tilde{\kappa}_{M, j} =\;\ln\bigg(\frac{p_j-s_j}{|M_j \setminus S_j| }\bigg) + f'_j.
\end{equation*}
(A14) \quad For each block $j$, there exists $l_j\to +\infty$ such that for $n$ sufficiently large,
    \begin{equation*}
        \sqrt{\frac{(1-\gamma) n\rho(\bs X)}{6}}{{\beta_{\min}^*}_j} \,- \,\sqrt{\ln(p_j-s_j) + f'_j} \;=\; \sqrt{\ln(s_j)} + l_j
    \end{equation*}
where $\gamma:=\frac12(1+\max_j \frac{\ln(p_j-s_j)}{\ln(p_j-s_j) + f'_j})\in (\tfrac12,1)$.\\[5pt] 

We can now state the main result of this supplement: 
\begin{thm}\label{theo:suffcondlinearmodelnonlinpen}

Under assumptions A1, A2, A13 and, A14, we have
$$\sum_{M \in \M\setminus \{S\}}\E\left(NC(M)\right)\to 0 \qquad \mbox{and} \qquad P(\hat{S}^b = S) \to 1.$$
\end{thm}
Theorem~\ref{theo:suffcondlinearmodelnonlinpen} is consistent with results in the literature. A standard nonlinear standard penalty in high-dimensional variable selection is the EBIC penalty \cite{EBIC} that sets for some $\zeta\geq 0$:
\begin{equation}\label{eq:EBICpen}
    \eta(|M|)=\zeta\ln{\binom{p}{|M|}} + \tfrac{1}{2}\ln(n).
\end{equation}
The penalty can be shown to satisfy assumption~A13 under a restriction on the number of active signals. A corollary of Theorem~\ref{theo:suffcondlinearmodelnonlinpen} is as follows.
\begin{cor}\label{cor:ebic}
    Assume A1, A2, $\eta$ as in \eqref{eq:EBICpen} with $\zeta\geq1$, $s^{\zeta}=o(\sqrt{n})$, and that there exists $l\to +\infty$ such that for $n$ sufficiently large
    \begin{equation}\label{eq:betaminebic}
        \sqrt{\frac{\ln \big(\tfrac{\sqrt{n}}{(s+1)^{\zeta}}\big)-2\zeta}{\ln(p-s) + \ln \big(\tfrac{\sqrt{n}}{(s+1)^{\zeta}}\big)-2\zeta}\frac{n\rho(\bs X)}{12}}{{\beta_{\min}^*}} \,- \,\sqrt{\zeta\ln(p-s) +\ln(\sqrt{n}2e^{\zeta})} \;=\; \sqrt{\ln(s)} + l .
    \end{equation}
    Then A13 and A14 hold and: $$\sum_{M \in \M\setminus \{S\}}\E\left(NC(M)\right)\to 0 \qquad \mbox{and} \qquad P(\widehat{S}^b = S) \to 1.$$    
\end{cor}
Corollary~\ref{cor:ebic} shows the consistency of the standard EBIC penalty under milder betamin conditions than the literature \cite{EBIC,ebicdiverging}. It also shows the EBIC is \emph{strongly} selection consistent.

Observe that the assumptions and proof strategy of Theorem~\ref{theo:suffcondlinearmodelnonlinpen} are similar to those of Theorem~\ref{theo:suffcondlinearmodel}. We do not develop them here but results analogous to Theorem~\ref{prop:fullrecovreg} on convergence rates and the necessary conditions in Section~\ref{ssec:neccond_linreg} can be obtained for the nonlinear penalties. We also expect the benefits of linear block penalties to extend to the nonlinear ones.
\subsection{Proofs}
\subsubsection{Proof of Theorem~\ref{theo:suffcondlinearmodelnonlinpen}}
The proof is essentially the same as the proof of Theorem~\ref{theo:suffcondlinearmodel} replacing $\kappa_j$ by $\tilde{\kappa}_{M, j}$ for every $j=1,\ldots,b$. We redevelop only the steps altered by the non-linearity of the penalty. We refer the reader to the proof of Theorem~\ref{theo:suffcondlinearmodel} for more details.

Observe that $\Delta_{MS}$, defined in \eqref{eq:DeltaMbis}, can be rewritten as $\Delta_{MS}=\sum_{j=1}^b(|M_j\setminus S_j|-|S_j\setminus M_j|)\tilde{\kappa}_{M}$. By  Lemma~\ref{lem:boundrho}, for every $n\in \N$ we have 
\begin{equation}\label{eq:nonoverfittingcondgen2}
   A_S\;=\; \gamma \Delta_{MS}+\frac{1-\gamma}{6}\mu_{Q_SM}\;\geq\; \gamma\sum_{j=1}^b|M_j \setminus S_j|\tilde{\kappa}_{M} + \sum_{j=1}^b |S_j\setminus M_j |\left(\tfrac{1-\gamma}{6} n \rho(\bs X) {\beta^*_{\min,j}}^{2} - \gamma\tilde{\kappa}_{M, j}\right) 
\end{equation} 
Proceeding as in the proof of Theorem~\ref{theo:suffcondlinearmodel}, we immediately have $\mu_{Q_S S}=o(A_S)$. We also have:
$$\frac{|M\setminus S|}{A_S}\leq \Big[ \gamma\min_{j=1,\ldots,b}\tilde{\kappa}_{M, j} \Big]^{-1} $$
and by assumption A13, $|M\setminus S|=o(A_S)$ when $|M\setminus S|\neq 0$ too.
We conclude that for every $M\in \mathcal M\setminus \{S\}$, by Lemma~\ref{lem:boundexpnc}, for any $\psi\in(0,1)$ and $n$ large enough:
$$
\E(NC(M))\;\leq\;e^{-\psi A_S}.
$$
Let 
$$A_S^*\;:=\;\gamma\sum_{j=1}^b|M_j \setminus S_j|\kappa_j + \sum_{j=1}^b |S_j\setminus M_j |\left(\tfrac{1-\gamma}{6} n \rho(\bs X) {\beta^*_{\min,j}}^{2} - \gamma\kappa_j\right).$$
By \eqref{eq:nonoverfittingcondgen2}, we have for $n$ large enough:
\begin{equation}\label{eq:boundexp3}
    \E(NC(M))\;\leq\;e^{-\psi A_S^*}.
\end{equation}
Since for all $M\in\M$, $\tilde{\kappa}_{M, j} \leq \ln(p_j-s_j)+f'_j$, assumption A14 implies there exist $g'_j \to +\infty$ such that:
$$\frac{(1-\gamma) n\rho(\bs X)}{6}{{\beta_{\min}^*}_j}^2 \,- \,\tilde{\kappa}_{M, j} \;=\; \ln(s_j) + g'_{j}.$$
To choose $\psi$ we proceed similarly to the proof of Theorem~\ref{theo:suffcondlinearmodel}. Let $\delta \in (0,1)$ and denote:
$$\bar{m}_j=
    \max\bigg\{\max_{M\in\M\text{s.t.}|M_j\setminus S_j|\neq 0}\tfrac{2\ln\big(\tfrac{p_j-s_j}{|M_j\setminus S_j|}\big)}{f'_j},\tfrac{2\ln(s_j)}{g'_{j}}\bigg\}$$
Take $\psi=\max_{j=1, \ldots,b}\frac{a + \delta + \bar{m}_j}{1+\bar{m}_j}$ for some $a \in(0, 1- \delta)$ then $\psi \in (0,1)$ and for every $j=1,\ldots,b$:
$$\delta \leq \psi-(1-\psi)\bar{m}_j.$$ 
By definition of $\gamma$ and since for any $M$, $\ln(p_j-s_j/|M_j\setminus S_j|) \leq \ln(p_j-s_j)$, we have for every $j=1,\ldots,b$ such that $|M_j \setminus S_j|\neq0$:
$$
\gamma\tilde{\kappa}_{M, j}\;\geq\;\tfrac12 \bigg(1+\frac{\ln(p_j-s_j)}{\ln(p_j-s_j) + f'_j}\bigg)\tilde{\kappa}_{M, j}\geq \tfrac12 \bigg(1+\tfrac{\ln\big(\tfrac{p_j-s_j}{|M_j\setminus S_j|}\big)}{\tilde{\kappa}_{M, j}}\bigg)\tilde{\kappa}_{M, j} \;=\; \ln\big(\tfrac{p_j-s_j}{|M_j\setminus S_j|}\big)+ \tfrac12f'_j.
$$
From our choice of $\psi$ and $\delta$, we get:
$$\psi\gamma\tilde{\kappa}_{M, j} \;\geq\; \psi\bigg(\ln\big(\tfrac{p_j-s_j}{|M_j\setminus S_j|}\big)+ \frac12f'_j\bigg) \;\geq\; \ln\big(\tfrac{p_j-s_j}{|M_j\setminus S_j|}\big)+ \delta\frac12f'_j.$$
Using that for any $M$, $\tilde{\kappa}_{M, j} \leq \ln(p_j-s_j)+f'_j$, by A14 and our choice of $\psi,\delta$, we have:
$$\psi\left(\tfrac{1-\gamma}{6} n \rho(\bs X) {\beta^*_{\min,j}}^{2} - \gamma\tilde{\kappa}_{M, j}\right) \;\geq\; \psi\left(\ln(s_j)+g_j'\right) \;\geq\; \ln(s_j)+\delta \frac{g'_j}{2}.
$$
Using the expression for $A_S^*$, the bound in \eqref{eq:boundexp3}, and the above inequalities, we get
$$
\E(NC(M))\;\leq\;\exp\{- \sum_{j\, \text{s.t.} |M_j \setminus S_j| > 0} |M_j\setminus S_j|\bigg(\ln\bigg(\frac{p_j-s_j}{|M_j\setminus S_j|}\bigg)+\delta \tfrac{f'_j}{2}\bigg)-\sum_{j=1}^b |S_j\setminus M_j|(\ln(s_j)+\delta \tfrac{g'_j}{2})\}.
$$
Following the same subsequent steps of the proof of Theorem~\ref{theo:suffcondlinearmodel} shows variable selection consistency.

\subsubsection{Proof of Corollary~\ref{cor:ebic}}
We start by showing that $\eta$ satisfies assumption~A13. Let $M\in\M$ and denote $|M\setminus S|=u$ and $|S\setminus M|=w$. We have $|M|=u+s-w$ and $|M|-|S|=u-w$. For $\eta$, assumption A13 deals with models $M$ such that $|M|\neq |S|$ and $|M\setminus S|>0$. We thus restrict our attention to $M\in\mathcal{M}$ such that $u- w\neq0$ and $u>0$. We consider first the case $u-w>0$, we have:
    \begin{equation*}
        \tilde{\kappa}_{M}=\frac{\zeta}{u-w}\ln\bigg[\frac{\binom{p}{|M|}}{\binom{p}{|S|}}\bigg]+ \frac{1}{2}\ln(n)
    \end{equation*}
    A well-known property of binomial coefficients is that for any positive integers $n$, $h$, $k$ we have:
    \begin{equation}\label{eq:idbinom}
        \binom{n}{h}\binom{n-h}{k}=\binom{n}{h+k}\binom{h+k}{h}.
    \end{equation}
    Taking $n=p$, $k=u-w$ and $h=s$ in \eqref{eq:idbinom}, we get:
    \begin{equation}\label{eq:identitybinom}
        \tilde{\kappa}_{M}=\frac{\zeta}{u-w}\ln\bigg[\frac{\binom{p -s}{u-w}}{\binom{|M|}{s}}\bigg]+ \frac{1}{2}\ln(n)        .
    \end{equation}
Standard bounds on binomial coefficient for $1\leq k\leq n$ are:
\begin{equation}\label{eq:boundbinom}
  \left(\frac{n}{k}\right)^k \leq  \binom{n}{k} \leq \frac{n^n}{k^k(n-k)^{n-k}}
\end{equation}
Combining \eqref{eq:identitybinom} and \eqref{eq:boundbinom}, we get:
\begin{equation}\label{eq:gggg}
    \frac{\zeta}{u-w}\ln\bigg[\frac{\binom{p}{|M|}}{\binom{p}{s}}\bigg] \geq \zeta\ln\bigg(\frac{p-s}{u-w}\bigg)-\zeta\bigg[\frac{s}{u-w}\ln\bigg(\frac{u-w}{s}+1\bigg) +\ln\bigg(1+\frac{s}{u-w}\bigg) \bigg].
\end{equation}
We have:
\begin{equation}\label{eq:gggg2}
\zeta\ln\bigg(\frac{p-s}{u-w}\bigg) \geq \zeta\ln\bigg(\frac{p-s}{u}\bigg) \geq \ln\bigg(\frac{p-s}{u}\bigg).\end{equation}
Observe also that $f:x \mapsto \frac{\ln(x+1)}{x}+\ln(1+{x}^{-1})$ for $x>0$ is decreasing and that $\frac{u-w}{s} \geq s^{-1}$. By \eqref{eq:gggg} and \eqref{eq:gggg2}, we then have: 
\begin{equation}\label{eq:lowerboundb1ebic2}
        \tilde{\kappa}_{M} \geq \ln\big(\frac{p-s}{u}\big)-\zeta s\ln(s^{-1}+1)+ \ln \big(\tfrac{\sqrt{n}}{(1+s)^{\zeta}}\big)
\end{equation}
Since $\lim_{n\to+\infty}s\ln(s^{-1}+1)=1$ and by assumption $s^{\zeta}=o(\sqrt{n})$, $\eta$ satisfies assumption~A13 in the case $u-w>0$.
Consider now the case where $u-w<0$. We have:
     \begin{equation*}
        \tilde{\kappa}_{M}=\frac{\zeta}{w-u}\ln\bigg[\frac{\binom{p}{s}}{\binom{p}{|M|}}\bigg]+ \frac{1}{2}\ln n
    \end{equation*}
 Taking $n=p$, $k=w-u$ and $h=|M|$ in \eqref{eq:idbinom} gives:
     \begin{equation}\label{eq:identitybinom2}
     \tilde{\kappa}_{M}=\frac{\zeta}{w-u}\ln\bigg[\frac{\binom{p - |M|}{w-u}}{\binom{s}{|M|}}\bigg]+ \frac{1}{2}\ln n.
    \end{equation}
By the bounds in \eqref{eq:boundbinom}, we get:
$$ \frac{\zeta}{w-u}\ln\bigg[\frac{\binom{p}{s}}{\binom{p}{|M|}}\bigg] \geq \zeta\ln\bigg(\frac{p-|M|}{w-u}\bigg)+\zeta\bigg[\frac{s}{w-u}\ln\bigg(1-\frac{w-u}{s}\bigg) -\ln\bigg(\frac{s}{w-u}-1\bigg) \bigg]$$
Since $u-w<0$, $|M|<s$. Using also that $\zeta\geq 1$, we have:
$$\zeta\ln\bigg(\frac{p-|M|}{w-u}\bigg) \geq \zeta\ln\bigg(\frac{p-s}{w-u}\bigg) \geq \zeta\ln\bigg(\frac{p-s}{w}\bigg)\geq \ln\bigg(\frac{p-s}{u}\bigg)$$
Observe that $g:x \mapsto \frac{\ln(1-x)}{x}-\ln({x}^{-1}-1)$ for $x\in(0,1)$ is increasing and that $1>\frac{w-u}{s} \geq s^{-1}$. We get: 
\begin{equation}\label{eq:bounbkappaebicunderfit}
        \tilde{\kappa}_{M} \geq \ln\big(\frac{p-s}{u}\big)+\zeta s\ln(1-s^{-1})+ \ln \big(\tfrac{\sqrt{n}}{(s-1)^{\zeta}}\big)
\end{equation}
Note that the right-hand side above is well defined. Indeed, since $u>0$ and $u-w<0$, we have $2\leq w\leq s$. Since $\lim_{n\to+\infty}s\ln(1-s^{-1})=-1$ and by assumption $s^{\zeta}=o(\sqrt{n})$, $\eta$ satisfies assumption~A13 in the case $u-w<0$ too.

We now show that if \eqref{eq:betaminebic} holds then A14 holds. If $u-w>0$, by \eqref{eq:identitybinom}, the upper bound in \eqref{eq:upperboudbinom}, and the lower bound in \eqref{eq:boundbinom}:
 \begin{eqnarray*}
     \tilde{\kappa}_M &\leq& \zeta\ln\Big(\frac{(p-s)e}{u-w}\Big)-\zeta(1+\frac{s}{u-w})\ln\Big(1+\frac{u-w}{s}\Big)+\frac{1}{2}\ln(n)\\
     &\leq& \ln\big(\frac{p-s}{u}\big)+(\zeta-1)\ln(p-s)+\zeta+\frac{1}{2}\ln(n)
 \end{eqnarray*}
If $w-u>0$, by \eqref{eq:identitybinom2}, the upper bound in \eqref{eq:upperboudbinom}, and the lower bound in \eqref{eq:boundbinom}:
 \begin{eqnarray*}
     \tilde{\kappa}_M &\leq& 
     \zeta\ln\Big(\frac{(p-s+(w-u))e}{w-u}\Big)+\zeta(\frac{s}{w-u}-1)\ln\Big(1
     -\frac{w-u}{s}\Big)+\frac{1}{2}\ln(n)\\
     &\leq&\ln\big(\frac{p-s}{u}+1\big)+(\zeta-1)\ln(p-s+1)+\zeta+\frac{1}{2}\ln(n)
 \end{eqnarray*}
 We get that for all $M\in \mathcal{M}$ such that $|M|\neq |S|$ and $|M\setminus S|>0$:
 \begin{eqnarray*}
    f'= \tilde{\kappa}_{M} - \ln\Big(\frac{p-s}{|M\setminus S|}\Big) 
    &\leq& \ln(1+\frac{u}{p-s})+(\zeta-1)\ln(p-s+1)+\zeta+\frac{1}{2}\ln(n)\\
    &\leq& \ln(2)+(\zeta-1)\ln(p-s+1)+\zeta+\frac{1}{2}\ln(n).
 \end{eqnarray*}
Then, we have:
 \begin{eqnarray*}
    \ln(p-s)+ f'
    &\leq& \zeta\ln(p-s+1)+\ln(2e^{\zeta})+\frac{1}{2}\ln(n).
 \end{eqnarray*}
 and
\begin{equation}\label{eq:xxxxx}
     - \,\sqrt{\ln(p-s) + f'} \geq - \,\sqrt{\zeta\ln(p-s+1)+\ln(2e^{\zeta})+\frac{1}{2}\ln(n)} .
\end{equation}
If $u-w>0$, by \eqref{eq:lowerboundb1ebic2}:
\begin{equation*}
        \tilde{\kappa}_{M} \geq \ln\big(\frac{p-s}{u}\big)-\zeta+ \ln \big(\tfrac{\sqrt{n}}{(1+s)^{\zeta}}\big) \geq \ln\big(\frac{p-s}{u}\big)-2\zeta+ \ln \big(\tfrac{\sqrt{n}}{(1+s)^{\zeta}}\big).
\end{equation*}
If $w-u>0$, then $s\geq 2$ and by \eqref{eq:bounbkappaebicunderfit}:
$$\tilde{\kappa}_{M} \geq \ln\big(\frac{p-s}{u}\big)+\zeta 2\ln(\frac{1}{2})+ \ln \big(\tfrac{\sqrt{n}}{(s-1)^{\zeta}}\big) \geq \ln\big(\frac{p-s}{u}\big)-2\zeta + \ln \big(\tfrac{\sqrt{n}}{(s+1)^{\zeta}}\big).$$
We have that for all $M\in \mathcal{M}$ such that $|M|\neq |S|$ and $|M\setminus S|>0$:
$$f'= \tilde{\kappa}_{M} - \ln\Big(\frac{p-s}{|M\setminus S|}\Big) \geq \ln \big(\tfrac{\sqrt{n}}{(s+1)^{\zeta}}\big)-2\zeta$$
and,
\begin{eqnarray*}
    1-\frac12\Big(1+ \frac{\ln(p-s)}{\ln(p-s) + f'}\Big)
    &\geq& 1-\frac12\Big(1+\frac{\ln(p-s)}{\ln(p-s) + \ln \big(\tfrac{\sqrt{n}}{(s+1)^{\zeta}}\big)-2\zeta}\Big)\\
    &\geq&\frac12\Big(\frac{\ln \big(\tfrac{\sqrt{n}}{(s+1)^{\zeta}}\big)-2\zeta}{\ln(p-s) + \ln \big(\tfrac{\sqrt{n}}{(s+1)^{\zeta}}\big)-2\zeta}\Big)\\
\end{eqnarray*}
From the last inequality and \eqref{eq:xxxxx}, we get that when \eqref{eq:betaminebic} holds, A14 holds. The result then follows directly from Theorem \ref{theo:suffcondlinearmodelnonlinpen}.


\end{document}
